\input{preamble}

% OK, start here.
%
\begin{document}

\title{Topologies sur les schémas}

\maketitle

\phantomsection
\label{section-phantom}

\tableofcontents



\section{Introduction}
\label{section-introduction}

\noindent
Dans ce chapitre, nous décrivons les différentes topologies sur la
catégorie des schémas. On pourra consulter, entre autres, \cite{SGA1} et \cite{Ner}.
Avant de le faire, nous souhaitons signaler qu'il existe de nombreux
choix de sites (au sens de
Sites, définition \ref{sites-definition-site}) qui donnent lieu à
la même notion de faisceau sur la catégorie sous-jacente. Ainsi,
nos choix peuvent différer légèrement de ceux des références,
mais conduisent finalement aux mêmes groupes de cohomologie, etc.

\section{La méthode générale}
\label{section-procedure}

\noindent
Dans cette section, nous exposons une méthode générale de construction des
sites avec lesquels nous travaillerons. Supposons que l'on veuille étudier les faisceaux
sur les schémas relativement à une topologie $\tau$. Pour obtenir
un site, comme dans Sites, définition \ref{sites-definition-site},
formé de schémas munis de cette topologie, il faut prendre quelques précautions. En effet,
on ne peut pas simplement dire « considérons tous les schémas munis de la topologie de Zariski »,
car on obtiendrait ainsi une catégorie « grande ». Dans chaque section de
ce chapitre, nous procéderons plutôt comme suit :
\begin{enumerate}
\item On définit une classe $\text{Cov}_\tau$ de recouvrements de schémas
satisfaisant aux axiomes de Sites, définition \ref{sites-definition-site}.
Un recouvrement ouvert de Zariski d'un
schéma sera toujours un recouvrement pour $\tau$.
\item On distingue, dans la catégorie des schémas affines, une notion de
$\tau$-recouvrement standard.
\item On définit ce qu'est un gros $\tau$-site « absolu » $\Sch_\tau$.
Il s'agit des sites obtenus en choisissant convenablement un ensemble de schémas
et un ensemble de recouvrements.
\item Pour tout objet $S$ de $\Sch_\tau$,
on définit le gros $\tau$-site $(\Sch/S)_\tau$ et, lorsque
$\tau$ s'y prête, le petit\footnote{Les mots gros et
petit ne renvoient pas ici à la taille des catégories
correspondantes.} $\tau$-site $S_\tau$.
\item En outre, il existe un site $(\textit{Aff}/S)_\tau$ fondé sur la
notion de $\tau$-recouvrement standard d'affines\footnote{Dans le cas
de la topologie ph, nous nous écartons très légèrement de cette méthode ; voir la
définition \ref{definition-big-small-ph} et la discussion qui l'entoure.}
dont la catégorie des faisceaux
est équivalente à la catégorie des faisceaux sur $(\Sch/S)_\tau$.
\end{enumerate}
Cette construction est quelque peu malcommode, car elle ne fournit pas de choix canonique
du gros $\tau$-site d'un schéma, ni même du petit
$\tau$-site d'un schéma. Si l'on accepte d'ignorer les difficultés
ensemblistes, on peut travailler avec des classes et obtenir
des gros et petits sites canoniques…







\section{La topologie de Zariski}
\label{section-zariski}

\begin{definition}
\label{definition-zariski-covering}
Soit $T$ un schéma. Un {\it recouvrement de Zariski de $T$} est une famille
de morphismes de schémas $\{f_i : T_i \to T\}_{i \in I}$
telle que chaque $f_i$ soit une immersion ouverte et que
$T = \bigcup f_i(T_i)$.
\end{definition}

\noindent
Cela définit une classe (propre) de recouvrements.
Montrons à présent que cette notion satisfait aux conditions de
Sites, définition \ref{sites-definition-site}.

\begin{lemma}
\label{lemma-zariski}
Soit $T$ un schéma.
\begin{enumerate}
\item Si $T' \to T$ est un isomorphisme, alors $\{T' \to T\}$
est un recouvrement de Zariski de $T$.
\item Si $\{T_i \to T\}_{i\in I}$ est un recouvrement de Zariski et si, pour tout
$i$, on dispose d'un recouvrement de Zariski $\{T_{ij} \to T_i\}_{j\in J_i}$, alors
$\{T_{ij} \to T\}_{i \in I, j\in J_i}$ est un recouvrement de Zariski.
\item Si $\{T_i \to T\}_{i\in I}$ est un recouvrement de Zariski
et si $T' \to T$ est un morphisme de schémas, alors
$\{T' \times_T T_i \to T'\}_{i\in I}$ est un recouvrement de Zariski.
\end{enumerate}
\end{lemma}

\begin{proof}
Omis.
\end{proof}

\begin{lemma}
\label{lemma-zariski-affine}
Soit $T$ un schéma affine. Soit $\{T_i \to T\}_{i \in I}$ un
recouvrement de Zariski de $T$. Il existe alors un recouvrement de Zariski
$\{U_j \to T\}_{j = 1, \ldots, m}$ qui raffine
$\{T_i \to T\}_{i \in I}$ et tel que chaque $U_j$ soit un ouvert
principal de $T$ ; voir
Schémas, définition \ref{schemes-definition-standard-covering}.
On peut en outre choisir chaque $U_j$ comme ouvert de l'un des $T_i$.
\end{lemma}

\begin{proof}
Cela résulte de ce que $T$ est quasi-compact et que les ouverts principaux forment une base
de sa topologie. C'est également démontré dans
Schémas, lemme \ref{schemes-lemma-standard-open}.
\end{proof}

\noindent
On définit donc comme suit les recouvrements standards d'affines correspondants.

\begin{definition}
\label{definition-standard-Zariski}
Comparer à Schémas, définition \ref{schemes-definition-standard-covering}.
Soit $T$ un schéma affine. Un {\it recouvrement de Zariski standard}
de $T$ est un recouvrement de Zariski $\{U_j \to T\}_{j = 1, \ldots, m}$
tel que chaque $U_j \to T$ induise un isomorphisme sur un ouvert affine principal
de $T$.
\end{definition}

\begin{definition}
\label{definition-big-zariski-site}
Un {\it gros site de Zariski} est un site $\Sch_{Zar}$, au sens de
Sites, définition \ref{sites-definition-site}, construit comme suit :
\begin{enumerate}
\item Choisir un ensemble quelconque de schémas $S_0$, ainsi qu'un ensemble de recouvrements de Zariski
$\text{Cov}_0$ entre ces schémas.
\item Prendre pour catégorie sous-jacente à $\Sch_{Zar}$
une catégorie $\Sch_\alpha$ construite comme dans
Ensembles, lemme \ref{sets-lemma-construct-category}, à partir de l'ensemble $S_0$.
\item Prendre pour recouvrements de $\Sch_{Zar}$ un ensemble de recouvrements obtenu comme dans
Ensembles, lemme \ref{sets-lemma-coverings-site}, à partir de la
catégorie $\Sch_\alpha$, de la classe des recouvrements de Zariski
et de l'ensemble $\text{Cov}_0$ choisi ci-dessus.
\end{enumerate}
\end{definition}

\noindent
Sites, lemme \ref{sites-lemma-choice-set-coverings-immaterial}, montre
qu'une fois la catégorie $\Sch_\alpha$ choisie, la
catégorie des faisceaux sur $\Sch_\alpha$ ne dépend pas du
choix des recouvrements effectué au point (3) ci-dessus. Autrement dit, le topos
$\Sh(\Sch_{Zar})$ ne dépend que du choix de
la catégorie $\Sch_\alpha$. Il est démontré dans
Ensembles, lemme \ref{sets-lemma-what-is-in-it}, que ces catégories
sont stables par de nombreuses constructions de géométrie algébrique, par exemple
les produits fibrés et le passage aux sous-schémas ouverts ou fermés. On peut également montrer
que le choix précis de $\Sch_\alpha$ importe
assez peu ; voir la section \ref{section-change-alpha}.

\medskip\noindent
Une autre méthode consisterait à supposer l'existence d'un
cardinal fortement inaccessible et à définir $\Sch_{Zar}$
comme la catégorie des schémas contenus dans un univers choisi, les
recouvrements étant les recouvrements de Zariski contenus dans ce même
univers.

\medskip\noindent
Avant de poursuivre par l'introduction du gros site de Zariski d'un
schéma $S$, observons que la topologie d'un gros site de Zariski
$\Sch_{Zar}$ est, en un certain sens, induite par la topologie de Zariski
sur la catégorie de tous les schémas.

\begin{lemma}
\label{lemma-zariski-induced}
Soit $\Sch_{Zar}$ un gros site de Zariski comme dans la
définition \ref{definition-big-zariski-site}.
Soit $T \in \Ob(\Sch_{Zar})$.
Soit $\{T_i \to T\}_{i \in I}$ un recouvrement de Zariski quelconque de $T$.
Il existe un recouvrement $\{U_j \to T\}_{j \in J}$ de $T$ dans le site
$\Sch_{Zar}$ qui soit tautologiquement équivalent (voir
Sites, définition \ref{sites-definition-combinatorial-tautological})
à $\{T_i \to T\}_{i \in I}$.
\end{lemma}

\begin{proof}
Puisque chaque $T_i \to T$ est une immersion ouverte, il résulte de
Ensembles, lemme \ref{sets-lemma-what-is-in-it},
que chaque $T_i$ est isomorphe à un objet $V_i$ de $\Sch_{Zar}$.
Le recouvrement $\{V_i \to T\}_{i \in I}$ est tautologiquement équivalent
à $\{T_i \to T\}_{i \in I}$ (en utilisant dans les deux sens l'application identité de $I$).
En outre, $\{V_i \to T\}_{i \in I}$ est combinatoirement équivalent à un
recouvrement $\{U_j \to T\}_{j \in J}$ de $T$ dans le site $\Sch_{Zar}$ d'après
Ensembles, lemme \ref{sets-lemma-coverings-site}.
\end{proof}

\begin{definition}
\label{definition-big-small-Zariski}
Soit $S$ un schéma. Soit $\Sch_{Zar}$ un gros site de Zariski
contenant $S$.
\begin{enumerate}
\item Le {\it gros site de Zariski de $S$}, noté
$(\Sch/S)_{Zar}$, est le site $\Sch_{Zar}/S$
introduit dans Sites, section \ref{sites-section-localize}.
\item Le {\it petit site de Zariski de $S$}, que nous notons
$S_{Zar}$, est la sous-catégorie pleine de $(\Sch/S)_{Zar}$
dont les objets sont les $U/S$ tels que $U \to S$ soit une immersion ouverte.
Un recouvrement de $S_{Zar}$ est tout recouvrement $\{U_i \to U\}$ de
$(\Sch/S)_{Zar}$ avec $U \in \Ob(S_{Zar})$.
\item Le {\it gros site de Zariski affine de $S$}, noté
$(\textit{Aff}/S)_{Zar}$, est la sous-catégorie pleine de
$(\Sch/S)_{Zar}$ formée des objets $U/S$ tels que $U$ soit un
schéma affine. Un recouvrement de $(\textit{Aff}/S)_{Zar}$ est tout recouvrement
$\{U_i \to U\}$ de $(\Sch/S)_{Zar}$ avec $U \in \Ob((\textit{Aff}/S)_{Zar})$
qui soit un recouvrement de Zariski standard.
\item Le {\it petit site de Zariski affine de $S$}, noté
$S_{affine, Zar}$, est la sous-catégorie pleine de $S_{Zar}$
dont les objets sont les $U/S$ tels que $U$ soit un schéma affine.
Un recouvrement de $S_{affine, Zar}$ est tout recouvrement $\{U_i \to U\}$ de
$S_{Zar}$ avec $U \in \Ob(S_{affine, Zar})$
qui soit un recouvrement de Zariski standard.
\end{enumerate}
\end{definition}

\noindent
Il n'est pas tout à fait évident que le petit site de Zariski,
le gros site de Zariski affine et le petit site de Zariski affine
soient bien des sites. Vérifions-le.

\begin{lemma}
\label{lemma-verify-site-Zariski}
Soit $S$ un schéma. Soit $\Sch_{Zar}$ un gros site de Zariski
contenant $S$. Les structures $S_{Zar}$, $(\textit{Aff}/S)_{Zar}$
et $S_{affine, Zar}$ définies ci-dessus sont des sites.
\end{lemma}

\begin{proof}
Montrons que $S_{Zar}$ est un site. C'est une catégorie munie d'un
ensemble donné de familles de morphismes de but fixé. Il faut donc
établir les propriétés (1), (2) et (3) de
Sites, définition \ref{sites-definition-site}.
Puisque $(\Sch/S)_{Zar}$ est un site, il suffit de prouver
que, pour tout recouvrement $\{U_i \to U\}$ de $(\Sch/S)_{Zar}$
avec $U \in \Ob(S_{Zar})$, on a aussi $U_i \in \Ob(S_{Zar})$.
Cela résulte des définitions,
car la composée de deux immersions ouvertes est une immersion ouverte.

\medskip\noindent
Montrons que $(\textit{Aff}/S)_{Zar}$ est un site.
En raisonnant comme ci-dessus, il suffit de montrer que l'ensemble
des recouvrements de Zariski standards d'affines satisfait aux propriétés
(1), (2) et (3) de
Sites, définition \ref{sites-definition-site}.
Soit $R$ un anneau. Soient $f_1, \ldots, f_n \in R$ engendrant l'idéal unité.
Pour chaque $i \in \{1, \ldots, n\}$, soient $g_{i1}, \ldots, g_{in_i} \in R_{f_i}$
des éléments engendrant l'idéal unité de $R_{f_i}$. On peut écrire
$g_{ij} = f_{ij}/f_i^{e_{ij}}$. Après avoir remplacé, si nécessaire,
$f_{ij}$ par $f_i f_{ij}$, on a
$D(f_{ij}) \subset D(f_i) \cong \Spec(R_{f_i})$
égal à $D(g_{ij}) \subset \Spec(R_{f_i})$. On voit donc que
la famille de morphismes $\{D(g_{ij}) \to \Spec(R)\}$
est un recouvrement de Zariski standard. Ces considérations
montrent que (2) vaut pour les recouvrements de Zariski standards.
Nous omettons la vérification de (1) et de (3).

\medskip\noindent
Nous omettons la démonstration du fait que $S_{affine, Zar}$ est un site.
\end{proof}

\begin{lemma}
\label{lemma-fibre-products-Zariski}
Soit $S$ un schéma. Soit $\Sch_{Zar}$ un gros site de Zariski
contenant $S$. Les catégories sous-jacentes aux sites
$\Sch_{Zar}$, $(\Sch/S)_{Zar}$,
$S_{Zar}$, $(\textit{Aff}/S)_{Zar}$ et $S_{affine, Zar}$ possèdent des produits fibrés.
Dans chaque cas, le foncteur évident vers la catégorie $\Sch$ de
tous les schémas commute aux produits fibrés. Les catégories
$(\Sch/S)_{Zar}$ et $S_{Zar}$ possèdent toutes deux un objet final,
à savoir $S/S$.
\end{lemma}

\begin{proof}
Pour $\Sch_{Zar}$, c'est vrai par construction ; voir
Ensembles, lemme \ref{sets-lemma-what-is-in-it}.
Supposons donnés des morphismes de schémas $U \to S$, $V \to U$, $W \to U$
avec $U, V, W \in \Ob(\Sch_{Zar})$.
Le produit fibré $V \times_U W$ dans $\Sch_{Zar}$
est un produit fibré dans $\Sch$ ; c'est aussi
le produit fibré de $V/S$ et $W/S$ au-dessus de $U/S$ dans
la catégorie de tous les schémas au-dessus de $S$, et donc également un
produit fibré dans $(\Sch/S)_{Zar}$.
Cela démontre le résultat pour $(\Sch/S)_{Zar}$.
Si $U \to S$, $V \to U$ et $W \to U$ sont des immersions ouvertes, alors
$V \times_U W \to S$ l'est aussi, d'où le résultat pour $S_{Zar}$.
Si $U, V, W$ sont affines, $V \times_U W$ l'est également, d'où le
résultat pour $(\textit{Aff}/S)_{Zar}$ et $S_{affine, Zar}$.
\end{proof}

\noindent
Vérifions à présent que le gros, resp.\ petit, site affine définit le même
topos que le gros, resp.\ petit, site.

\begin{lemma}
\label{lemma-affine-big-site-Zariski}
Soit $S$ un schéma. Soit $\Sch_{Zar}$ un gros site de Zariski
contenant $S$.
Le foncteur $(\textit{Aff}/S)_{Zar} \to (\Sch/S)_{Zar}$
est un foncteur cocontinu spécial. Il induit donc une équivalence
de topos de $\Sh((\textit{Aff}/S)_{Zar})$ vers
$\Sh((\Sch/S)_{Zar})$.
\end{lemma}

\begin{proof}
La notion de foncteur cocontinu spécial est introduite dans
Sites, définition \ref{sites-definition-special-cocontinuous-functor}.
Il faut donc vérifier les hypothèses (1) -- (5) de
Sites, lemme \ref{sites-lemma-equivalence}.
Notons le foncteur d'inclusion
$u : (\textit{Aff}/S)_{Zar} \to (\Sch/S)_{Zar}$.
La cocontinuité signifie simplement que tout recouvrement de Zariski de
$T/S$, avec $T$ affine, peut être raffiné par un recouvrement de Zariski standard de $T$.
C'est le contenu du
lemme \ref{lemma-zariski-affine}.
Ainsi, (1) est satisfaite. Le foncteur $u$ est continu, simplement parce qu'un recouvrement
de Zariski standard est un recouvrement de Zariski. Ainsi, (2) est satisfaite.
Les assertions (3) et (4) résultent immédiatement de ce que $u$ est
pleinement fidèle. Enfin, la condition (5) résulte de ce que
tout schéma possède un recouvrement ouvert affine.
\end{proof}

\begin{lemma}
\label{lemma-alternative-zariski}
Soit $S$ un schéma. Soit $\Sch_{Zar}$ un gros site de Zariski
contenant $S$. Le foncteur $S_{affine, Zar} \to S_{Zar}$
est un foncteur cocontinu spécial. Il induit donc une équivalence
de topos de $\Sh(S_{affine, Zar})$ vers $\Sh(S_{Zar})$.
\end{lemma}

\begin{proof}
Omis. Indication : comparer à la démonstration du
lemme \ref{lemma-affine-big-site-Zariski}.
\end{proof}

\noindent
Vérifions que la notion de faisceau sur le petit site de Zariski
correspond à celle de faisceau sur $S$.

\begin{lemma}
\label{lemma-Zariski-usual}
La catégorie des faisceaux sur $S_{Zar}$ est équivalente à la
catégorie des faisceaux sur l'espace topologique sous-jacent à $S$.
\end{lemma}

\begin{proof}
Nous utiliserons à plusieurs reprises le fait que, pour tout objet
$U/S$ de $S_{Zar}$, le morphisme $U \to S$ est un isomorphisme
sur un sous-schéma ouvert.
Soit $\mathcal{F}$ un faisceau sur $S$. On définit alors un faisceau
sur $S_{Zar}$ par la règle $\mathcal{F}'(U/S) = \mathcal{F}(\Im(U \to S))$.
Réciproquement, choisissons pour tout sous-schéma ouvert $U \subset S$ un objet
$U'/S \in \Ob(S_{Zar})$ tel que $\Im(U' \to S) = U$
(il faut ici utiliser Ensembles, lemme \ref{sets-lemma-what-is-in-it}).
Étant donné un faisceau $\mathcal{G}$ sur $S_{Zar}$, on définit un faisceau sur $S$ en posant
$\mathcal{G}'(U) = \mathcal{G}(U'/S)$. Pour voir que $\mathcal{G}'$ est
un faisceau, on utilise que, pour tout recouvrement ouvert $U = \bigcup_{i \in I} U_i$,
le recouvrement $\{U_i \to U\}_{i \in I}$
est combinatoirement équivalent à un recouvrement $\{U_j' \to U'\}_{j \in J}$
dans $S_{Zar}$ d'après Ensembles, lemme \ref{sets-lemma-coverings-site},
et l'on applique Sites, lemme \ref{sites-lemma-tautological-same-sheaf}.
Détails omis.
\end{proof}

\noindent
Désormais, nous ne ferons plus de distinction entre un faisceau sur
$S_{Zar}$ et un faisceau sur $S$. Nous utiliserons toujours les procédés
de la démonstration du lemme pour passer d'une notion à l'autre.
Établissons maintenant quelques relations entre les topos
associés à ces sites.

\begin{lemma}
\label{lemma-put-in-T}
Soit $\Sch_{Zar}$ un gros site de Zariski.
Soit $f : T \to S$ un morphisme dans $\Sch_{Zar}$.
Le foncteur $T_{Zar} \to (\Sch/S)_{Zar}$
est cocontinu et induit un morphisme de topos
$$
i_f :
\Sh(T_{Zar})
\longrightarrow
\Sh((\Sch/S)_{Zar})
$$
Pour tout faisceau $\mathcal{G}$ sur $(\Sch/S)_{Zar}$,
on a la formule $(i_f^{-1}\mathcal{G})(U/T) = \mathcal{G}(U/S)$.
Le foncteur $i_f^{-1}$ possède aussi un adjoint à gauche $i_{f, !}$ qui commute
aux produits fibrés et aux égalisateurs.
\end{lemma}

\begin{proof}
Notons le foncteur $u : T_{Zar} \to (\Sch/S)_{Zar}$.
Autrement dit, étant donnée une immersion ouverte $j : U \to T$ correspondant
à un objet de $T_{Zar}$, on pose $u(U \to T) = (f \circ j : U \to S)$.
Ce foncteur commute aux produits fibrés ; voir le
lemme \ref{lemma-fibre-products-Zariski}.
De plus, $T_{Zar}$ possède des égalisateurs (car deux morphismes de même
source et de même but sont égaux) et $u$ commute à ceux-ci.
Il est manifestement cocontinu.
Il est aussi continu, puisque $u$ transforme les recouvrements en recouvrements et
commute aux produits fibrés. Le lemme résulte donc de
Sites, lemmes \ref{sites-lemma-when-shriek}
et \ref{sites-lemma-preserve-equalizers}.
\end{proof}

\begin{lemma}
\label{lemma-at-the-bottom}
Soit $S$ un schéma. Soit $\Sch_{Zar}$ un gros site de Zariski
contenant $S$.
Le foncteur d'inclusion $S_{Zar} \to (\Sch/S)_{Zar}$
satisfait aux hypothèses de Sites, lemme \ref{sites-lemma-bigger-site},
et induit donc un morphisme de sites
$$
\pi_S : (\Sch/S)_{Zar} \longrightarrow S_{Zar}
$$
et un morphisme de topos
$$
i_S : \Sh(S_{Zar}) \longrightarrow \Sh((\Sch/S)_{Zar})
$$
tels que $\pi_S \circ i_S = \text{id}$. De plus, $i_S = i_{\text{id}_S}$,
où $i_{\text{id}_S}$ est celui du lemme \ref{lemma-put-in-T}. En particulier, le
foncteur $i_S^{-1} = \pi_{S, *}$ est décrit par la règle
$i_S^{-1}(\mathcal{G})(U/S) = \mathcal{G}(U/S)$.
\end{lemma}

\begin{proof}
Dans ce cas, le foncteur $u : S_{Zar} \to (\Sch/S)_{Zar}$,
outre les propriétés relevées dans la démonstration du
lemme \ref{lemma-put-in-T} ci-dessus, est pleinement fidèle
et transforme l'objet final en l'objet final.
Le lemme en résulte.
\end{proof}

\begin{definition}
\label{definition-restriction-small-zariski}
Dans la situation du
lemme \ref{lemma-at-the-bottom},
le foncteur $i_S^{-1} = \pi_{S, *}$ est souvent
appelé la {\it restriction au petit site de Zariski} et, pour un faisceau
$\mathcal{F}$ sur le gros site de Zariski, on note $\mathcal{F}|_{S_{Zar}}$
cette restriction.
\end{definition}

\noindent
Avec cette notation, pour un faisceau $\mathcal{F}$ sur le
gros site et un faisceau $\mathcal{G}$ sur le petit site, on a
\begin{align*}
\Mor_{\Sh(S_{Zar})}(\mathcal{F}|_{S_{Zar}}, \mathcal{G})
& =
\Mor_{\Sh((\Sch/S)_{Zar})}(\mathcal{F},
i_{S, *}\mathcal{G}) \\
\Mor_{\Sh(S_{Zar})}(\mathcal{G}, \mathcal{F}|_{S_{Zar}})
& =
\Mor_{\Sh((\Sch/S)_{Zar})}(\pi_S^{-1}\mathcal{G},
\mathcal{F})
\end{align*}
De plus, on a $(i_{S, *}\mathcal{G})|_{S_{Zar}} = \mathcal{G}$
et $(\pi_S^{-1}\mathcal{G})|_{S_{Zar}} = \mathcal{G}$.

\begin{lemma}
\label{lemma-morphism-big}
Soit $\Sch_{Zar}$ un gros site de Zariski.
Soit $f : T \to S$ un morphisme dans $\Sch_{Zar}$.
Le foncteur
$$
u : (\Sch/T)_{Zar} \longrightarrow (\Sch/S)_{Zar},
\quad
V/T \longmapsto V/S
$$
est cocontinu et possède un adjoint à droite continu
$$
v : (\Sch/S)_{Zar} \longrightarrow (\Sch/T)_{Zar},
\quad
(U \to S) \longmapsto (U \times_S T \to T).
$$
Ils induisent le même morphisme de topos
$$
f_{big} :
\Sh((\Sch/T)_{Zar})
\longrightarrow
\Sh((\Sch/S)_{Zar})
$$
On a $f_{big}^{-1}(\mathcal{G})(U/T) = \mathcal{G}(U/S)$.
On a $f_{big, *}(\mathcal{F})(U/S) = \mathcal{F}(U \times_S T/T)$.
De plus, $f_{big}^{-1}$ possède un adjoint à gauche $f_{big!}$ qui commute aux
produits fibrés et aux égalisateurs.
\end{lemma}

\begin{proof}
Le foncteur $u$ est cocontinu, continu et commute aux produits fibrés
et aux égalisateurs (détails omis ; comparer à la démonstration du
lemme \ref{lemma-put-in-T}).
Par conséquent,
Sites, lemmes \ref{sites-lemma-when-shriek} et
\ref{sites-lemma-preserve-equalizers},
s'appliquent et donnent la formule
pour $f_{big}^{-1}$ ainsi que l'existence de $f_{big!}$. De plus,
le foncteur $v$ est adjoint à droite car, pour $U/T$ et $V/S$ donnés,
on a $\Mor_S(u(U), V) = \Mor_T(U, V \times_S T)$,
comme voulu. On peut donc appliquer
Sites, lemmes \ref{sites-lemma-have-functor-other-way} et
\ref{sites-lemma-have-functor-other-way-morphism},
pour obtenir la formule pour $f_{big, *}$.
\end{proof}

\begin{lemma}
\label{lemma-morphism-big-small}
Soit $\Sch_{Zar}$ un gros site de Zariski.
Soit $f : T \to S$ un morphisme dans $\Sch_{Zar}$.
\begin{enumerate}
\item On a $i_f = f_{big} \circ i_T$, où $i_f$ est celui du
lemme \ref{lemma-put-in-T} et $i_T$ celui du
lemme \ref{lemma-at-the-bottom}.
\item Le foncteur $S_{Zar} \to T_{Zar}$,
$(U \to S) \mapsto (U \times_S T \to T)$, est continu et induit
un morphisme de sites
$$
f_{small} :
T_{Zar}
\longrightarrow
S_{Zar}
$$
Les foncteurs $f_{small}^{-1}$ et $f_{small, *}$ coïncident avec
les notions usuelles $f^{-1}$ et $f_*$ lorsque l'on identifie les faisceaux
sur $T_{Zar}$, resp.\ $S_{Zar}$, aux faisceaux sur $T$, resp.\ $S$,
au moyen du lemme \ref{lemma-Zariski-usual}.
\item On a un diagramme commutatif de morphismes de sites
$$
\xymatrix{
T_{Zar} \ar[d]_{f_{small}} &
(\Sch/T)_{Zar} \ar[d]^{f_{big}} \ar[l]^{\pi_T} \\
S_{Zar} &
(\Sch/S)_{Zar} \ar[l]_{\pi_S}
}
$$
de sorte que $f_{small} \circ \pi_T = \pi_S \circ f_{big}$ comme morphismes de topos.
\item On a $f_{small} = \pi_S \circ f_{big} \circ i_T = \pi_S \circ i_f$.
\end{enumerate}
\end{lemma}

\begin{proof}
L'égalité $i_f = f_{big} \circ i_T$ résulte de
l'égalité $i_f^{-1} = i_T^{-1} \circ f_{big}^{-1}$, qui est
claire d'après les descriptions de ces foncteurs données ci-dessus.
Cela démontre (1).

\medskip\noindent
Assertion (2) : voir Sites, exemple \ref{sites-example-continuous-map}.

\medskip\noindent
L'assertion (3) résulte de ce que $\pi_S$ et $\pi_T$ sont donnés par
les foncteurs d'inclusion, tandis que $f_{small}$ et $f_{big}$ le sont par le
foncteur de changement de base $U \mapsto U \times_S T$.

\medskip\noindent
L'assertion (4) se déduit de (3) par précomposition avec $i_T$.
\end{proof}

\noindent
Dans la situation du lemme, avec la terminologie de la
définition \ref{definition-restriction-small-zariski},
on a, pour tout faisceau $\mathcal{F}$ sur le gros site de Zariski de $T$,
$$
(f_{big, *}\mathcal{F})|_{S_{Zar}} =
f_{small, *}(\mathcal{F}|_{T_{Zar}}),
$$
Cette égalité résulte immédiatement de la commutativité du diagramme de
sites du lemme, puisque la restriction au petit site de Zariski de
$T$, resp.\ $S$, est donnée par $\pi_{T, *}$, resp.\ $\pi_{S, *}$. La formule
analogue faisant intervenir images inverses et restrictions est fausse.

\begin{lemma}
\label{lemma-composition}
Étant donnés des schémas $X$, $Y$, $Z$ dans $(\Sch/S)_{Zar}$
et des morphismes $f : X \to Y$, $g : Y \to Z$, on a
$g_{big} \circ f_{big} = (g \circ f)_{big}$ et
$g_{small} \circ f_{small} = (g \circ f)_{small}$.
\end{lemma}

\begin{proof}
Cela résulte de la description simple de l'image directe
et de l'image inverse pour les foncteurs entre gros sites donnée par le
lemme \ref{lemma-morphism-big}. Pour les foncteurs
entre petits sites, c'est le contenu de
Faisceaux, lemme \ref{sheaves-lemma-pushforward-composition},
via l'identification du lemme \ref{lemma-Zariski-usual}.
\end{proof}

\begin{lemma}
\label{lemma-morphism-big-small-cartesian-diagram}
Soit $\Sch_{Zar}$ un gros site de Zariski. Considérons un diagramme cartésien
$$
\xymatrix{
T' \ar[r]_{g'} \ar[d]_{f'} & T \ar[d]^f \\
S' \ar[r]^g & S
}
$$
dans $\Sch_{Zar}$. Alors
$i_g^{-1} \circ f_{big, *} = f'_{small, *} \circ (i_{g'})^{-1}$
et $g_{big}^{-1} \circ f_{big, *} = f'_{big, *} \circ (g'_{big})^{-1}$.
\end{lemma}

\begin{proof}
Puisque le diagramme est cartésien, on a, pour $U'/S'$,
l'égalité $U' \times_{S'} T' = U' \times_S T$. Ainsi,
$i_g^{-1} \circ f_{big, *}$ et $f'_{small, *} \circ (i_{g'})^{-1}$
envoient tous deux un faisceau $\mathcal{F}$ sur $(\Sch/T)_{Zar}$ sur le faisceau
$U' \mapsto \mathcal{F}(U' \times_{S'} T')$ sur $S'_{Zar}$
(utiliser les lemmes \ref{lemma-put-in-T} et \ref{lemma-morphism-big-small}).
La seconde égalité se démontre de la même manière ou peut se
déduire du très général
Sites, lemme \ref{sites-lemma-localize-morphism}.
\end{proof}

\noindent
On peut considérer un faisceau sur le gros site de Zariski de $S$ comme une collection
de faisceaux « usuels » sur tous les schémas au-dessus de $S$.

\begin{lemma}
\label{lemma-characterize-sheaf-big}
Soit $S$ un schéma contenu dans un gros site de Zariski $\Sch_{Zar}$.
Un faisceau $\mathcal{F}$ sur le gros site de Zariski $(\Sch/S)_{Zar}$
est donné par les données suivantes :
\begin{enumerate}
\item pour tout $T/S \in \Ob((\Sch/S)_{Zar})$, un faisceau
$\mathcal{F}_T$ sur $T$ ;
\item pour tout $f : T' \to T$ dans
$(\Sch/S)_{Zar}$, une application
$c_f : f^{-1}\mathcal{F}_T \to \mathcal{F}_{T'}$.
\end{enumerate}
Ces données sont soumises aux conditions suivantes :
\begin{enumerate}
\item[(a)] pour tous $f : T' \to T$ et $g : T'' \to T'$ dans
$(\Sch/S)_{Zar}$, la composée $c_g \circ g^{-1}c_f$
est égale à $c_{f \circ g}$ ;
\item[(b)] si $f : T' \to T$ dans $(\Sch/S)_{Zar}$ est une
immersion ouverte, alors $c_f$ est un isomorphisme.
\end{enumerate}
\end{lemma}

\begin{proof}
Ce lemme résulte d'un énoncé purement faisceautique discuté dans
Sites, remarque \ref{sites-remark-variant-glue-sheaves-absolute}.
Nous en donnons aussi une démonstration directe dans ce cas.

\medskip\noindent
Étant donné un faisceau $\mathcal{F}$ sur $\Sh((\Sch/S)_{Zar})$,
on pose $\mathcal{F}_T = i_p^{-1}\mathcal{F}$, où $p : T \to S$
est le morphisme structural. Remarquons que
$\mathcal{F}_T(U) = \mathcal{F}(U'/S)$ pour tout ouvert $U \subset T$
et toute immersion ouverte $U' \to T$ dans $(\Sch/T)_{Zar}$
d'image $U$ ; voir les lemmes \ref{lemma-Zariski-usual} et \ref{lemma-put-in-T}.
Ainsi, étant donnés $f : T' \to T$ au-dessus de $S$ et $U, U' \to T$, on obtient une
application canonique $\mathcal{F}_T(U) = \mathcal{F}(U'/S) \to \mathcal{F}(U'\times_T T'/S)
= \mathcal{F}_{T'}(f^{-1}(U))$, celle du milieu étant l'application de restriction
de $\mathcal{F}$ relativement au morphisme
$U' \times_T T' \to U'$ au-dessus de $S$. La famille de ces applications est
compatible aux restrictions et définit donc une $f$-application $c_f$
de $\mathcal{F}_T$ vers $\mathcal{F}_{T'}$ ; voir
Faisceaux, définition \ref{sheaves-definition-f-map} et la discussion
qui l'entoure. Il est clair que $c_{f \circ g}$ est la composée de
$c_f$ et $c_g$, puisque composer les applications de restriction de $\mathcal{F}$
donne encore une application de restriction.

\medskip\noindent
Réciproquement, étant donné un système $(\mathcal{F}_T, c_f)$ comme dans le lemme,
on peut définir un préfaisceau $\mathcal{F}$ sur $\Sh((\Sch/S)_{Zar})$
en posant simplement $\mathcal{F}(T/S) = \mathcal{F}_T(T)$. Comme application de restriction,
pour $f : T' \to T$ donné, on définit, pour $s \in \mathcal{F}(T)$,
l'image inverse $f^*(s)$ comme étant $c_f(s)$ (où l'on considère de nouveau $c_f$ comme
une $f$-application). La condition imposée aux $c_f$ garantit que
les images inverses satisfont à la propriété de fonctorialité requise.
Nous omettons la vérification qu'il s'agit d'un faisceau.
Il est clair que les constructions ainsi définies sont mutuellement inverses.
\end{proof}























\section{La topologie \'etale}
\label{section-etale}

\noindent
Soit $S$ un schéma. Nous voulons définir la topologie \'etale sur
la catégorie des schémas au-dessus de $S$. Conformément à notre principe général,
nous introduisons d'abord la notion de recouvrement \'etale.

\begin{definition}
\label{definition-etale-covering}
Soit $T$ un schéma. Un {\it recouvrement \'etale de $T$} est une famille
de morphismes de schémas $\{f_i : T_i \to T\}_{i \in I}$
telle que chaque $f_i$ soit \'etale et que $T = \bigcup f_i(T_i)$.
\end{definition}

\begin{lemma}
\label{lemma-zariski-etale}
Tout recouvrement de Zariski est un recouvrement \'etale.
\end{lemma}

\begin{proof}
Cela résulte immédiatement des définitions et du fait qu'une immersion ouverte
est un morphisme \'etale ; voir
Morphismes, lemme \ref{morphisms-lemma-open-immersion-etale}.
\end{proof}

\noindent
Montrons maintenant que cette notion satisfait aux conditions de
Sites, définition \ref{sites-definition-site}.

\begin{lemma}
\label{lemma-etale}
Soit $T$ un schéma.
\begin{enumerate}
\item Si $T' \to T$ est un isomorphisme, alors $\{T' \to T\}$
est un recouvrement \'etale de $T$.
\item Si $\{T_i \to T\}_{i\in I}$ est un recouvrement \'etale et si, pour tout
$i$, on a un recouvrement \'etale $\{T_{ij} \to T_i\}_{j\in J_i}$, alors
$\{T_{ij} \to T\}_{i \in I, j\in J_i}$ est un recouvrement \'etale.
\item Si $\{T_i \to T\}_{i\in I}$ est un recouvrement \'etale
et si $T' \to T$ est un morphisme de schémas, alors
$\{T' \times_T T_i \to T'\}_{i\in I}$ est un recouvrement \'etale.
\end{enumerate}
\end{lemma}

\begin{proof}
Omis.
\end{proof}

\begin{lemma}
\label{lemma-etale-affine}
Soit $T$ un schéma affine.
Soit $\{T_i \to T\}_{i \in I}$ un recouvrement \'etale de $T$.
Il existe alors un recouvrement \'etale
$\{U_j \to T\}_{j = 1, \ldots, m}$ qui raffine
$\{T_i \to T\}_{i \in I}$ et tel que chaque $U_j$ soit un schéma
affine. On peut en outre choisir chaque $U_j$ comme ouvert affine
de l'un des $T_i$.
\end{lemma}

\begin{proof}
Omis.
\end{proof}

\noindent
On définit donc comme suit les recouvrements standards d'affines correspondants.

\begin{definition}
\label{definition-standard-etale}
Soit $T$ un schéma affine. Un {\it recouvrement \'etale standard}
de $T$ est une famille $\{f_j : U_j \to T\}_{j = 1, \ldots, m}$
telle que chaque $U_j$ soit affine et \'etale sur $T$, et que
$T = \bigcup f_j(U_j)$.
\end{definition}

\noindent
Dans la définition ci-dessus, nous ne supposons {\bf pas} que les morphismes $f_j$ soient
étales standards. En effet, sous cette hypothèse, les recouvrements étales
standards ne définiraient pas un site sur $\textit{Aff}/S$, notamment à cause de
Algèbre, lemme \ref{algebra-lemma-standard-etale}, partie (4).
En revanche, un morphisme étale d'affines est automatiquement
lisse standard ; voir
Algèbre, lemme \ref{algebra-lemma-etale-standard-smooth}.
Ainsi, un recouvrement étale standard est un recouvrement lisse
standard et un recouvrement syntomique standard.

\begin{definition}
\label{definition-big-etale-site}
Un {\it gros site étale} est un site $\Sch_\etale$, au sens de
Sites, définition \ref{sites-definition-site}, construit comme suit :
\begin{enumerate}
\item Choisir un ensemble quelconque de schémas $S_0$, ainsi qu'un ensemble de recouvrements étales
$\text{Cov}_0$ entre ces schémas.
\item Prendre comme catégorie sous-jacente une catégorie $\Sch_\alpha$
construite comme dans Ensembles, lemme \ref{sets-lemma-construct-category},
à partir de l'ensemble $S_0$.
\item Choisir un ensemble de recouvrements comme dans
Ensembles, lemme \ref{sets-lemma-coverings-site}, à partir de la
catégorie $\Sch_\alpha$, de la classe des recouvrements étales
et de l'ensemble $\text{Cov}_0$ choisi ci-dessus.
\end{enumerate}
\end{definition}

\noindent
Voir les remarques suivant la définition \ref{definition-big-zariski-site}
pour la motivation et les explications relatives à la définition des gros sites.

\medskip\noindent
Avant de poursuivre par l'introduction du gros site étale d'un
schéma $S$, observons que la topologie d'un gros site étale
$\Sch_\etale$ est, en un certain sens, induite par la topologie
étale sur la catégorie de tous les schémas.

\begin{lemma}
\label{lemma-etale-induced}
Soit $\Sch_\etale$ un gros site étale comme dans la
définition \ref{definition-big-etale-site}.
Soit $T \in \Ob(\Sch_\etale)$.
Soit $\{T_i \to T\}_{i \in I}$ un recouvrement étale quelconque de $T$.
\begin{enumerate}
\item Il existe un recouvrement $\{U_j \to T\}_{j \in J}$ de $T$ dans le site
$\Sch_\etale$ qui raffine $\{T_i \to T\}_{i \in I}$.
\item Si $\{T_i \to T\}_{i \in I}$ est un recouvrement étale standard, alors
il est tautologiquement équivalent à un recouvrement dans $\Sch_\etale$.
\item Si $\{T_i \to T\}_{i \in I}$ est un recouvrement de Zariski, alors
il est tautologiquement équivalent à un recouvrement dans $\Sch_\etale$.
\end{enumerate}
\end{lemma}

\begin{proof}
Pour chaque $i$, choisissons un recouvrement ouvert affine $T_i = \bigcup_{j \in J_i} T_{ij}$
tel que chaque $T_{ij}$ s'envoie dans un sous-schéma ouvert affine de $T$. D'après le
lemme \ref{lemma-etale},
le raffinement $\{T_{ij} \to T\}_{i \in I, j \in J_i}$ est lui aussi un recouvrement étale
de $T$. On peut donc supposer chaque $T_i$ affine et s'envoyant dans
un ouvert affine $W_i$ de $T$. En appliquant
Ensembles, lemme \ref{sets-lemma-what-is-in-it},
on voit que $W_i$ est isomorphe à un objet de $\Sch_\etale$.
Mais alors $T_i$, en tant que schéma de type fini au-dessus de $W_i$,
est isomorphe à un objet $V_i$ de $\Sch_\etale$ par une seconde
application de
Ensembles, lemme \ref{sets-lemma-what-is-in-it}.
Le recouvrement $\{V_i \to T\}_{i \in I}$ raffine $\{T_i \to T\}_{i \in I}$
(puisqu'ils sont isomorphes).
De plus, $\{V_i \to T\}_{i \in I}$ est combinatoirement équivalent à un
recouvrement $\{U_j \to T\}_{j \in J}$ de $T$ dans le site
$\Sch_\etale$ d'après
Ensembles, lemme \ref{sets-lemma-what-is-in-it}.
Le recouvrement $\{U_j \to T\}_{j \in J}$ est un raffinement comme en (1).
Dans la situation de (2), (3), chacun des
schémas $T_i$ est isomorphe à un objet de $\Sch_\etale$ d'après
Ensembles, lemme \ref{sets-lemma-what-is-in-it},
et une nouvelle application de
Ensembles, lemme \ref{sets-lemma-coverings-site},
donne le résultat voulu.
\end{proof}

\begin{definition}
\label{definition-big-small-etale}
Soit $S$ un schéma. Soit $\Sch_\etale$ un gros site
étale contenant $S$.
\begin{enumerate}
\item Le {\it gros site étale de $S$}, noté
$(\Sch/S)_\etale$, est le site
$\Sch_\etale/S$ introduit dans
Sites, section \ref{sites-section-localize}.
\item Le {\it petit site étale de $S$}, que nous notons
$S_\etale$, est la sous-catégorie pleine de
$(\Sch/S)_\etale$
dont les objets sont les $U/S$ tels que $U \to S$ soit étale.
Un recouvrement de $S_\etale$ est tout recouvrement $\{U_i \to U\}$ de
$(\Sch/S)_\etale$ avec $U \in \Ob(S_\etale)$.
\item Le {\it gros site étale affine de $S$}, noté
$(\textit{Aff}/S)_\etale$, est la sous-catégorie pleine de
$(\Sch/S)_\etale$ dont les objets sont les $U/S$
tels que $U$ soit un schéma affine.
Un recouvrement de $(\textit{Aff}/S)_\etale$ est tout recouvrement
$\{U_i \to U\}$ de $(\Sch/S)_\etale$ avec $U \in \Ob((\textit{Aff}/S)_\etale)$
qui soit un recouvrement étale standard.
\item Le {\it petit site étale affine de $S$}, noté
$S_{affine, \etale}$, est la sous-catégorie pleine de $S_\etale$
dont les objets sont les $U/S$ tels que $U$ soit un schéma affine.
Un recouvrement de $S_{affine, \etale}$ est tout recouvrement $\{U_i \to U\}$
de $S_\etale$ avec $U \in \Ob(S_{affine, \etale})$ qui soit un recouvrement
étale standard.
\end{enumerate}
\end{definition}

\noindent
Il n'est pas tout à fait évident que le gros site étale affine,
le petit site étale et le petit site étale affine soient des sites.
Vérifions-le.

\begin{lemma}
\label{lemma-verify-site-etale}
Soit $S$ un schéma. Soit $\Sch_\etale$ un gros site
étale contenant $S$.
Les structures $S_\etale$, $(\textit{Aff}/S)_\etale$ et $S_{affine, \etale}$
sont des sites.
\end{lemma}

\begin{proof}
Montrons que $S_\etale$ est un site. C'est une catégorie munie d'un
ensemble donné de familles de morphismes de but fixé. Il faut donc
établir les propriétés (1), (2) et (3) de
Sites, définition \ref{sites-definition-site}.
Puisque $(\Sch/S)_\etale$ est un site, il suffit de prouver
que, pour tout recouvrement $\{U_i \to U\}$ de $(\Sch/S)_\etale$
avec $U \in \Ob(S_\etale)$, on a aussi
$U_i \in \Ob(S_\etale)$.
Cela résulte des définitions, car la composée de morphismes étales
est un morphisme étale.

\medskip\noindent
Montrons que $(\textit{Aff}/S)_\etale$ est un site.
En raisonnant comme ci-dessus, il suffit de montrer que l'ensemble
des recouvrements étales standards d'affines satisfait aux propriétés
(1), (2) et (3) de
Sites, définition \ref{sites-definition-site}.
C'est clair : par exemple, étant donné un recouvrement étale
standard $\{T_i \to T\}_{i\in I}$ et, pour chaque
$i$, un recouvrement étale standard $\{T_{ij} \to T_i\}_{j\in J_i}$,
$\{T_{ij} \to T\}_{i \in I, j\in J_i}$ est un recouvrement étale standard
car $\bigcup_{i\in I} J_i$ est fini et chaque $T_{ij}$ est affine.

\medskip\noindent
Nous omettons la démonstration du fait que $S_{affine, \'etale}$ est un site.
\end{proof}

\begin{lemma}
\label{lemma-fibre-products-etale}
Soit $S$ un schéma. Soit $\Sch_\etale$ un gros site
étale contenant $S$. Les catégories sous-jacentes des sites
$\Sch_\etale$, $(\Sch/S)_\etale$, $S_\etale$, $(\textit{Aff}/S)_\etale$,
et $S_{affine, \etale}$ admettent des produits fibrés.
Dans chaque cas, le foncteur évident vers la catégorie $\Sch$ de
tous les schémas commute à la formation des produits fibrés. Les catégories
$(\Sch/S)_\etale$ et $S_\etale$ possèdent toutes deux un
objet final, à savoir $S/S$.
\end{lemma}

\begin{proof}
Pour $\Sch_\etale$, cela est vrai par construction ; voir
Ensembles, lemme \ref{sets-lemma-what-is-in-it}.
Supposons donnés des morphismes $U \to S$, $V \to U$, $W \to U$
de schémas, avec $U, V, W \in \Ob(\Sch_\etale)$.
Le produit fibré $V \times_U W$ dans $\Sch_\etale$
est un produit fibré dans $\Sch$ et
est le produit fibré de $V/S$ et $W/S$ au-dessus de $U/S$ dans
la catégorie de tous les schémas au-dessus de $S$, donc aussi un
produit fibré dans $(\Sch/S)_\etale$.
Cela prouve le résultat pour $(\Sch/S)_\etale$.
Si $U \to S$, $V \to U$ et $W \to U$ sont étales, alors
$V \times_U W \to S$ l'est aussi, d'où le résultat pour $S_\etale$.
Si $U, V, W$ sont affines, $V \times_U W$ l'est aussi, d'où le
résultat pour $(\textit{Aff}/S)_\etale$ et $S_{affine, \etale}$.
\end{proof}

\noindent
Vérifions ensuite que le gros, resp.\ petit, site affine définit le même
topos que le gros, resp.\ petit, site.

\begin{lemma}
\label{lemma-affine-big-site-etale}
Soit $S$ un schéma. Soit $\Sch_\etale$ un gros site
étale contenant $S$.
Le foncteur
$(\textit{Aff}/S)_\etale \to (\Sch/S)_\etale$
est cocontinu spécial et induit une équivalence de topos de
$\Sh((\textit{Aff}/S)_\etale)$ vers
$\Sh((\Sch/S)_\etale)$.
\end{lemma}

\begin{proof}
La notion de foncteur cocontinu spécial est introduite dans
Sites, définition \ref{sites-definition-special-cocontinuous-functor}.
Il faut donc vérifier les hypothèses (1) -- (5) de
Sites, lemme \ref{sites-lemma-equivalence}.
Notons le foncteur d'inclusion
$u : (\textit{Aff}/S)_\etale \to (\Sch/S)_\etale$.
La cocontinuité signifie simplement que tout recouvrement étale de
$T/S$, avec $T$ affine, admet un raffinement par un recouvrement étale standard de $T$.
C'est le contenu du
lemme \ref{lemma-etale-affine}.
Ainsi, (1) est vérifiée. Le foncteur $u$ est continu simplement parce qu'un recouvrement
étale standard est un recouvrement étale. Ainsi, (2) est vérifiée.
Les points (3) et (4) découlent immédiatement du fait que $u$ est
pleinement fidèle. Enfin, la condition (5) découle du
fait que tout schéma admet un recouvrement ouvert affine.
\end{proof}

\begin{lemma}
\label{lemma-alternative}
Soit $S$ un schéma. Soit $\Sch_\etale$ un gros site
étale contenant $S$. Le foncteur $S_{affine, \etale} \to S_\etale$
est cocontinu spécial et induit une équivalence de topos de
$\Sh(S_{affine, \etale})$ vers $\Sh(S_\etale)$.
\end{lemma}

\begin{proof}
Omis. Indication : comparer avec la démonstration du
lemme \ref{lemma-affine-big-site-etale}.
\end{proof}

\noindent
Établissons maintenant quelques relations entre les topos
associés à ces sites.

\begin{lemma}
\label{lemma-put-in-T-etale}
Soit $\Sch_\etale$ un gros site étale.
Soit $f : T \to S$ un morphisme dans $\Sch_\etale$.
Le foncteur $T_\etale \to (\Sch/S)_\etale$
est cocontinu et induit un morphisme de topos
$$
i_f :
\Sh(T_\etale)
\longrightarrow
\Sh((\Sch/S)_\etale)
$$
Pour un faisceau $\mathcal{G}$ sur $(\Sch/S)_\etale$,
on a la formule $(i_f^{-1}\mathcal{G})(U/T) = \mathcal{G}(U/S)$.
Le foncteur $i_f^{-1}$ admet aussi un adjoint à gauche $i_{f, !}$ qui commute
aux produits fibrés et aux égalisateurs.
\end{lemma}

\begin{proof}
Notons le foncteur $u : T_\etale \to (\Sch/S)_\etale$.
Autrement dit, à un morphisme étale $j : U \to T$ correspondant
à un objet de $T_\etale$, on associe $u(U \to T) = (f \circ j : U \to S)$.
Ce foncteur commute aux produits fibrés ; voir le
lemme \ref{lemma-fibre-products-etale}.
Soient $a, b : U \to V$ deux morphismes dans $T_\etale$.
Dans ce cas, l'égalisateur de $a$ et $b$ (dans la catégorie des schémas) est
$$
V \times_{\Delta_{V/T}, V \times_T V, (a, b)} U
$$
qui est un produit fibré de schémas étales au-dessus de $T$, donc étale
au-dessus de $T$. Ainsi, $T_\etale$ admet des égalisateurs et $u$ commute avec eux.
Il est clairement cocontinu.
Il est aussi continu, car $u$ envoie les recouvrements sur des recouvrements et
commute aux produits fibrés. Le lemme résulte donc de
Sites, lemmes \ref{sites-lemma-when-shriek}
et \ref{sites-lemma-preserve-equalizers}.
\end{proof}

\begin{lemma}
\label{lemma-at-the-bottom-etale}
Soit $S$ un schéma. Soit $\Sch_\etale$ un gros site
étale contenant $S$.
Le foncteur d'inclusion $S_\etale \to (\Sch/S)_\etale$
satisfait aux hypothèses de Sites, lemme \ref{sites-lemma-bigger-site},
et induit donc un morphisme de sites
$$
\pi_S : (\Sch/S)_\etale \longrightarrow S_\etale
$$
et un morphisme de topos
$$
i_S : \Sh(S_\etale) \longrightarrow \Sh((\Sch/S)_\etale)
$$
tel que $\pi_S \circ i_S = \text{id}$. De plus, $i_S = i_{\text{id}_S}$,
où $i_{\text{id}_S}$ est défini comme dans le lemme \ref{lemma-put-in-T-etale}.
En particulier, le foncteur $i_S^{-1} = \pi_{S, *}$ est décrit par la règle
$i_S^{-1}(\mathcal{G})(U/S) = \mathcal{G}(U/S)$.
\end{lemma}

\begin{proof}
Dans ce cas, le foncteur
$u : S_\etale \to (\Sch/S)_\etale$,
outre les propriétés établies dans la démonstration du
lemme \ref{lemma-put-in-T-etale} ci-dessus, est pleinement fidèle
et envoie l'objet final sur l'objet final.
Le lemme résulte de Sites, lemme \ref{sites-lemma-bigger-site}.
\end{proof}

\begin{definition}
\label{definition-restriction-small-etale}
Dans la situation du
lemme \ref{lemma-at-the-bottom-etale},
le foncteur $i_S^{-1} = \pi_{S, *}$ est souvent
appelé la {\it restriction au petit site étale} et, pour un faisceau
$\mathcal{F}$ sur le gros site étale, nous notons
$\mathcal{F}|_{S_\etale}$ cette restriction.
\end{definition}

\noindent
Avec cette notation, pour un faisceau $\mathcal{F}$ sur le
gros site et un faisceau $\mathcal{G}$ sur le petit site, on a
\begin{align*}
\Mor_{\Sh(S_\etale)}(
\mathcal{F}|_{S_\etale},
\mathcal{G})
& =
\Mor_{\Sh((\Sch/S)_\etale)}(
\mathcal{F},
i_{S, *}\mathcal{G}) \\
\Mor_{\Sh(S_\etale)}(
\mathcal{G},
\mathcal{F}|_{S_\etale})
& =
\Mor_{\Sh((\Sch/S)_\etale)}(
\pi_S^{-1}\mathcal{G},
\mathcal{F})
\end{align*}
De plus, on a $(i_{S, *}\mathcal{G})|_{S_\etale} = \mathcal{G}$
et $(\pi_S^{-1}\mathcal{G})|_{S_\etale} = \mathcal{G}$.

\begin{lemma}
\label{lemma-morphism-big-etale}
Soit $\Sch_\etale$ un gros site étale.
Soit $f : T \to S$ un morphisme dans $\Sch_\etale$.
Le foncteur
$$
u :
(\Sch/T)_\etale
\longrightarrow
(\Sch/S)_\etale,
\quad
V/T \longmapsto V/S
$$
est cocontinu et admet un adjoint à droite continu
$$
v :
(\Sch/S)_\etale
\longrightarrow
(\Sch/T)_\etale,
\quad
(U \to S) \longmapsto (U \times_S T \to T).
$$
Ils induisent le même morphisme de topos
$$
f_{big} :
\Sh((\Sch/T)_\etale)
\longrightarrow
\Sh((\Sch/S)_\etale)
$$
On a $f_{big}^{-1}(\mathcal{G})(U/T) = \mathcal{G}(U/S)$.
On a $f_{big, *}(\mathcal{F})(U/S) = \mathcal{F}(U \times_S T/T)$.
En outre, $f_{big}^{-1}$ admet un adjoint à gauche $f_{big!}$ qui commute aux
produits fibrés et aux égalisateurs.
\end{lemma}

\begin{proof}
Le foncteur $u$ est cocontinu, continu et commute aux produits fibrés
et aux égalisateurs (détails omis ; comparer à la démonstration du
lemme \ref{lemma-put-in-T-etale}).
Par conséquent,
Sites, lemmes \ref{sites-lemma-when-shriek} et
\ref{sites-lemma-preserve-equalizers}
s'appliquent et donnent la formule
pour $f_{big}^{-1}$ ainsi que l'existence de $f_{big!}$. De plus,
le foncteur $v$ est un adjoint à droite car, étant donnés $U/T$ et $V/S$,
on a $\Mor_S(u(U), V) = \Mor_T(U, V \times_S T)$,
comme voulu. On peut donc appliquer
Sites, lemmes \ref{sites-lemma-have-functor-other-way} et
\ref{sites-lemma-have-functor-other-way-morphism} pour obtenir la
formule pour $f_{big, *}$.
\end{proof}

\begin{lemma}
\label{lemma-morphism-big-small-etale}
Soit $\Sch_\etale$ un gros site étale.
Soit $f : T \to S$ un morphisme dans $\Sch_\etale$.
\begin{enumerate}
\item On a $i_f = f_{big} \circ i_T$, où $i_f$ est défini comme dans le
lemme \ref{lemma-put-in-T-etale} et $i_T$ comme dans le
lemme \ref{lemma-at-the-bottom-etale}.
\item Le foncteur $S_\etale \to T_\etale$,
$(U \to S) \mapsto (U \times_S T \to T)$, est continu et induit
un morphisme de sites
$$
f_{small} : T_\etale \longrightarrow S_\etale
$$
On a $f_{small, *}(\mathcal{F})(U/S) = \mathcal{F}(U \times_S T/T)$.
\item On a un diagramme commutatif de morphismes de sites
$$
\xymatrix{
T_\etale \ar[d]_{f_{small}} &
(\Sch/T)_\etale \ar[d]^{f_{big}} \ar[l]^{\pi_T}\\
S_\etale &
(\Sch/S)_\etale \ar[l]_{\pi_S}
}
$$
de sorte que $f_{small} \circ \pi_T = \pi_S \circ f_{big}$ comme morphismes de topos.
\item On a $f_{small} = \pi_S \circ f_{big} \circ i_T = \pi_S \circ i_f$.
\end{enumerate}
\end{lemma}

\begin{proof}
L'égalité $i_f = f_{big} \circ i_T$ résulte de
l'égalité $i_f^{-1} = i_T^{-1} \circ f_{big}^{-1}$, qui découle
des descriptions de ces foncteurs données ci-dessus.
Cela prouve (1).

\medskip\noindent
Le foncteur
$u :
S_\etale
\to
T_\etale$, $u(U \to S) = (U \times_S T \to T)$
envoie les recouvrements sur des recouvrements et commute aux produits fibrés ;
voir le lemme \ref{lemma-etale} (3) et \ref{lemma-fibre-products-etale}.
De plus, $S_\etale$ et $T_\etale$ admettent tous deux des objets finaux,
à savoir $S/S$ et $T/T$, et $u(S/S) = T/T$. Par conséquent, d'après
Sites, proposition \ref{sites-proposition-get-morphism},
le foncteur $u$ correspond à un morphisme de sites
$T_\etale \to S_\etale$. Celui-ci donne à son tour naissance au
morphisme de topos ; voir
Sites, lemme \ref{sites-lemma-morphism-sites-topoi}. La description
de l'image directe résulte immédiatement de ces références.

\medskip\noindent
Le point (3) résulte du fait que $\pi_S$ et $\pi_T$ sont donnés par les
foncteurs d'inclusion, et $f_{small}$ et $f_{big}$ par les
foncteurs de changement de base $U \mapsto U \times_S T$.

\medskip\noindent
L'assertion (4) résulte de (3) en précomposant avec $i_T$.
\end{proof}

\noindent
Dans la situation du lemme, avec la terminologie de la
définition \ref{definition-restriction-small-etale},
on a, pour $\mathcal{F}$ faisceau sur le gros site étale de $T$,
$$
(f_{big, *}\mathcal{F})|_{S_\etale} =
f_{small, *}(\mathcal{F}|_{T_\etale}),
$$
Cette égalité résulte de la commutativité du diagramme de
sites du lemme, puisque la restriction au petit site étale de
$T$, resp.\ $S$, est donnée par $\pi_{T, *}$, resp.\ $\pi_{S, *}$. Une formule
analogue faisant intervenir images inverses et restrictions est fausse.

\begin{lemma}
\label{lemma-composition-etale}
Étant donnés des schémas $X$, $Y$, $Y$ de $\Sch_\etale$
et des morphismes $f : X \to Y$, $g : Y \to Z$, on a
$g_{big} \circ f_{big} = (g \circ f)_{big}$ et
$g_{small} \circ f_{small} = (g \circ f)_{small}$.
\end{lemma}

\begin{proof}
Cela résulte de la description simple de l'image directe
et de l'image inverse pour les foncteurs sur les gros sites donnée par le
lemme \ref{lemma-morphism-big-etale}. Pour les foncteurs
sur les petits sites, cela résulte de la description des
foncteurs image directe dans le lemme \ref{lemma-morphism-big-small-etale}.
\end{proof}

\begin{lemma}
\label{lemma-morphism-big-small-cartesian-diagram-etale}
Soit $\Sch_\etale$ un gros site étale. Considérons un diagramme cartésien
$$
\xymatrix{
T' \ar[r]_{g'} \ar[d]_{f'} & T \ar[d]^f \\
S' \ar[r]^g & S
}
$$
dans $\Sch_\etale$. Alors
$i_g^{-1} \circ f_{big, *} = f'_{small, *} \circ (i_{g'})^{-1}$
et $g_{big}^{-1} \circ f_{big, *} = f'_{big, *} \circ (g'_{big})^{-1}$.
\end{lemma}

\begin{proof}
Le diagramme étant cartésien, on a, pour $U'/S'$,
$U' \times_{S'} T' = U' \times_S T$. Par conséquent, les deux foncteurs
$i_g^{-1} \circ f_{big, *}$ et $f'_{small, *} \circ (i_{g'})^{-1}$
envoient un faisceau $\mathcal{F}$ sur $(\Sch/T)_\etale$ sur le faisceau
$U' \mapsto \mathcal{F}(U' \times_{S'} T')$ sur $S'_\etale$
(utiliser les lemmes \ref{lemma-put-in-T-etale} et \ref{lemma-morphism-big-etale}).
La seconde égalité peut se démontrer de la même manière ou se
déduire du très général
Sites, lemme \ref{sites-lemma-localize-morphism}.
\end{proof}

\noindent
On peut considérer un faisceau sur le gros site étale de $S$ comme une collection
de faisceaux « usuels » sur tous les schémas au-dessus de $S$.

\begin{lemma}
\label{lemma-characterize-sheaf-big-etale}
Soit $S$ un schéma contenu dans un gros site étale
$\Sch_\etale$.
Un faisceau $\mathcal{F}$ sur le gros site étale
$(\Sch/S)_\etale$ est décrit par les données suivantes :
\begin{enumerate}
\item pour tout $T/S \in \Ob((\Sch/S)_\etale)$, un faisceau
$\mathcal{F}_T$ sur $T_\etale$,
\item pour tout $f : T' \to T$ dans
$(\Sch/S)_\etale$, une application
$c_f : f_{small}^{-1}\mathcal{F}_T \to \mathcal{F}_{T'}$.
\end{enumerate}
Ces données sont soumises aux conditions suivantes :
\begin{enumerate}
\item[(a)] pour tous $f : T' \to T$ et $g : T'' \to T'$ dans
$(\Sch/S)_\etale$, la composée
$c_g \circ g_{small}^{-1}c_f$ est égale à $c_{f \circ g}$, et
\item[(b)] si $f : T' \to T$ dans $(\Sch/S)_\etale$
est étale, alors $c_f$ est un isomorphisme.
\end{enumerate}
\end{lemma}

\begin{proof}
Ce lemme résulte d'un énoncé purement faisceautique discuté dans
Sites, remarque \ref{sites-remark-variant-glue-sheaves-absolute}.
Nous donnons aussi une démonstration directe dans ce cas.

\medskip\noindent
Étant donné un faisceau $\mathcal{F}$ sur $\Sh((\Sch/S)_\etale)$,
on pose $\mathcal{F}_T = i_p^{-1}\mathcal{F}$, où $p : T \to S$
est le morphisme structural. Remarquons que
$\mathcal{F}_T(U) = \mathcal{F}(U/S)$ pour tout $U \to T$
dans $T_\etale$ ; voir le lemme \ref{lemma-put-in-T-etale}.
Ainsi, étant donnés $f : T' \to T$ au-dessus de $S$ et $U \to T$, on obtient une
application canonique $\mathcal{F}_T(U) = \mathcal{F}(U/S) \to \mathcal{F}(U \times_T T'/S)
= \mathcal{F}_{T'}(U \times_T T')$, où l'application médiane est l'application de restriction
de $\mathcal{F}$ relativement au morphisme
$U \times_T T' \to U$ au-dessus de $S$. La famille de ces applications est
compatible aux restrictions et définit donc une application
$c'_f : \mathcal{F}_T \to f_{small, *}\mathcal{F}_{T'}$, où
$u : T_\etale \to T'_\etale$ est le foncteur de changement de base
associé à $f$. Par l'adjonction de $f_{small, *}$ (voir
Sites, section \ref{sites-section-continuous-functors}) avec
$f_{small}^{-1}$, cela revient à une application
$c_f : f_{small}^{-1}\mathcal{F}_T \to \mathcal{F}_{T'}$.
Il est clair que $c'_{f \circ g}$ est la composée de
$c'_f$ et $f_{small, *}c'_g$, puisque la composée d'applications de restriction
de $\mathcal{F}$ est encore une application de restriction ; on obtient ainsi la
relation voulue entre $c_f$, $c_g$ et $c_{f \circ g}$.

\medskip\noindent
Réciproquement, étant donné un système $(\mathcal{F}_T, c_f)$ comme dans le lemme,
on peut définir un préfaisceau $\mathcal{F}$ sur
$\Sh((\Sch/S)_\etale)$
en posant simplement $\mathcal{F}(T/S) = \mathcal{F}_T(T)$. Comme application de restriction,
étant donné $f : T' \to T$, on définit, pour $s \in \mathcal{F}(T)$,
l'image inverse $f^*(s)$ comme $c_f(s)$, où l'on considère de nouveau $c_f$ comme
une application $\mathcal{F}_T \to f_{small, *}\mathcal{F}_{T'}$.
La condition sur les $c_f$ garantit que
les images inverses satisfont à la fonctorialité requise.
Nous omettons la vérification qu'il s'agit d'un faisceau.
Il est clair que les constructions ainsi définies sont inverses l'une de l'autre.
\end{proof}























\section{La topologie lisse}
\label{section-smooth}

\noindent
Dans cette section, nous définissons la topologie lisse.
Cela est quelque peu gratuit, puisqu'on verra plus loin (voir
Compléments sur les morphismes, section \ref{more-morphisms-section-etale-over-smooth})
que cette topologie définit le même topos que la
topologie étale. Elle demeure néanmoins naturelle et intervient
à l'occasion.

\begin{definition}
\label{definition-smooth-covering}
Soit $T$ un schéma. Un {\it recouvrement lisse de $T$} est une famille
de morphismes de schémas $\{f_i : T_i \to T\}_{i \in I}$
telle que chaque $f_i$ soit lisse et
que $T = \bigcup f_i(T_i)$.
\end{definition}

\begin{lemma}
\label{lemma-zariski-etale-smooth}
Tout recouvrement étale est un recouvrement lisse et, à plus forte raison,
tout recouvrement de Zariski est un recouvrement lisse.
\end{lemma}

\begin{proof}
Cela résulte des définitions et du fait qu'un morphisme étale est
lisse ; voir
Morphismes, définition \ref{morphisms-definition-etale}
et le lemme \ref{lemma-zariski-etale}.
\end{proof}

\noindent
Montrons maintenant que cette notion satisfait aux conditions de
Sites, définition \ref{sites-definition-site}.

\begin{lemma}
\label{lemma-smooth}
Soit $T$ un schéma.
\begin{enumerate}
\item Si $T' \to T$ est un isomorphisme, alors $\{T' \to T\}$
est un recouvrement lisse de $T$.
\item Si $\{T_i \to T\}_{i\in I}$ est un recouvrement lisse et si, pour tout
$i$, on a un recouvrement lisse $\{T_{ij} \to T_i\}_{j\in J_i}$, alors
$\{T_{ij} \to T\}_{i \in I, j\in J_i}$ est un recouvrement lisse.
\item Si $\{T_i \to T\}_{i\in I}$ est un recouvrement lisse
et si $T' \to T$ est un morphisme de schémas, alors
$\{T' \times_T T_i \to T'\}_{i\in I}$ est un recouvrement lisse.
\end{enumerate}
\end{lemma}

\begin{proof}
Omis.
\end{proof}

\begin{lemma}
\label{lemma-smooth-affine}
Soit $T$ un schéma affine.
Soit $\{T_i \to T\}_{i \in I}$ un recouvrement lisse de $T$.
Il existe alors un recouvrement lisse
$\{U_j \to T\}_{j = 1, \ldots, m}$ qui raffine
$\{T_i \to T\}_{i \in I}$, tel que chaque $U_j$ soit un schéma
affine et que chaque morphisme $U_j \to T$ soit lisse
standard ; voir Morphismes, définition \ref{morphisms-definition-smooth}.
On peut en outre choisir chaque $U_j$ comme ouvert affine de l'un des $T_i$.
\end{lemma}

\begin{proof}
Omis ; voir toutefois Algèbre, lemme \ref{algebra-lemma-smooth-syntomic}.
\end{proof}

\noindent
On définit donc comme suit les recouvrements standards d'affines correspondants.

\begin{definition}
\label{definition-standard-smooth}
Soit $T$ un schéma affine. Un {\it recouvrement lisse standard}
de $T$ est une famille $\{f_j : U_j \to T\}_{j = 1, \ldots, m}$
telle que chaque $U_j$ soit affine, que $U_j \to T$ soit lisse standard
et que $T = \bigcup f_j(U_j)$.
\end{definition}

\begin{definition}
\label{definition-big-smooth-site}
Un {\it gros site lisse} est un site $\Sch_{smooth}$ au sens de
Sites, définition \ref{sites-definition-site}, construit comme suit :
\begin{enumerate}
\item Choisir un ensemble quelconque de schémas $S_0$ et un ensemble de recouvrements lisses
$\text{Cov}_0$ entre ces schémas.
\item Prendre comme catégorie sous-jacente une catégorie $\Sch_\alpha$
construite comme dans Ensembles, lemme \ref{sets-lemma-construct-category},
à partir de l'ensemble $S_0$.
\item Choisir un ensemble de recouvrements comme dans
Ensembles, lemme \ref{sets-lemma-coverings-site}, à partir de la
catégorie $\Sch_\alpha$, de la classe des recouvrements lisses
et de l'ensemble $\text{Cov}_0$ choisi ci-dessus.
\end{enumerate}
\end{definition}

\noindent
Voir les remarques suivant la définition \ref{definition-big-zariski-site}
pour la motivation et les explications relatives à la définition des gros sites.

\medskip\noindent
Avant de poursuivre par l'introduction du gros site lisse d'un
schéma $S$, observons que la topologie d'un gros site lisse
$\Sch_{smooth}$ est, en un certain sens, induite par la topologie lisse
sur la catégorie de tous les schémas.

\begin{lemma}
\label{lemma-smooth-induced}
Soit $\Sch_{smooth}$ un gros site lisse comme dans la
définition \ref{definition-big-smooth-site}.
Soit $T \in \Ob(\Sch_{smooth})$.
Soit $\{T_i \to T\}_{i \in I}$ un recouvrement lisse quelconque de $T$.
\begin{enumerate}
\item Il existe un recouvrement $\{U_j \to T\}_{j \in J}$ de $T$ dans le site
$\Sch_{smooth}$ qui raffine $\{T_i \to T\}_{i \in I}$.
\item Si $\{T_i \to T\}_{i \in I}$ est un recouvrement lisse standard, alors
il est tautologiquement équivalent à un recouvrement de $\Sch_{smooth}$.
\item Si $\{T_i \to T\}_{i \in I}$ est un recouvrement de Zariski, alors
il est tautologiquement équivalent à un recouvrement de $\Sch_{smooth}$.
\end{enumerate}
\end{lemma}

\begin{proof}
Pour chaque $i$, choisissons un recouvrement ouvert affine $T_i = \bigcup_{j \in J_i} T_{ij}$
tel que chaque $T_{ij}$ s'envoie dans un sous-schéma ouvert affine de $T$. D'après le
lemme \ref{lemma-smooth},
le raffinement $\{T_{ij} \to T\}_{i \in I, j \in J_i}$ est lui aussi un recouvrement lisse
de $T$. On peut donc supposer chaque $T_i$ affine et s'envoyant dans
un ouvert affine $W_i$ de $T$. En appliquant
Ensembles, lemme \ref{sets-lemma-what-is-in-it},
on voit que $W_i$ est isomorphe à un objet de $\Sch_{smooth}$.
Mais alors $T_i$, comme schéma de type fini au-dessus de $W_i$,
est isomorphe à un objet $V_i$ de $\Sch_{smooth}$ par une seconde
application de
Ensembles, lemme \ref{sets-lemma-what-is-in-it}.
Le recouvrement $\{V_i \to T\}_{i \in I}$ raffine $\{T_i \to T\}_{i \in I}$
(puisqu'ils sont isomorphes).
De plus, $\{V_i \to T\}_{i \in I}$ est combinatoirement équivalent à un
recouvrement $\{U_j \to T\}_{j \in J}$ de $T$ dans le site
$\Sch_{smooth}$ d'après
Ensembles, lemme \ref{sets-lemma-what-is-in-it}.
Le recouvrement $\{U_j \to T\}_{j \in J}$ est un raffinement comme en (1).
Dans la situation de (2), (3), chacun des
schémas $T_i$ est isomorphe à un objet de $\Sch_{smooth}$ d'après
Ensembles, lemme \ref{sets-lemma-what-is-in-it},
et une nouvelle application de
Ensembles, lemme \ref{sets-lemma-coverings-site},
donne le résultat voulu.
\end{proof}

\begin{definition}
\label{definition-big-small-smooth}
Soit $S$ un schéma. Soit $\Sch_{smooth}$ un gros site
lisse contenant $S$.
\begin{enumerate}
\item Le {\it gros site lisse de $S$}, noté
$(\Sch/S)_{smooth}$, est le site $\Sch_{smooth}/S$
introduit dans Sites, section \ref{sites-section-localize}.
\item Le {\it gros site lisse affine de $S$}, noté
$(\textit{Aff}/S)_{smooth}$, est la sous-catégorie pleine de
$(\Sch/S)_{smooth}$ dont les objets sont les $U/S$ affines.
Un recouvrement de $(\textit{Aff}/S)_{smooth}$ est tout recouvrement
$\{U_i \to U\}$ de $(\Sch/S)_{smooth}$ qui soit un
recouvrement lisse standard.
\end{enumerate}
\end{definition}

\noindent
Vérifions ensuite que le gros site affine définit le même
topos que le gros site.

\begin{lemma}
\label{lemma-affine-big-site-smooth}
Soit $S$ un schéma. Soit $\Sch_{smooth}$ un gros site
lisse contenant $S$.
Le foncteur
$(\textit{Aff}/S)_{smooth} \to (\Sch/S)_{smooth}$
est cocontinu spécial et induit une équivalence de topos de
$\Sh((\textit{Aff}/S)_{smooth})$ vers
$\Sh((\Sch/S)_{smooth})$.
\end{lemma}

\begin{proof}
La notion de foncteur cocontinu spécial est introduite dans
Sites, définition \ref{sites-definition-special-cocontinuous-functor}.
Il faut donc vérifier les hypothèses (1) -- (5) de
Sites, lemme \ref{sites-lemma-equivalence}.
Notons le foncteur d'inclusion
$u : (\textit{Aff}/S)_{smooth} \to (\Sch/S)_{smooth}$.
La cocontinuité signifie simplement que tout recouvrement lisse de
$T/S$, avec $T$ affine, admet un raffinement par un recouvrement lisse standard de $T$.
C'est le contenu du
lemme \ref{lemma-smooth-affine}.
Ainsi, (1) est vérifiée. Le foncteur $u$ est continu simplement parce qu'un recouvrement
lisse standard est un recouvrement lisse. Ainsi, (2) est vérifiée.
Les points (3) et (4) découlent immédiatement du fait que $u$ est
pleinement fidèle. Enfin, la condition (5) découle du
fait que tout schéma admet un recouvrement ouvert affine.
\end{proof}

\noindent
À suivre...

\begin{lemma}
\label{lemma-morphism-big-smooth}
Soit $\Sch_{smooth}$ un gros site lisse.
Soit $f : T \to S$ un morphisme dans $\Sch_{smooth}$.
Le foncteur
$$
u : (\Sch/T)_{smooth} \longrightarrow (\Sch/S)_{smooth},
\quad
V/T \longmapsto V/S
$$
est cocontinu et admet un adjoint à droite continu
$$
v : (\Sch/S)_{smooth} \longrightarrow (\Sch/T)_{smooth},
\quad
(U \to S) \longmapsto (U \times_S T \to T).
$$
Ils induisent le même morphisme de topos
$$
f_{big} :
\Sh((\Sch/T)_{smooth})
\longrightarrow
\Sh((\Sch/S)_{smooth})
$$
On a $f_{big}^{-1}(\mathcal{G})(U/T) = \mathcal{G}(U/S)$.
On a $f_{big, *}(\mathcal{F})(U/S) = \mathcal{F}(U \times_S T/T)$.
En outre, $f_{big}^{-1}$ admet un adjoint à gauche $f_{big!}$ qui commute aux
produits fibrés et aux égalisateurs.
\end{lemma}

\begin{proof}
Le foncteur $u$ est cocontinu, continu et commute aux produits fibrés
et aux égalisateurs. Par conséquent,
Sites, lemmes \ref{sites-lemma-when-shriek} et
\ref{sites-lemma-preserve-equalizers}
s'appliquent et donnent la formule
pour $f_{big}^{-1}$ ainsi que l'existence de $f_{big!}$. De plus,
le foncteur $v$ est un adjoint à droite car, étant donnés $U/T$ et $V/S$,
on a $\Mor_S(u(U), V) = \Mor_T(U, V \times_S T)$,
comme voulu. On peut donc appliquer
Sites, lemmes \ref{sites-lemma-have-functor-other-way} et
\ref{sites-lemma-have-functor-other-way-morphism} pour obtenir la
formule pour $f_{big, *}$.
\end{proof}











\section{La topologie syntomique}
\label{section-syntomic}

\noindent
Dans cette section, nous définissons la topologie syntomique.
Cette topologie est assez intéressante : elle possède souvent
les mêmes groupes de cohomologie que la topologie fppf,
tout en étant techniquement plus facile à manier.

\begin{definition}
\label{definition-syntomic-covering}
Soit $T$ un schéma. Un {\it recouvrement syntomique de $T$} est une famille
de morphismes de schémas $\{f_i : T_i \to T\}_{i \in I}$
telle que chaque $f_i$ soit syntomique et
que $T = \bigcup f_i(T_i)$.
\end{definition}

\begin{lemma}
\label{lemma-zariski-etale-smooth-syntomic}
Tout recouvrement lisse est un recouvrement syntomique et, à plus forte raison,
tout recouvrement étale ou de Zariski est un recouvrement syntomique.
\end{lemma}

\begin{proof}
Cela résulte des définitions et du fait qu'un
morphisme lisse est syntomique ; voir
Morphismes, lemme \ref{morphisms-lemma-smooth-syntomic}
et le lemme \ref{lemma-zariski-etale-smooth}.
\end{proof}

\noindent
Montrons maintenant que cette notion satisfait aux conditions de
Sites, définition \ref{sites-definition-site}.

\begin{lemma}
\label{lemma-syntomic}
Soit $T$ un schéma.
\begin{enumerate}
\item Si $T' \to T$ est un isomorphisme, alors $\{T' \to T\}$
est un recouvrement syntomique de $T$.
\item Si $\{T_i \to T\}_{i\in I}$ est un recouvrement syntomique et si, pour tout
$i$, on a un recouvrement syntomique $\{T_{ij} \to T_i\}_{j\in J_i}$, alors
$\{T_{ij} \to T\}_{i \in I, j\in J_i}$ est un recouvrement syntomique.
\item Si $\{T_i \to T\}_{i\in I}$ est un recouvrement syntomique
et si $T' \to T$ est un morphisme de schémas, alors
$\{T' \times_T T_i \to T'\}_{i\in I}$ est un recouvrement syntomique.
\end{enumerate}
\end{lemma}

\begin{proof}
Omis.
\end{proof}

\begin{lemma}
\label{lemma-syntomic-affine}
Soit $T$ un schéma affine.
Soit $\{T_i \to T\}_{i \in I}$ un recouvrement syntomique de $T$.
Il existe alors un recouvrement syntomique
$\{U_j \to T\}_{j = 1, \ldots, m}$ qui raffine
$\{T_i \to T\}_{i \in I}$, tel que chaque $U_j$ soit un schéma
affine et que chaque morphisme $U_j \to T$ soit syntomique
standard ; voir Morphismes, définition \ref{morphisms-definition-syntomic}.
On peut en outre choisir chaque $U_j$ comme ouvert affine de l'un des $T_i$.
\end{lemma}

\begin{proof}
Omis ; voir toutefois Algèbre, lemme \ref{algebra-lemma-syntomic}.
\end{proof}

\noindent
On définit donc comme suit les recouvrements standards d'affines correspondants.

\begin{definition}
\label{definition-standard-syntomic}
Soit $T$ un schéma affine. Un {\it recouvrement syntomique standard} de $T$ est
une famille $\{f_j : U_j \to T\}_{j = 1, \ldots, m}$ telle que chaque $U_j$ soit
affine, que $U_j \to T$ soit syntomique standard et que $T = \bigcup f_j(U_j)$.
\end{definition}

\begin{definition}
\label{definition-big-syntomic-site}
Un {\it gros site syntomique} est un site $\Sch_{syntomic}$ au sens de
Sites, définition \ref{sites-definition-site}, construit comme suit :
\begin{enumerate}
\item Choisir un ensemble quelconque de schémas $S_0$ et un ensemble de recouvrements syntomiques
$\text{Cov}_0$ entre ces schémas.
\item Prendre comme catégorie sous-jacente une catégorie $\Sch_\alpha$
construite comme dans Ensembles, lemme \ref{sets-lemma-construct-category},
à partir de l'ensemble $S_0$.
\item Choisir un ensemble de recouvrements comme dans
Ensembles, lemme \ref{sets-lemma-coverings-site}, à partir de la
catégorie $\Sch_\alpha$, de la classe des recouvrements syntomiques
et de l'ensemble $\text{Cov}_0$ choisi ci-dessus.
\end{enumerate}
\end{definition}

\noindent
Voir les remarques suivant la définition \ref{definition-big-zariski-site}
pour la motivation et les explications relatives à la définition des gros sites.

\medskip\noindent
Avant de poursuivre par l'introduction du gros site syntomique d'un
schéma $S$, observons que la topologie d'un gros site syntomique
$\Sch_{syntomic}$ est, en un certain sens, induite par la topologie syntomique
sur la catégorie de tous les schémas.

\begin{lemma}
\label{lemma-syntomic-induced}
Soit $\Sch_{syntomic}$ un gros site syntomique comme dans la
définition \ref{definition-big-syntomic-site}.
Soit $T \in \Ob(\Sch_{syntomic})$.
Soit $\{T_i \to T\}_{i \in I}$ un recouvrement syntomique quelconque de $T$.
\begin{enumerate}
\item Il existe un recouvrement $\{U_j \to T\}_{j \in J}$ de $T$ dans le site
$\Sch_{syntomic}$ qui raffine $\{T_i \to T\}_{i \in I}$.
\item Si $\{T_i \to T\}_{i \in I}$ est un recouvrement syntomique standard, alors
il est tautologiquement équivalent à un recouvrement dans $\Sch_{syntomic}$.
\item Si $\{T_i \to T\}_{i \in I}$ est un recouvrement de Zariski, alors
il est tautologiquement équivalent à un recouvrement dans $\Sch_{syntomic}$.
\end{enumerate}
\end{lemma}

\begin{proof}
Pour chaque $i$, choisissons un recouvrement ouvert affine $T_i = \bigcup_{j \in J_i} T_{ij}$
tel que chaque $T_{ij}$ s'envoie dans un sous-schéma ouvert affine de $T$. D'après le
lemme \ref{lemma-syntomic},
le raffinement $\{T_{ij} \to T\}_{i \in I, j \in J_i}$ est lui aussi un recouvrement syntomique
de $T$. On peut donc supposer chaque $T_i$ affine et s'envoyant dans
un ouvert affine $W_i$ de $T$. En appliquant
Ensembles, lemme \ref{sets-lemma-what-is-in-it},
on voit que $W_i$ est isomorphe à un objet de $\Sch_{syntomic}$.
Mais alors $T_i$, comme schéma de type fini au-dessus de $W_i$,
est isomorphe à un objet $V_i$ de $\Sch_{syntomic}$ par une seconde
application de
Ensembles, lemme \ref{sets-lemma-what-is-in-it}.
Le recouvrement $\{V_i \to T\}_{i \in I}$ raffine $\{T_i \to T\}_{i \in I}$
(puisqu'ils sont isomorphes).
De plus, $\{V_i \to T\}_{i \in I}$ est combinatoirement équivalent à un
recouvrement $\{U_j \to T\}_{j \in J}$ de $T$ dans le site
$\Sch_{syntomic}$ d'après
Ensembles, lemme \ref{sets-lemma-what-is-in-it}.
Le recouvrement $\{U_j \to T\}_{j \in J}$ est un recouvrement comme en (1).
Dans la situation de (2), (3), chacun des
schémas $T_i$ est isomorphe à un objet de $\Sch_{syntomic}$ d'après
Ensembles, lemme \ref{sets-lemma-what-is-in-it},
et une nouvelle application de
Ensembles, lemme \ref{sets-lemma-coverings-site},
donne le résultat voulu.
\end{proof}

\begin{definition}
\label{definition-big-small-syntomic}
Soit $S$ un schéma. Soit $\Sch_{syntomic}$ un gros site
syntomique contenant $S$.
\begin{enumerate}
\item Le {\it gros site syntomique de $S$}, noté
$(\Sch/S)_{syntomic}$, est le site $\Sch_{syntomic}/S$
introduit dans Sites, section \ref{sites-section-localize}.
\item Le {\it gros site syntomique affine de $S$}, noté
$(\textit{Aff}/S)_{syntomic}$, est la sous-catégorie pleine de
$(\Sch/S)_{syntomic}$ dont les objets sont les $U/S$ affines.
Un recouvrement de $(\textit{Aff}/S)_{syntomic}$ est tout recouvrement
$\{U_i \to U\}$ de $(\Sch/S)_{syntomic}$ qui soit un
recouvrement syntomique standard.
\end{enumerate}
\end{definition}

\noindent
Vérifions ensuite que le gros site affine définit le même
topos que le gros site.

\begin{lemma}
\label{lemma-affine-big-site-syntomic}
Soit $S$ un schéma. Soit $\Sch_{syntomic}$ un gros site
syntomique contenant $S$.
Le foncteur
$(\textit{Aff}/S)_{syntomic} \to (\Sch/S)_{syntomic}$
est cocontinu spécial et induit une équivalence de topos de
$\Sh((\textit{Aff}/S)_{syntomic})$ vers
$\Sh((\Sch/S)_{syntomic})$.
\end{lemma}

\begin{proof}
La notion de foncteur cocontinu spécial est introduite dans
Sites, définition \ref{sites-definition-special-cocontinuous-functor}.
Il faut donc vérifier les hypothèses (1) -- (5) de
Sites, lemme \ref{sites-lemma-equivalence}.
Notons le foncteur d'inclusion
$u : (\textit{Aff}/S)_{syntomic} \to (\Sch/S)_{syntomic}$.
La cocontinuité signifie simplement que tout recouvrement syntomique de
$T/S$, avec $T$ affine, admet un raffinement par un recouvrement syntomique standard de $T$.
C'est le contenu du
lemme \ref{lemma-syntomic-affine}.
Ainsi, (1) est vérifiée. Le foncteur $u$ est continu simplement parce qu'un recouvrement
syntomique standard est un recouvrement syntomique. Ainsi, (2) est vérifiée.
Les points (3) et (4) découlent immédiatement du fait que $u$ est
pleinement fidèle. Enfin, la condition (5) découle du
fait que tout schéma admet un recouvrement ouvert affine.
\end{proof}

\noindent
À suivre...

\begin{lemma}
\label{lemma-morphism-big-syntomic}
Soit $\Sch_{syntomic}$ un gros site syntomique.
Soit $f : T \to S$ un morphisme dans $\Sch_{syntomic}$.
Le foncteur
$$
u : (\Sch/T)_{syntomic} \longrightarrow (\Sch/S)_{syntomic},
\quad
V/T \longmapsto V/S
$$
est cocontinu et admet un adjoint à droite continu
$$
v : (\Sch/S)_{syntomic} \longrightarrow (\Sch/T)_{syntomic},
\quad
(U \to S) \longmapsto (U \times_S T \to T).
$$
Ils induisent le même morphisme de topos
$$
f_{big} :
\Sh((\Sch/T)_{syntomic})
\longrightarrow
\Sh((\Sch/S)_{syntomic})
$$
On a $f_{big}^{-1}(\mathcal{G})(U/T) = \mathcal{G}(U/S)$.
On a $f_{big, *}(\mathcal{F})(U/S) = \mathcal{F}(U \times_S T/T)$.
En outre, $f_{big}^{-1}$ admet un adjoint à gauche $f_{big!}$ qui commute aux
produits fibrés et aux égalisateurs.
\end{lemma}

\begin{proof}
Le foncteur $u$ est cocontinu, continu et commute aux produits fibrés
et aux égalisateurs. Par conséquent,
Sites, lemmes \ref{sites-lemma-when-shriek} et
\ref{sites-lemma-preserve-equalizers}
s'appliquent et donnent la formule
pour $f_{big}^{-1}$ ainsi que l'existence de $f_{big!}$. De plus,
le foncteur $v$ est un adjoint à droite car, étant donnés $U/T$ et $V/S$,
on a $\Mor_S(u(U), V) = \Mor_T(U, V \times_S T)$,
comme voulu. On peut donc appliquer
Sites, lemmes \ref{sites-lemma-have-functor-other-way} et
\ref{sites-lemma-have-functor-other-way-morphism} pour obtenir la
formule pour $f_{big, *}$.
\end{proof}













\section{La topologie fppf}
\label{section-fppf}

\noindent
Soit $S$ un schéma. Nous voulons définir la topologie fppf\footnote{
Les lettres fppf signifient « fid\`element plat de pr\'esentation finie ».} sur
la catégorie des schémas au-dessus de $S$. Conformément à notre principe général,
nous introduisons d'abord la notion de recouvrement fppf.

\begin{definition}
\label{definition-fppf-covering}
Soit $T$ un schéma. Un {\it recouvrement fppf de $T$} est une famille
de morphismes de schémas $\{f_i : T_i \to T\}_{i \in I}$
telle que chaque $f_i$ soit plat, localement de présentation finie, et
que $T = \bigcup f_i(T_i)$.
\end{definition}

\begin{lemma}
\label{lemma-zariski-etale-smooth-syntomic-fppf}
Tout recouvrement syntomique est un recouvrement fppf et, à plus forte raison,
tout recouvrement lisse, étale ou de Zariski est un recouvrement fppf.
\end{lemma}

\begin{proof}
Cela résulte des définitions et du fait qu'un morphisme syntomique
est plat et localement de présentation finie ; voir
Morphismes, lemmes
\ref{morphisms-lemma-syntomic-locally-finite-presentation} et
\ref{morphisms-lemma-syntomic-flat},
ainsi que
le lemme \ref{lemma-zariski-etale-smooth-syntomic}.
\end{proof}

\noindent
Montrons maintenant que cette notion satisfait aux conditions de
Sites, définition \ref{sites-definition-site}.

\begin{lemma}
\label{lemma-fppf}
Soit $T$ un schéma.
\begin{enumerate}
\item Si $T' \to T$ est un isomorphisme, alors $\{T' \to T\}$
est un recouvrement fppf de $T$.
\item Si $\{T_i \to T\}_{i\in I}$ est un recouvrement fppf et si, pour tout
$i$, on a un recouvrement fppf $\{T_{ij} \to T_i\}_{j\in J_i}$, alors
$\{T_{ij} \to T\}_{i \in I, j\in J_i}$ est un recouvrement fppf.
\item Si $\{T_i \to T\}_{i\in I}$ est un recouvrement fppf
et si $T' \to T$ est un morphisme de schémas, alors
$\{T' \times_T T_i \to T'\}_{i\in I}$ est un recouvrement fppf.
\end{enumerate}
\end{lemma}

\begin{proof}
La première assertion est claire.
La deuxième résulte de ce que la composée de morphismes plats est plate
(voir Morphismes, lemme \ref{morphisms-lemma-composition-flat})
et que la composée de morphismes de présentation finie est
de présentation finie
(voir Morphismes, lemme \ref{morphisms-lemma-composition-finite-presentation}).
La troisième résulte de ce que le changement de base d'un morphisme plat est plat
(voir Morphismes, lemme \ref{morphisms-lemma-base-change-flat})
et que le changement de base d'un morphisme de présentation finie est
de présentation finie
(voir Morphismes, lemme \ref{morphisms-lemma-base-change-finite-presentation}).
De plus, le changement de base d'une famille surjective de morphismes est surjectif
(démonstration omise).
\end{proof}

\begin{lemma}
\label{lemma-fppf-affine}
Soit $T$ un schéma affine.
Soit $\{T_i \to T\}_{i \in I}$ un recouvrement fppf de $T$.
Il existe alors un recouvrement fppf
$\{U_j \to T\}_{j = 1, \ldots, m}$ qui raffine
$\{T_i \to T\}_{i \in I}$ et tel que chaque $U_j$ soit un schéma
affine. On peut en outre choisir chaque $U_j$ comme ouvert affine
de l'un des $T_i$.
\end{lemma}

\begin{proof}
Cela résulte directement des définitions, puisqu'un
morphisme plat et localement de présentation finie est ouvert ;
voir Morphismes, lemme \ref{morphisms-lemma-fppf-open}.
\end{proof}

\noindent
On définit donc comme suit les recouvrements standards d'affines correspondants.

\begin{definition}
\label{definition-standard-fppf}
Soit $T$ un schéma affine. Un {\it recouvrement fppf standard}
de $T$ est une famille $\{f_j : U_j \to T\}_{j = 1, \ldots, m}$
telle que chaque $U_j$ soit affine, plat et de présentation finie sur $T$,
et que $T = \bigcup f_j(U_j)$.
\end{definition}

\begin{definition}
\label{definition-big-fppf-site}
Un {\it gros site fppf} est un site $\Sch_{fppf}$ au sens de
Sites, définition \ref{sites-definition-site}, construit comme suit :
\begin{enumerate}
\item Choisir un ensemble quelconque de schémas $S_0$ et un ensemble de recouvrements fppf
$\text{Cov}_0$ entre ces schémas.
\item Prendre comme catégorie sous-jacente une catégorie $\Sch_\alpha$
construite comme dans Ensembles, lemme \ref{sets-lemma-construct-category},
à partir de l'ensemble $S_0$.
\item Choisir un ensemble de recouvrements comme dans
Ensembles, lemme \ref{sets-lemma-coverings-site}, à partir de la
catégorie $\Sch_\alpha$, de la classe des recouvrements fppf
et de l'ensemble $\text{Cov}_0$ choisi ci-dessus.
\end{enumerate}
\end{definition}

\noindent
Voir les remarques suivant la définition \ref{definition-big-zariski-site}
pour la motivation et les explications relatives à la définition des gros sites.

\medskip\noindent
Avant de poursuivre par l'introduction du gros site fppf d'un
schéma $S$, observons que la topologie d'un gros site fppf
$\Sch_{fppf}$ est, en un certain sens, induite par la topologie fppf
sur la catégorie de tous les schémas.

\begin{lemma}
\label{lemma-fppf-induced}
Soit $\Sch_{fppf}$ un gros site fppf comme dans la
définition \ref{definition-big-fppf-site}.
Soit $T \in \Ob(\Sch_{fppf})$.
Soit $\{T_i \to T\}_{i \in I}$ un recouvrement fppf quelconque de $T$.
\begin{enumerate}
\item Il existe un recouvrement $\{U_j \to T\}_{j \in J}$ de $T$ dans le site
$\Sch_{fppf}$ qui raffine $\{T_i \to T\}_{i \in I}$.
\item Si $\{T_i \to T\}_{i \in I}$ est un recouvrement fppf standard, alors
il est tautologiquement équivalent à un recouvrement de $\Sch_{fppf}$.
\item Si $\{T_i \to T\}_{i \in I}$ est un recouvrement de Zariski, alors
il est tautologiquement équivalent à un recouvrement de $\Sch_{fppf}$.
\end{enumerate}
\end{lemma}

\begin{proof}
Pour chaque $i$, choisissons un recouvrement ouvert affine $T_i = \bigcup_{j \in J_i} T_{ij}$
tel que chaque $T_{ij}$ s'envoie dans un sous-schéma ouvert affine de $T$. D'après le
lemme \ref{lemma-fppf},
le raffinement $\{T_{ij} \to T\}_{i \in I, j \in J_i}$ est lui aussi un recouvrement fppf
de $T$. On peut donc supposer chaque $T_i$ affine et s'envoyant dans
un ouvert affine $W_i$ de $T$. En appliquant
Ensembles, lemme \ref{sets-lemma-what-is-in-it},
on voit que $W_i$ est isomorphe à un objet de $\Sch_{fppf}$.
Mais alors $T_i$, comme schéma de type fini au-dessus de $W_i$,
est isomorphe à un objet $V_i$ de $\Sch_{fppf}$ par une seconde
application de
Ensembles, lemme \ref{sets-lemma-what-is-in-it}.
Le recouvrement $\{V_i \to T\}_{i \in I}$ raffine $\{T_i \to T\}_{i \in I}$
(puisqu'ils sont isomorphes).
De plus, $\{V_i \to T\}_{i \in I}$ est combinatoirement équivalent à un
recouvrement $\{U_j \to T\}_{j \in J}$ de $T$ dans le site
$\Sch_{fppf}$ d'après
Ensembles, lemme \ref{sets-lemma-what-is-in-it}.
Le recouvrement $\{U_j \to T\}_{j \in J}$ est un raffinement comme en (1).
Dans la situation de (2), (3), chacun des
schémas $T_i$ est isomorphe à un objet de $\Sch_{fppf}$ d'après
Ensembles, lemme \ref{sets-lemma-what-is-in-it},
et une nouvelle application de
Ensembles, lemme \ref{sets-lemma-coverings-site},
donne le résultat voulu.
\end{proof}

\begin{definition}
\label{definition-big-small-fppf}
Soit $S$ un schéma. Soit $\Sch_{fppf}$ un gros site
fppf contenant $S$.
\begin{enumerate}
\item Le {\it gros site fppf de $S$}, noté
$(\Sch/S)_{fppf}$, est le site $\Sch_{fppf}/S$
introduit dans Sites, section \ref{sites-section-localize}.
\item Le {\it gros site fppf affine de $S$}, noté
$(\textit{Aff}/S)_{fppf}$, est la sous-catégorie pleine de
$(\Sch/S)_{fppf}$ dont les objets sont les $U/S$ affines.
Un recouvrement de $(\textit{Aff}/S)_{fppf}$ est tout recouvrement
$\{U_i \to U\}$ de $(\Sch/S)_{fppf}$ qui soit un
recouvrement fppf standard.
\end{enumerate}
\end{definition}

\noindent
Il n'est pas tout à fait évident que
le gros site fppf affine soit un site. Vérifions-le.

\begin{lemma}
\label{lemma-verify-site-fppf}
Soit $S$ un schéma. Soit $\Sch_{fppf}$ un gros site
fppf contenant $S$. Alors $(\textit{Aff}/S)_{fppf}$ est un site.
\end{lemma}

\begin{proof}
Montrons que $(\textit{Aff}/S)_{fppf}$ est un site.
En raisonnant comme dans la démonstration du lemme \ref{lemma-verify-site-etale},
il suffit de montrer que l'ensemble
des recouvrements fppf standards d'affines satisfait aux propriétés
(1), (2) et (3) de
Sites, définition \ref{sites-definition-site}.
C'est clair : par exemple, étant donné un recouvrement fppf
standard $\{T_i \to T\}_{i\in I}$ et, pour chaque
$i$, un recouvrement fppf standard $\{T_{ij} \to T_i\}_{j\in J_i}$,
$\{T_{ij} \to T\}_{i \in I, j\in J_i}$ est un recouvrement fppf standard
car $\bigcup_{i\in I} J_i$ est fini et chaque $T_{ij}$ est affine.
\end{proof}

\begin{lemma}
\label{lemma-fibre-products-fppf}
Soit $S$ un schéma. Soit $\Sch_{fppf}$ un gros site
fppf contenant $S$. Les catégories sous-jacentes des sites
$\Sch_{fppf}$, $(\Sch/S)_{fppf}$,
et $(\textit{Aff}/S)_{fppf}$ admettent des produits fibrés.
Dans chaque cas, le foncteur évident vers la catégorie $\Sch$ de
tous les schémas commute à la formation des produits fibrés. La catégorie
$(\Sch/S)_{fppf}$ possède un objet final, à savoir $S/S$.
\end{lemma}

\begin{proof}
Pour $\Sch_{fppf}$, cela est vrai par construction ; voir
Ensembles, lemme \ref{sets-lemma-what-is-in-it}.
Supposons donnés des morphismes $U \to S$, $V \to U$, $W \to U$
de schémas, avec $U, V, W \in \Ob(\Sch_{fppf})$.
Le produit fibré $V \times_U W$ dans $\Sch_{fppf}$
est un produit fibré dans $\Sch$ et
est le produit fibré de $V/S$ et $W/S$ au-dessus de $U/S$ dans
la catégorie de tous les schémas au-dessus de $S$, donc aussi un
produit fibré dans $(\Sch/S)_{fppf}$.
Cela prouve le résultat pour $(\Sch/S)_{fppf}$.
Si $U, V, W$ sont affines, $V \times_U W$ l'est aussi, d'où le
résultat pour $(\textit{Aff}/S)_{fppf}$.
\end{proof}

\noindent
Vérifions ensuite que le gros site affine définit le même
topos que le gros site.

\begin{lemma}
\label{lemma-affine-big-site-fppf}
Soit $S$ un schéma. Soit $\Sch_{fppf}$ un gros site
fppf contenant $S$.
Le foncteur $(\textit{Aff}/S)_{fppf} \to (\Sch/S)_{fppf}$
est cocontinu et induit une équivalence de topos de
$\Sh((\textit{Aff}/S)_{fppf})$ vers
$\Sh((\Sch/S)_{fppf})$.
\end{lemma}

\begin{proof}
La notion de foncteur cocontinu spécial est introduite dans
Sites, définition \ref{sites-definition-special-cocontinuous-functor}.
Il faut donc vérifier les hypothèses (1) -- (5) de
Sites, lemme \ref{sites-lemma-equivalence}.
Notons le foncteur d'inclusion
$u : (\textit{Aff}/S)_{fppf} \to (\Sch/S)_{fppf}$.
La cocontinuité signifie simplement que tout recouvrement fppf de
$T/S$, avec $T$ affine, admet un raffinement par un recouvrement fppf standard de $T$.
C'est le contenu du
lemme \ref{lemma-fppf-affine}.
Ainsi, (1) est vérifiée. Le foncteur $u$ est continu simplement parce qu'un recouvrement
fppf standard est un recouvrement fppf. Ainsi, (2) est vérifiée.
Les points (3) et (4) découlent immédiatement du fait que $u$ est
pleinement fidèle. Enfin, la condition (5) découle du
fait que tout schéma admet un recouvrement ouvert affine.
\end{proof}

\noindent
Établissons maintenant quelques relations entre les topos
associés à ces sites.

\begin{lemma}
\label{lemma-morphism-big-fppf}
Soit $\Sch_{fppf}$ un gros site fppf.
Soit $f : T \to S$ un morphisme dans $\Sch_{fppf}$.
Le foncteur
$$
u : (\Sch/T)_{fppf} \longrightarrow (\Sch/S)_{fppf},
\quad
V/T \longmapsto V/S
$$
est cocontinu et admet un adjoint à droite continu
$$
v : (\Sch/S)_{fppf} \longrightarrow (\Sch/T)_{fppf},
\quad
(U \to S) \longmapsto (U \times_S T \to T).
$$
Ils induisent le même morphisme de topos
$$
f_{big} :
\Sh((\Sch/T)_{fppf})
\longrightarrow
\Sh((\Sch/S)_{fppf})
$$
On a $f_{big}^{-1}(\mathcal{G})(U/T) = \mathcal{G}(U/S)$.
On a $f_{big, *}(\mathcal{F})(U/S) = \mathcal{F}(U \times_S T/T)$.
En outre, $f_{big}^{-1}$ admet un adjoint à gauche $f_{big!}$ qui commute aux
produits fibrés et aux égalisateurs.
\end{lemma}

\begin{proof}
Le foncteur $u$ est cocontinu, continu et commute aux produits fibrés
et aux égalisateurs. Par conséquent,
Sites, lemmes \ref{sites-lemma-when-shriek} et
\ref{sites-lemma-preserve-equalizers}
s'appliquent et donnent la formule
pour $f_{big}^{-1}$ ainsi que l'existence de $f_{big!}$. De plus,
le foncteur $v$ est un adjoint à droite car, étant donnés $U/T$ et $V/S$,
on a $\Mor_S(u(U), V) = \Mor_T(U, V \times_S T)$,
comme voulu. On peut donc appliquer
Sites, lemmes \ref{sites-lemma-have-functor-other-way} et
\ref{sites-lemma-have-functor-other-way-morphism} pour obtenir la
formule pour $f_{big, *}$.
\end{proof}

\begin{lemma}
\label{lemma-composition-fppf}
Étant donnés des schémas $X$, $Y$, $Z$ dans $(\Sch/S)_{fppf}$
et des morphismes $f : X \to Y$, $g : Y \to Z$, on a
$g_{big} \circ f_{big} = (g \circ f)_{big}$.
\end{lemma}

\begin{proof}
Cela résulte de la description simple de l'image directe
et de l'image inverse pour les foncteurs sur les gros sites donnée par le
lemme \ref{lemma-morphism-big-fppf}.
\end{proof}



















































\section{La topologie ph}
\label{section-ph}

\noindent
Dans cette section, nous définissons la topologie ph. C'est la topologie
engendrée par les recouvrements de Zariski et les morphismes propres surjectifs, voir
le lemme \ref{lemma-characterize-sheaf}.

\medskip\noindent
Nous reprenons nos notations et notre terminologie
de l'article \cite{ph} de Goodwillie et Lichtenbaum.
Ces auteurs montrent que, si l'on se restreint à la sous-catégorie des
schémas noethériens, la topologie ph coïncide avec la
« topologie h » telle que Voevodsky l'a définie à l'origine : c'est
la topologie engendrée par les recouvrements ouverts de Zariski et les
morphismes de type fini qui sont universellement submersifs.
Ils montrent aussi que les deux topologies ne coïncident pas sur les
schémas non noethériens, voir \cite[exemple 4.5]{ph}.
Nous revenons à (notre version de) la topologie h dans
Compléments sur la platitude, section \ref{flat-section-h}.

\medskip\noindent
Avant de pouvoir définir les recouvrements de notre topologie, il nous faut
quelques préliminaires.

\begin{definition}
\label{definition-standard-ph-covering}
Soit $T$ un schéma affine. Un {\it recouvrement ph standard} est une famille
$\{f_j : U_j \to T\}_{j = 1, \ldots, m}$ construite à partir d'un
morphisme propre surjectif $f : U \to T$ et d'un recouvrement ouvert affine
$U = \bigcup_{j = 1, \ldots, m} U_j$, en posant $f_j = f|_{U_j}$.
\end{definition}

\noindent
Il résulte immédiatement du lemme de Chow que l'on peut raffiner un
recouvrement ph standard par un recouvrement ph standard correspondant à
un morphisme projectif surjectif.

\begin{lemma}
\label{lemma-base-change-standard-ph}
Soit $\{f_j : U_j \to T\}_{j = 1, \ldots, m}$ un recouvrement ph standard.
Soit $T' \to T$ un morphisme de schémas affines.
Alors $\{U_j \times_T T' \to T'\}_{j = 1, \ldots, m}$ est un
recouvrement ph standard.
\end{lemma}

\begin{proof}
Soient $f : U \to T$ un morphisme propre surjectif et
$U = \bigcup_{j = 1, \ldots, m} U_j$ un recouvrement ouvert affine comme dans la
définition \ref{definition-standard-ph-covering}.
Alors $U \times_T T' \to T'$ est propre et surjectif
(Morphismes, lemmes \ref{morphisms-lemma-base-change-surjective} et
\ref{morphisms-lemma-base-change-proper}).
De plus, $U \times_T T' = \bigcup_{j = 1, \ldots, m} U_j \times_T T'$
est un recouvrement ouvert affine.
Ceci achève la démonstration.
\end{proof}

\begin{lemma}
\label{lemma-refine-by-standard-ph}
Soit $T$ un schéma affine. Chacun des types de familles suivants
de morphismes de but $T$ admet un raffinement qui est un recouvrement ph standard :
\begin{enumerate}
\item tout recouvrement ouvert de Zariski de $T$,
\item $\{W_{ji} \to T\}_{j = 1, \ldots, m, i = 1, \ldots n_j}$
où $\{W_{ji} \to U_j\}_{i = 1, \ldots, n_j}$
et $\{U_j \to T\}_{j = 1, \ldots, m}$ sont des recouvrements ph standards.
\end{enumerate}
\end{lemma}

\begin{proof}
L'assertion (1) résulte du fait que tout recouvrement ouvert de Zariski de $T$
peut être raffiné par un recouvrement ouvert affine fini.

\medskip\noindent
Preuve de (3). Choisissons $U \to T$ propre et surjectif et
$U = \bigcup_{j = 1, \ldots, m} U_j$ comme dans la
définition \ref{definition-standard-ph-covering}.
Choisissons $W_j \to U_j$ propre et surjectif et $W_j = \bigcup W_{ji}$
comme dans la définition \ref{definition-standard-ph-covering}.
Par le lemme de Chow (Limits, lemme \ref{limits-lemma-chow-finite-type}),
il existe $W'_j \to W_j$ propre et surjectif
et des immersions fermées $W'_j \to \mathbf{P}^{e_j}_{U_j}$.
Ainsi, après avoir remplacé $W_j$ par $W'_j$ et $W_j = \bigcup W_{ji}$
par un recouvrement ouvert affine convenable de $W'_j$, on peut supposer
qu'il existe une immersion fermée $W_j \subset \mathbf{P}^{e_j}_{U_j}$
pour tout $j = 1, \ldots, m$.

\medskip\noindent
Soit $\overline{W}_j \subset \mathbf{P}^{e_j}_U$ l'adhérence
schématique de $W_j$. Alors $W_j \subset \overline{W}_j$
est un sous-schéma ouvert ; en fait, $W_j$ est l'image inverse de
$U_j \subset U$ par le morphisme $\overline{W}_j \to U$.
(Pour le voir, utilisons le fait que $W_j \to \mathbf{P}^{e_j}_U$ est quasi-compact
et que, par suite, la formation de l'image schématique commute
à la restriction aux ouverts, voir
Morphismes, section \ref{morphisms-section-scheme-theoretic-image}.)
Munissons $Z_j = U \setminus U_j$
de la structure de sous-schéma fermé réduit induite.
Alors
$$
V_j = \overline{W}_j \amalg Z_j \to U
$$
est propre et surjectif, et le sous-schéma ouvert $W_j \subset V_j$
est l'image inverse de $U_j$. Ainsi, pour $v \in V_j$, $v \not \in W_j$,
on peut choisir un voisinage ouvert affine $v \in V_{j, v} \subset V_j$
dont l'image est contenue dans $U_{j'}$ pour un certain $1 \leq j' \leq m$.

\medskip\noindent
Pour achever la démonstration, considérons le morphisme propre surjectif
$$
V = V_1 \times_U V_2 \times_U \ldots \times_U V_m
\longrightarrow U \longrightarrow T
$$
et le recouvrement de $V$ par les ouverts affines
$$
V_{1, v_1} \times_U \ldots \times_U V_{j - 1, v_{j - 1}}
\times_U W_{j i} \times_U
V_{j + 1, v_{j + 1}} \times_U \ldots \times_U V_{m, v_m}
$$
Ceux-ci forment bien un recouvrement, car tout point de $U$ appartient
à un certain $U_j$ et l'image inverse de $U_j$ dans $V$ est égale à
$V_1 \times \ldots \times V_{j - 1} \times W_j \times V_{j + 1} \times
\ldots \times V_m$. Remarquons que le morphisme de
l'ouvert affine ci-dessus vers $T$ se factorise par $W_{ji}$ ;
on obtient donc un raffinement. Enfin, il suffit d'un nombre fini
de ces ouverts affines puisque $V$ est quasi-compact (comme schéma propre
sur le schéma affine $T$).
\end{proof}

\begin{definition}
\label{definition-ph-covering}
Soit $T$ un schéma. Un {\it recouvrement ph de $T$} est une famille
de morphismes de schémas $\{f_i : T_i \to T\}_{i \in I}$ telle
que $f_i$ soit localement de type fini et que, pour tout
ouvert affine $U \subset T$, il existe un recouvrement ph standard
$\{U_j \to U\}_{j = 1, \ldots, m}$ qui raffine la famille
$\{T_i \times_T U \to U\}_{i \in I}$.
\end{definition}

\noindent
Un recouvrement ph standard est un recouvrement ph d'après le
lemme \ref{lemma-base-change-standard-ph}.

\begin{lemma}
\label{lemma-zariski-ph}
Un recouvrement de Zariski est un recouvrement ph\footnote{Nous verrons
dans Compléments sur les morphismes, lemme \ref{more-morphisms-lemma-fppf-ph}, que les
recouvrements fppf (et donc les recouvrements syntomiques, lisses ou \'etales)
sont eux aussi des recouvrements ph.}.
\end{lemma}

\begin{proof}
Cela résulte du fait qu'un recouvrement de Zariski d'un schéma affine
peut être raffiné par un recouvrement ph standard d'après le
lemme \ref{lemma-refine-by-standard-ph}.
\end{proof}

\begin{lemma}
\label{lemma-surjective-proper-ph}
Soit $f : Y \to X$ un morphisme propre surjectif de schémas.
Alors $\{Y \to X\}$ est un recouvrement ph.
\end{lemma}

\begin{proof}
Omise.
\end{proof}

\begin{lemma}
\label{lemma-refine-by-ph}
Soit $T$ un schéma. Soit $\{f_i : T_i \to T\}_{i \in I}$ une famille
de morphismes telle que $f_i$ soit localement de type fini pour tout $i$.
Les conditions suivantes sont équivalentes :
\begin{enumerate}
\item $\{T_i \to T\}_{i \in I}$ est un recouvrement ph,
\item il existe un recouvrement ph qui raffine $\{T_i \to T\}_{i \in I}$, et
\item $\{\coprod_{i \in I} T_i \to T\}$ est un recouvrement ph.
\end{enumerate}
\end{lemma}

\begin{proof}
L'équivalence de (1) et (2) résulte immédiatement de la
définition \ref{definition-ph-covering}
et du fait qu'un raffinement d'un raffinement est un raffinement.
En vertu de l'équivalence de (1) et (2), et puisque
$\{T_i \to T\}_{i \in I}$ raffine $\{\coprod_{i \in I} T_i \to T\}$,
on voit que (1) implique (3). Supposons enfin que (3) soit vérifiée.
Soit $U \subset T$ un ouvert affine et soit
$\{U_j \to U\}_{j = 1, \ldots, m}$ un recouvrement ph standard
qui raffine $\{U \times_T \coprod_{i \in I} T_i \to U\}$.
Cela signifie que, pour tout $j$, on dispose d'un morphisme
$$
h_j :
U_j
\longrightarrow
U \times_T \coprod\nolimits_{i \in I} T_i =
\coprod\nolimits_{i \in I} U \times_T T_i
$$
au-dessus de $U$. Comme $U_j$ est quasi-compact, on obtient
des décompositions en réunion disjointe $U_j = \coprod_{i \in I} U_{j, i}$
en sous-schémas ouverts et fermés, presque tous vides,
telles que $h_j|_{U_{j, i}}$ envoie $U_{j, i}$ dans $U \times_T T_i$.
Il s'ensuit que
$$
\{U_{j, i} \to U\}_{j = 1, \ldots, m,\ i \in I,\ U_{j, i} \not = \emptyset}
$$
est un recouvrement ph standard (un petit détail est omis) qui raffine
$\{U \times_T T_i \to U\}_{i \in I}$. Ainsi, (1) est vérifiée.
\end{proof}

\noindent
Montrons ensuite que cette notion satisfait les conditions de
Sites, définition \ref{sites-definition-site}.

\begin{lemma}
\label{lemma-ph}
Soit $T$ un schéma.
\begin{enumerate}
\item Si $T' \to T$ est un isomorphisme, alors $\{T' \to T\}$
est un recouvrement ph de $T$.
\item Si $\{T_i \to T\}_{i\in I}$ est un recouvrement ph et si, pour tout
$i$, on a un recouvrement ph $\{T_{ij} \to T_i\}_{j\in J_i}$, alors
$\{T_{ij} \to T\}_{i \in I, j\in J_i}$ est un recouvrement ph.
\item Si $\{T_i \to T\}_{i\in I}$ est un recouvrement ph
et si $T' \to T$ est un morphisme de schémas, alors
$\{T' \times_T T_i \to T'\}_{i\in I}$ est un recouvrement ph.
\end{enumerate}
\end{lemma}

\begin{proof}
L'assertion (1) est claire.

\medskip\noindent
Preuve de (3). Le changement de base $T_i \times_T T' \to T'$ est
localement de type fini d'après
Morphismes, lemme \ref{morphisms-lemma-base-change-finite-type}.
Il suffit donc de vérifier la condition sur les ouverts affines.
Soit $U' \subset T'$ un sous-schéma ouvert affine.
Comme $U'$ est quasi-compact, on peut trouver un recouvrement ouvert
affine fini $U' = U'_1 \cup \ldots \cup U'$ tel que
l'image de $U'_j \to T$ soit contenue dans un ouvert affine $U_j \subset T$.
Choisissons un recouvrement ph standard $\{U_{jl} \to U_j\}_{l = 1, \ldots, n_j}$
qui raffine $\{T_i \times_T U_j \to U_j\}$.
D'après le lemme \ref{lemma-base-change-standard-ph},
le changement de base $\{U_{jl} \times_{U_j} U'_j \to U'_j\}$
est un recouvrement ph standard. Remarquons que $\{U'_j \to U'\}$
est lui aussi un recouvrement ph standard.
D'après le lemme \ref{lemma-refine-by-standard-ph},
la famille $\{U_{jl} \times_{U_j} U'_j \to U'\}$
peut être raffinée par un recouvrement ph standard. Comme
$\{U_{jl} \times_{U_j} U'_j \to U'\}$ raffine $\{T_i \times_T U' \to U'\}$,
on conclut.

\medskip\noindent
Preuve de (2). La composition préserve la propriété d'être localement de type fini,
voir Morphismes, lemme \ref{morphisms-lemma-composition-finite-type}.
Il suffit donc de vérifier la condition sur les ouverts affines.
Soit $U \subset T$ un ouvert affine. Choisissons d'abord un recouvrement ph standard
$\{U_k \to U\}_{k = 1, \ldots, m}$ qui raffine
$\{T_i \times_T U \to U\}$.
Supposons que le raffinement soit donné par des morphismes $U_k \to T_{i_k}$ sur $T$.
Alors
$$
\{T_{i_kj} \times_{T_{i_k}} U_k \to U_k\}_{j \in J_{i_k}}
$$
est un recouvrement ph d'après l'assertion (3). Comme $U_k$ est affine,
on peut trouver un recouvrement ph standard
$\{U_{ka} \to U_k\}_{a = 1, \ldots, b_k}$ qui raffine cette famille.
On applique alors le lemme \ref{lemma-refine-by-standard-ph}
pour voir que $\{U_{ka} \to U\}$ peut être raffiné par un
recouvrement ph standard. Comme $\{U_{ka} \to U\}$
raffine $\{T_{ij} \times_T U \to U\}$, ceci achève la démonstration.
\end{proof}

\begin{definition}
\label{definition-big-ph-site}
Un {\it gros site ph} est un site $\Sch_{ph}$ comme dans
Sites, définition \ref{sites-definition-site}, construit comme suit :
\begin{enumerate}
\item Choisir un ensemble quelconque de schémas $S_0$ et un ensemble quelconque
$\text{Cov}_0$ de recouvrements ph entre ces schémas.
\item Prendre pour catégorie sous-jacente une catégorie quelconque $\Sch_\alpha$
construite comme dans Ensembles, lemme \ref{sets-lemma-construct-category},
à partir de l'ensemble $S_0$.
\item Choisir un ensemble quelconque de recouvrements comme dans
Ensembles, lemme \ref{sets-lemma-coverings-site}, à partir de la
catégorie $\Sch_\alpha$, de la classe des recouvrements ph
et de l'ensemble $\text{Cov}_0$ choisi ci-dessus.
\end{enumerate}
\end{definition}

\noindent
Voir les remarques qui suivent la définition \ref{definition-big-zariski-site}
pour la motivation et les explications concernant la définition des gros sites.

\medskip\noindent
Avant de poursuivre l'introduction du gros site ph d'un
schéma $S$, signalons que la topologie d'un gros site ph
$\Sch_{ph}$ est, en un certain sens, induite par la topologie ph
sur la catégorie de tous les schémas.

\begin{lemma}
\label{lemma-ph-induced}
Soit $\Sch_{ph}$ un gros site ph comme dans la
définition \ref{definition-big-ph-site}.
Soit $T \in \Ob(\Sch_{ph})$.
Soit $\{T_i \to T\}_{i \in I}$ un recouvrement ph quelconque de $T$.
\begin{enumerate}
\item Il existe un recouvrement $\{U_j \to T\}_{j \in J}$ de $T$ dans le site
$\Sch_{ph}$ qui raffine $\{T_i \to T\}_{i \in I}$.
\item Si $\{T_i \to T\}_{i \in I}$ est un recouvrement ph standard, alors
il est tautologiquement équivalent à un recouvrement de $\Sch_{ph}$.
\item Si $\{T_i \to T\}_{i \in I}$ est un recouvrement de Zariski, alors
il est tautologiquement équivalent à un recouvrement de $\Sch_{ph}$.
\end{enumerate}
\end{lemma}

\begin{proof}
Pour tout $i$, choisissons un recouvrement ouvert affine $T_i = \bigcup_{j \in J_i} T_{ij}$
tel que chaque $T_{ij}$ s'envoie dans un sous-schéma ouvert affine de $T$. D'après les
lemmes \ref{lemma-zariski-ph} et \ref{lemma-ph},
le raffinement $\{T_{ij} \to T\}_{i \in I, j \in J_i}$ est lui aussi un recouvrement ph
de $T$. On peut donc supposer que chaque $T_i$ est affine et s'envoie dans
un ouvert affine $W_i$ de $T$. En appliquant
Ensembles, lemme \ref{sets-lemma-what-is-in-it}
on voit que $W_i$ est isomorphe à un objet de $\Sch_{ph}$.
Mais alors $T_i$, comme schéma de type fini sur $W_i$,
est isomorphe à un objet $V_i$ de $\Sch_{ph}$ par une seconde
application de
Ensembles, lemme \ref{sets-lemma-what-is-in-it}.
Le recouvrement $\{V_i \to T\}_{i \in I}$ raffine $\{T_i \to T\}_{i \in I}$
(car ils sont isomorphes).
De plus, $\{V_i \to T\}_{i \in I}$ est combinatoirement équivalent à un
recouvrement $\{U_j \to T\}_{j \in J}$ de $T$ dans le site
$\Sch_{ph}$ d'après
Ensembles, lemme \ref{sets-lemma-what-is-in-it}.
Le recouvrement $\{U_j \to T\}_{j \in J}$ est un raffinement comme en (1).
Dans les situations de (2) et (3), chacun des
schémas $T_i$ est isomorphe à un objet de $\Sch_{ph}$ d'après
Ensembles, lemme \ref{sets-lemma-what-is-in-it},
et une autre application de
Ensembles, lemme \ref{sets-lemma-coverings-site}
donne le résultat voulu.
\end{proof}

\begin{definition}
\label{definition-big-small-ph}
Soit $S$ un schéma. Soit $\Sch_{ph}$ un gros site ph contenant $S$.
\begin{enumerate}
\item Le {\it gros site ph de $S$}, noté
$(\Sch/S)_{ph}$, est le site $\Sch_{ph}/S$
introduit dans Sites, section \ref{sites-section-localize}.
\item Le {\it gros site ph affine de $S$}, noté
$(\textit{Aff}/S)_{ph}$, est la sous-catégorie pleine de
$(\Sch/S)_{ph}$ dont les objets sont les $U/S$ affines.
Un recouvrement de $(\textit{Aff}/S)_{ph}$ est tout recouvrement fini
$\{U_i \to U\}$ de $(\Sch/S)_{ph}$ tel que $U_i$ et $U$ soient affines.
\end{enumerate}
\end{definition}

\noindent
Remarquons que les recouvrements de $(\textit{Aff}/S)_{ph}$ ne sont {\it pas}
donnés par les recouvrements ph standards. La raison en est simplement que
le deuxième axiome de Sites, définition \ref{sites-definition-site}, échouerait.
Les recouvrements de $(\textit{Aff}/S)_{ph}$ sont plutôt les familles finies
$\{U_i \to U\}$ de morphismes de type fini entre objets affines
de $(\Sch/S)_{ph}$ qui peuvent être raffinées par un recouvrement ph standard.
Énonçons et démontrons explicitement que le gros site ph affine est un site.

\begin{lemma}
\label{lemma-verify-site-ph}
Soit $S$ un schéma. Soit $\Sch_{ph}$ un gros site ph
contenant $S$. Alors $(\textit{Aff}/S)_{ph}$ est un site.
\end{lemma}

\begin{proof}
En raisonnant comme dans la démonstration du lemme \ref{lemma-verify-site-etale},
il suffit de montrer que la famille des recouvrements ph finis
$\{U_i \to U\}$, où $U$ et les $U_i$ sont affines, satisfait les propriétés
(1), (2) et (3) de
Sites, définition \ref{sites-definition-site}.
C'est clair car, par exemple, étant donné un recouvrement ph fini
$\{T_i \to T\}_{i\in I}$ où $T_i, T$ sont affines, et, pour tout
$i$, un recouvrement ph fini $\{T_{ij} \to T_i\}_{j\in J_i}$ où $T_{ij}$ est
affine, $\{T_{ij} \to T\}_{i \in I, j\in J_i}$ est alors un recouvrement ph
(lemme \ref{lemma-ph}), $\bigcup_{i\in I} J_i$ est fini et
chaque $T_{ij}$ est affine.
\end{proof}

\begin{lemma}
\label{lemma-fibre-products-ph}
Soit $S$ un schéma. Soit $\Sch_{ph}$ un gros site ph
contenant $S$. Les catégories sous-jacentes aux sites
$\Sch_{ph}$, $(\Sch/S)_{ph}$ et $(\textit{Aff}/S)_{ph}$ ont des produits fibrés.
Dans chaque cas, le foncteur évident vers la catégorie $\Sch$ de
tous les schémas commute aux produits fibrés. La catégorie
$(\Sch/S)_{ph}$ a un objet final, à savoir $S/S$.
\end{lemma}

\begin{proof}
Pour $\Sch_{ph}$, c'est vrai par construction, voir
Ensembles, lemme \ref{sets-lemma-what-is-in-it}.
Supposons donnés des morphismes de schémas $U \to S$, $V \to U$, $W \to U$
avec $U, V, W \in \Ob(\Sch_{ph})$.
Le produit fibré $V \times_U W$ dans $\Sch_{ph}$
est un produit fibré dans $\Sch$ et
est le produit fibré de $V/S$ et $W/S$ au-dessus de $U/S$ dans
la catégorie de tous les schémas au-dessus de $S$, et donc aussi un
produit fibré dans $(\Sch/S)_{ph}$.
Ceci démontre le résultat pour $(\Sch/S)_{ph}$.
Si $U, V, W$ sont affines, $V \times_U W$ l'est aussi, d'où le
résultat pour $(\textit{Aff}/S)_{ph}$.
\end{proof}

\noindent
Vérifions ensuite que le gros site affine définit le même
topos que le gros site.

\begin{lemma}
\label{lemma-affine-big-site-ph}
Soit $S$ un schéma. Soit $\Sch_{ph}$ un gros site ph
contenant $S$.
Le foncteur $(\textit{Aff}/S)_{ph} \to (\Sch/S)_{ph}$
est cocontinu et induit une équivalence de topos de
$\Sh((\textit{Aff}/S)_{ph})$ vers
$\Sh((\Sch/S)_{ph})$.
\end{lemma}

\begin{proof}
La notion de foncteur cocontinu spécial est introduite dans
Sites, définition \ref{sites-definition-special-cocontinuous-functor}.
Il faut donc vérifier les hypothèses (1) -- (5) de
Sites, lemme \ref{sites-lemma-equivalence}.
Notons le foncteur d'inclusion
$u : (\textit{Aff}/S)_{ph} \to (\Sch/S)_{ph}$.
La cocontinuité résulte du fait que tout recouvrement ph de
$T/S$, avec $T$ affine, peut être raffiné par un recouvrement ph standard de $T$
par définition. Ainsi, (1) est vérifiée. Le foncteur $u$ est continu simplement
parce qu'un recouvrement ph fini d'un affine par des affines est un recouvrement ph.
Ainsi, (2) est vérifiée.
Les assertions (3) et (4) résultent immédiatement du fait que $u$ est
pleinement fidèle. Enfin, la condition (5) résulte du
fait que tout schéma possède un recouvrement ouvert affine (qui est
un recouvrement ph).
\end{proof}

\begin{lemma}
\label{lemma-characterize-sheaf}
Soit $\mathcal{F}$ un préfaisceau sur $(\Sch/S)_{ph}$.
Alors $\mathcal{F}$ est un faisceau si et seulement si
\begin{enumerate}
\item $\mathcal{F}$ satisfait la condition de faisceau pour les
recouvrements de Zariski, et
\item si $f : V \to U$ est propre et surjectif, alors
$\mathcal{F}(U)$ s'envoie bijectivement sur l'égalisateur
des deux applications $\mathcal{F}(V) \to \mathcal{F}(V \times_U V)$.
\end{enumerate}
De plus, en présence de (1), la propriété (2) équivaut à la propriété
\begin{enumerate}
\item[(2')] de faisceau pour $\{V \to U\}$ comme en (2), avec $U$ affine.
\end{enumerate}
\end{lemma}

\begin{proof}
Montrons que, si (1) et (2) sont vérifiées, alors $\mathcal{F}$ est un faisceau.
Soit $\{T_i \to T\}$ un recouvrement ph, c'est-à-dire un recouvrement dans $(\Sch/S)_{ph}$.
Vérifions la condition de faisceau pour ce recouvrement.
Soient $s_i \in \mathcal{F}(T_i)$ des sections dont les restrictions donnent la même
section sur $T_i \times_T T_{i'}$. Montrons qu'il existe une
unique section $s \in \mathcal{F}(T)$ dont la restriction est $s_i$ sur $T_i$.
Soit $T = \bigcup U_j$ un recouvrement ouvert affine.
D'après la propriété (1), il suffit de produire des sections $s_j \in \mathcal{F}(U_j)$
qui coïncident sur $U_j \cap U_{j'}$ pour construire $s$.
Considérons les recouvrements ph $\{T_i \times_T U_j \to U_j\}$.
Alors les $s_{ji} = s_i|_{T_i \times_T U_j}$ sont des sections qui coïncident
sur $(T_i \times_T U_j) \times_{U_j} (T_{i'} \times_T U_j)$.
Choisissons un morphisme propre surjectif $V_j \to U_j$ et un recouvrement ouvert
affine fini $V_j = \bigcup V_{jk}$ tel que le recouvrement ph standard
$\{V_{jk} \to U_j\}$ raffine $\{T_i \times_T U_j \to U_j\}$.
Si $s_{jk} \in \mathcal{F}(V_{jk})$
désigne l'image inverse de $s_{ji}$ sur $V_{jk}$ par les
morphismes sous-entendus, les $s_{jk}$ se recollent en une section
$s'_j \in \mathcal{F}(V_j)$. En utilisant encore une fois l'accord sur les intersections,
on voit que $s'_j$ appartient à l'égalisateur des deux
applications $\mathcal{F}(V_j) \to \mathcal{F}(V_j \times_{U_j} V_j)$.
Ainsi, d'après (2), $s'_j$ provient d'une unique section
$s_j \in \mathcal{F}(U_j)$. Nous omettons de vérifier que ces
sections $s_j$ possèdent toutes les propriétés voulues.

\medskip\noindent
Preuve de l'équivalence de (2) et (2') en présence de (1).
Supposons que $V \to U$ soit un morphisme de $(\Sch/S)_{ph}$ qui soit
propre et surjectif. Choisissons un
recouvrement ouvert affine $U = \bigcup U_i$ et posons $V_i = V \times_U U_i$.
On voit alors que $\mathcal{F}(U) \to \mathcal{F}(V)$
est injective, car $\mathcal{F}(U_i) \to \mathcal{F}(V_i)$
est injective d'après (2') et $\mathcal{F}(U) \to \prod \mathcal{F}(U_i)$
est injective d'après (1). Supposons enfin donné un élément
$t \in \mathcal{F}(V)$ de l'égalisateur des deux applications
$\mathcal{F}(V) \to \mathcal{F}(V \times_U V)$.
Alors $t|_{V_i}$ appartient à l'égalisateur des deux applications
$\mathcal{F}(V_i) \to \mathcal{F}(V_i \times_{U_i} V_i)$
pour tout $i$. On obtient donc une unique section $s_i \in \mathcal{F}(U_i)$
qui s'envoie sur $t|_{V_i}$ pour tout $i$ d'après (2').
Nous omettons de vérifier que $s_i|_{U_i \cap U_j} = s_j|_{U_i \cap U_j}$
pour tous $i, j$ ; cela utilise la propriété d'unicité que l'on vient d'établir.
La propriété de faisceau pour le recouvrement $U = \bigcup U_i$ fournit
une section $s \in \mathcal{F}(U)$. Nous omettons de démontrer que $s$
s'envoie sur $t$ dans $\mathcal{F}(V)$.
\end{proof}

\noindent
Établissons ensuite quelques relations entre les topos
associés à ces sites.

\begin{lemma}
\label{lemma-morphism-big-ph}
Soit $\Sch_{ph}$ un gros site ph.
Soit $f : T \to S$ un morphisme de $\Sch_{ph}$.
Le foncteur
$$
u : (\Sch/T)_{ph} \longrightarrow (\Sch/S)_{ph},
\quad
V/T \longmapsto V/S
$$
est cocontinu et admet un adjoint à droite continu
$$
v : (\Sch/S)_{ph} \longrightarrow (\Sch/T)_{ph},
\quad
(U \to S) \longmapsto (U \times_S T \to T).
$$
Ils induisent le même morphisme de topos
$$
f_{big} :
\Sh((\Sch/T)_{ph})
\longrightarrow
\Sh((\Sch/S)_{ph})
$$
On a $f_{big}^{-1}(\mathcal{G})(U/T) = \mathcal{G}(U/S)$.
On a $f_{big, *}(\mathcal{F})(U/S) = \mathcal{F}(U \times_S T/T)$.
De plus, $f_{big}^{-1}$ admet un adjoint à gauche $f_{big!}$ qui commute aux
produits fibrés et aux égalisateurs.
\end{lemma}

\begin{proof}
Le foncteur $u$ est cocontinu, continu et commute aux produits fibrés
et aux égalisateurs. Par conséquent,
Sites, lemmes \ref{sites-lemma-when-shriek} et
\ref{sites-lemma-preserve-equalizers}
s'appliquent et donnent la formule
pour $f_{big}^{-1}$ ainsi que l'existence de $f_{big!}$. De plus,
le foncteur $v$ est adjoint à droite car, étant donnés $U/T$ et $V/S$,
on a $\Mor_S(u(U), V) = \Mor_T(U, V \times_S T)$,
comme voulu. On peut donc appliquer
Sites, lemmes \ref{sites-lemma-have-functor-other-way} et
\ref{sites-lemma-have-functor-other-way-morphism} pour obtenir la
formule pour $f_{big, *}$.
\end{proof}

\begin{lemma}
\label{lemma-composition-ph}
Étant donnés des schémas $X$, $Y$, $Y$ dans $(\Sch/S)_{ph}$
et des morphismes $f : X \to Y$, $g : Y \to Z$, on a
$g_{big} \circ f_{big} = (g \circ f)_{big}$.
\end{lemma}

\begin{proof}
Cela résulte de la description simple de l'image directe
et de l'image inverse des foncteurs sur les gros sites donnée dans le
lemme \ref{lemma-morphism-big-ph}.
\end{proof}



































\section{La topologie fpqc}
\label{section-fpqc}

\begin{definition}
\label{definition-fpqc-covering}
Soit $T$ un schéma. Un {\it recouvrement fpqc de $T$} est une famille
de morphismes de schémas $\{f_i : T_i \to T\}_{i \in I}$
telle que chaque $f_i$ soit plat et que, pour tout ouvert affine
$U \subset T$, il existe $n \geq 0$, une application
$a : \{1, \ldots, n\} \to I$ et des ouverts affines
$V_j \subset T_{a(j)}$, $j = 1, \ldots, n$,
tels que $\bigcup_{j = 1}^n f_{a(j)}(V_j) = U$.
\end{definition}

\noindent
Cette condition implique bien que $T = \bigcup f_i(T_i)$.
Il est un peu plus difficile de reconnaître un recouvrement fpqc ; nous donnons donc
quelques lemmes à cet effet.

\begin{lemma}
\label{lemma-recognize-fpqc-covering}
Soit $T$ un schéma. Soit $\{f_i : T_i \to T\}_{i \in I}$ une famille de
morphismes de schémas de but $T$. Les conditions suivantes sont équivalentes :
\begin{enumerate}
\item $\{f_i : T_i \to T\}_{i \in I}$ est un recouvrement fpqc,
\item chaque $f_i$ est plat et, pour tout ouvert affine $U \subset T$,
il existe des ouverts quasi-compacts
$U_i \subset T_i$, presque tous vides,
tels que $U = \bigcup f_i(U_i)$,
\item chaque $f_i$ est plat et il existe un recouvrement ouvert affine
$T = \bigcup_{\alpha \in A} U_\alpha$ et, pour tout $\alpha \in A$,
il existe $i_{\alpha, 1}, \ldots, i_{\alpha, n(\alpha)} \in I$
et des ouverts quasi-compacts $U_{\alpha, j} \subset T_{i_{\alpha, j}}$ tels que
$U_\alpha =
\bigcup_{j = 1, \ldots, n(\alpha)} f_{i_{\alpha, j}}(U_{\alpha, j})$.
\end{enumerate}
Si $T$ est quasi-séparé, ces conditions équivalent aussi à
\begin{enumerate}
\item[(4)] chaque $f_i$ est plat et, pour tout $t \in T$, il existe
$i_1, \ldots, i_n \in I$ et des ouverts quasi-compacts $U_j \subset T_{i_j}$
tels que $\bigcup_{j = 1, \ldots, n} f_{i_j}(U_j)$ soit un
voisinage (non nécessairement ouvert) de $t$ dans $T$.
\end{enumerate}
\end{lemma}

\begin{proof}
Nous omettons la démonstration de l'équivalence de (1), (2) et (3).
Supposons désormais que $T$ soit quasi-séparé.
Montrons que (4) implique (2). Soit $U \subset T$ un ouvert affine.
Pour démontrer (2), il suffit de montrer que, pour tout $t \in U$, il existe
un nombre fini d'ouverts quasi-compacts $U_j \subset T_{i_j}$ tels que
$f_{i_j}(U_j) \subset U$ et que $\bigcup f_{i_j}(U_j)$
soit un voisinage de $t$ dans $U$. Par hypothèse, il existe bien
un nombre fini d'ouverts quasi-compacts $U'_j \subset T_{i_j}$ tels que
$\bigcup f_{i_j}(U'_j)$ soit un voisinage de $t$ dans $T$.
Puisque $T$ est quasi-séparé, on voit que
$U_j = U'_j \cap f_{i_j}^{-1}(U) = (f_{i_j}|_{U'_j})^{-1}(U)$
est quasi-compact d'après
Schémas, lemme \ref{schemes-lemma-quasi-compact-permanence}.
Ainsi, (2) est vérifiée. Comme il est clair que (2) implique
(4), la démonstration est achevée.
\end{proof}

\begin{example}
\label{example-not-fpqc}
Soit $X = X_1 \cup X_2$, où $X_1 = X_2 = \Spec(k[t_1,t_2,\dots])$
est le schéma non quasi-séparé de
Schémas, exemple \ref{schemes-example-not-quasi-separated}, et soit
$Y = X_1 \amalg \Spec(\mathcal{O}_{X_2, 0})$.
Alors le morphisme évident $f : Y \to X$
est plat et surjectif et, puisque $Y$ est quasi-compact,
satisfait trivialement à (4) du lemme \ref{lemma-recognize-fpqc-covering}.
Mais $\{f : Y \to X\}$ n'est pas un recouvrement fpqc. En effet, affirmons qu'il
n'existe aucun ouvert quasi-compact $W$ de $Y$ dont l'image par $f$ soit $X_2$.
En effet, $f^{-1}(X_2) = X_1 \setminus \{0\} \amalg \Spec(\mathcal{O}_{X_2, 0})$.
Donc, si $W$ existe, alors $W = U \amalg \Spec(\mathcal{O}_{X_2, 0})$
pour un ouvert quasi-compact $U$ de $X_1 \setminus \{0\} = X_2 \setminus \{0\}$.
Mais on peut alors écrire $U = D(f_1) \cup \ldots \cup D(f_n)$ pour
certains $f_i \in k[t_1, t_2, \ldots]$ qui s'annulent en $0$.
Supposons que $f_1, \ldots, f_n$ n'utilisent que les variables
$x_1, \ldots, x_m$. Alors $p = (0, \ldots, 0, 1, 0, \ldots)$, où $1$
occupe la $(m + 1)$-ième place, est un point fermé de $X_2$ qui n'appartient pas à l'image
de $W$.
\end{example}

\begin{lemma}
\label{lemma-disjoint-union-is-fpqc-covering}
Soit $T$ un schéma. Soit $\{f_i : T_i \to T\}_{i \in I}$ une famille de
morphismes de schémas de but $T$. Les conditions suivantes sont équivalentes :
\begin{enumerate}
\item $\{f_i : T_i \to T\}_{i \in I}$ est un recouvrement fpqc, et
\item en posant $T' = \coprod_{i \in I} T_i$ et $f = \coprod_{i \in I} f_i$,
la famille $\{f : T' \to T\}$ est un recouvrement fpqc.
\end{enumerate}
\end{lemma}

\begin{proof}
Supposons que $U \subset T$ soit un ouvert affine. Si (1) est vérifiée, il existe
$i_1, \ldots, i_n \in I$ et des ouverts affines $U_j \subset T_{i_j}$ tels que
$U = \bigcup_{j = 1, \ldots, n} f_{i_j}(U_j)$. Alors
$U_1 \amalg \ldots \amalg U_n \subset T'$ est un ouvert quasi-compact dont l'image
est $U$. Ainsi, $\{f : T' \to T\}$ est un recouvrement fpqc d'après le
lemme \ref{lemma-recognize-fpqc-covering}.
Réciproquement, si (2) est vérifiée, il existe un ouvert quasi-compact
$U' \subset T'$ tel que $U = f(U')$. Alors $U_j = U' \cap T_j$ est un ouvert quasi-compact
de $T_j$ et est vide pour presque tout $j$. D'après le
lemme \ref{lemma-recognize-fpqc-covering}, on voit que (1) est vérifiée.
\end{proof}

\begin{lemma}
\label{lemma-family-flat-dominated-covering}
Soit $T$ un schéma. Soit $\{f_i : T_i \to T\}_{i \in I}$ une famille de
morphismes de schémas de but $T$. Supposons que
\begin{enumerate}
\item chaque $f_i$ soit plat, et
\item la famille $\{f_i : T_i \to T\}_{i \in I}$ puisse être raffinée par un
recouvrement fpqc de $T$.
\end{enumerate}
Alors $\{f_i : T_i \to T\}_{i \in I}$ est un recouvrement fpqc de $T$.
\end{lemma}

\begin{proof}
Soit $\{g_j : X_j \to T\}_{j \in J}$ un recouvrement fpqc qui raffine
$\{f_i : T_i \to T\}$. Supposons que $U \subset T$ soit un ouvert affine.
Choisissons $j_1, \ldots, j_m \in J$ et des ouverts affines $V_k \subset X_{j_k}$
tel que $U = \bigcup g_{j_k}(V_k)$. Pour tout $j$, choisissons $i_j \in I$
et un morphisme $h_j : X_j \to T_{i_j}$ tel que $g_j = f_{i_j} \circ h_j$.
Comme $h_{j_k}(V_k)$ est quasi-compact, on peut trouver un ouvert quasi-compact
$h_{j_k}(V_k) \subset U_k \subset f_{i_{j_k}}^{-1}(U)$.
Alors $U = \bigcup f_{i_{j_k}}(U_k)$. On en conclut que
$\{f_i : T_i \to T\}_{i \in I}$ est un recouvrement fpqc d'après le
lemme \ref{lemma-recognize-fpqc-covering}.
\end{proof}

\begin{lemma}
\label{lemma-family-flat-fpqc-local-covering}
Soit $T$ un schéma. Soit $\{f_i : T_i \to T\}_{i \in I}$ une famille de
morphismes de schémas de but $T$. Supposons que
\begin{enumerate}
\item chaque $f_i$ soit plat, et
\item il existe un recouvrement fpqc
$\{g_j : S_j \to T\}_{j \in J}$ tel que chacun des
$\{S_j \times_T T_i \to S_j\}_{i \in I}$ soit un recouvrement fpqc.
\end{enumerate}
Alors $\{f_i : T_i \to T\}_{i \in I}$ est un recouvrement fpqc de $T$.
\end{lemma}

\begin{proof}
Nous utiliserons le lemme \ref{lemma-recognize-fpqc-covering} sans autre
mention. Soit $U \subset T$ un ouvert affine. D'après (2), il existe
des ouverts quasi-compacts $V_j \subset S_j$ pour $j \in J$, presque tous vides, tels que
$U = \bigcup g_j(V_j)$. Pour tout $j$, on peut alors choisir des ouverts quasi-compacts
$W_{ij} \subset S_j \times_T T_i$ pour $i \in I$, presque tous vides,
tels que $V_j = \bigcup_i \text{pr}_1(W_{ij})$. Ainsi,
$\{S_j \times_T T_i \to T\}$ est un recouvrement fpqc.
Comme ce recouvrement raffine $\{f_i : T_i \to T\}$, on conclut par le
lemme \ref{lemma-family-flat-dominated-covering}.
\end{proof}

\begin{lemma}
\label{lemma-zariski-etale-smooth-syntomic-fppf-fpqc}
Tout recouvrement fppf est un recouvrement fpqc et, a fortiori,
tout recouvrement syntomique, lisse, \'etale ou de Zariski est un recouvrement fpqc.
\end{lemma}

\begin{proof}
Montrons qu'un recouvrement fppf est un recouvrement fpqc ; le
reste résultera alors du
lemme \ref{lemma-zariski-etale-smooth-syntomic-fppf}.
Soit $\{f_i : U_i \to U\}_{i \in I}$ un recouvrement fppf.
Par définition, les $f_i$ sont plats, ce qui vérifie la première
condition de la définition \ref{definition-fpqc-covering}. Pour vérifier la
seconde, soit $V \subset U$ un sous-ensemble ouvert affine.
Écrivons $f_i^{-1}(V) = \bigcup_{j \in J_i} V_{ij}$
pour des ouverts affines $V_{ij} \subset U_i$. Comme chaque $f_i$ est ouvert
(Morphismes, lemme \ref{morphisms-lemma-fppf-open}), on voit que
$V = \bigcup_{i\in I} \bigcup_{j \in J_i} f_i(V_{ij})$
est un recouvrement ouvert de $V$.
Comme $V$ est quasi-compact, ce recouvrement admet un
raffinement fini. Ceci achève la démonstration.
\end{proof}

\noindent
La topologie fpqc\footnote{Les lettres fpqc signifient
« fid\`element plat quasi-compacte ».}
ne peut pas être traitée de la même manière que la topologie
fppf\footnote{Plus précisément, l'analogue du
lemme \ref{lemma-fppf-induced} pour la topologie fpqc n'est pas valable.}.
En effet, supposons que $R$ soit un anneau non nul. Nous verrons au
lemme \ref{lemma-no-set-of-fpqc-covers-is-initial}
qu'il n'existe pas d'ensemble $A$ de recouvrements fpqc de $\Spec(R)$
tel que tout recouvrement fpqc puisse être raffiné par un élément de $A$.
Si $R = k$ est un corps, la raison de cette absence de borne est
qu'il n'existe aucune extension de corps de $k$ qui contienne toute extension de corps
de $k$.

\medskip\noindent
Si l'on néglige les difficultés ensemblistes, on rencontre des préfaisceaux
qui n'admettent pas de faisceau associé, voir
\cite[théorème 5.5]{Waterhouse-fpqc-sheafification}.
Une option d'un certain intérêt consiste à ne considérer que les extensions d'anneaux
fidèlement plates $R \to R'$ où le cardinal de $R'$ est convenablement borné.
(Et si l'on considère tous les schémas d'un univers fixé, comme dans SGA4, on
borne le cardinal par un cardinal fortement inaccessible.)
Il n'est toutefois pas très clair de savoir ce qui se passe si l'on remplace ce cardinal
par un cardinal plus grand.

\medskip\noindent
Pour ces raisons, nous n'introduisons pas de sites fpqc et ne considérons pas
la cohomologie relativement à la topologie fpqc.

\medskip\noindent
En revanche, étant donné un foncteur contravariant
$F : \Sch^{opp} \to \textit{Ensembles}$
il est légitime de demander si $F$ satisfait la propriété de faisceau
pour la topologie fpqc, voir ci-dessous.
On peut en outre s'interroger sur la descente des objets
pour la topologie fpqc, etc. En termes simples, pour certains résultats, le bon
degré de généralité consiste à travailler avec des recouvrements fpqc.

\begin{lemma}
\label{lemma-fpqc}
Soit $T$ un schéma.
\begin{enumerate}
\item Si $T' \to T$ est un isomorphisme, alors $\{T' \to T\}$
est un recouvrement fpqc de $T$.
\item Si $\{T_i \to T\}_{i\in I}$ est un recouvrement fpqc et si, pour tout
$i$, on a un recouvrement fpqc $\{T_{ij} \to T_i\}_{j\in J_i}$, alors
$\{T_{ij} \to T\}_{i \in I, j\in J_i}$ est un recouvrement fpqc.
\item Si $\{T_i \to T\}_{i\in I}$ est un recouvrement fpqc
et si $T' \to T$ est un morphisme de schémas, alors
$\{T' \times_T T_i \to T'\}_{i\in I}$ est un recouvrement fpqc.
\end{enumerate}
\end{lemma}

\begin{proof}
L'assertion (1) est immédiate. Rappelons que le composé de morphismes plats
est plat et que le changement de base d'un morphisme plat est plat
(Morphismes, lemmes \ref{morphisms-lemma-base-change-flat} et
\ref{morphisms-lemma-composition-flat}).
On peut donc appliquer le lemme \ref{lemma-recognize-fpqc-covering}
dans chaque cas pour vérifier que nos familles de morphismes sont des recouvrements fpqc.

\medskip\noindent
Preuve de (2). Supposons que $\{T_i \to T\}_{i\in I}$ soit un recouvrement fpqc et que, pour tout
$i$, on ait un recouvrement fpqc $\{f_{ij} : T_{ij} \to T_i\}_{j\in J_i}$.
Soit $U \subset T$ un ouvert affine. On peut trouver
des ouverts quasi-compacts $U_i \subset T_i$ pour $i \in I$, presque tous vides,
tels que $U = \bigcup f_i(U_i)$. Pour tout $i$, on peut alors choisir
des ouverts quasi-compacts $U_{ij} \subset T_{ij}$ pour $j \in J_i$, presque tous vides,
tels que $U_i = \bigcup_j f_{ij}(U_{ij})$. Ainsi,
$\{T_{ij} \to T\}$ est un recouvrement fpqc.

\medskip\noindent
Preuve de (3). Supposons que $\{T_i \to T\}_{i\in I}$ soit un recouvrement fpqc
et que $T' \to T$ soit un morphisme de schémas. Soit $U' \subset T'$ un ouvert affine
dont l'image est contenue dans l'ouvert affine $U \subset T$. Choisissons
des ouverts quasi-compacts $U_i \subset T_i$, presque tous vides,
tels que $U = \bigcup f_i(U_i)$. Alors $U' \times_U U_i$ est
un ouvert quasi-compact de $T' \times_T T_i$ et
$U' = \bigcup \text{pr}_1(U' \times_U U_i)$. Puisque $T'$
peut être recouvert par de tels ouverts affines $U' \subset T'$, on voit
que $\{T' \times_T T_i \to T'\}_{i\in I}$ est un recouvrement fpqc d'après le
lemme \ref{lemma-recognize-fpqc-covering}.
\end{proof}

\begin{lemma}
\label{lemma-fpqc-affine}
Soit $T$ un schéma affine.
Soit $\{T_i \to T\}_{i \in I}$ un recouvrement fpqc de $T$.
Il existe alors un recouvrement fpqc
$\{U_j \to T\}_{j = 1, \ldots, n}$ qui est un raffinement
de $\{T_i \to T\}_{i \in I}$ tel que chaque $U_j$ soit un schéma
affine. De plus, on peut choisir chaque $U_j$ ouvert affine
dans l'un des $T_i$.
\end{lemma}

\begin{proof}
Cela résulte directement de la définition.
\end{proof}

\begin{definition}
\label{definition-standard-fpqc}
Soit $T$ un schéma affine. Un {\it recouvrement fpqc standard}
de $T$ est une famille $\{f_j : U_j \to T\}_{j = 1, \ldots, n}$
où chaque $U_j$ est affine, plat sur $T$, et $T = \bigcup f_j(U_j)$.
\end{definition}

\noindent
Puisque nous n'introduisons pas le site affine, il faut montrer directement
que la famille de tous les recouvrements fpqc standards satisfait les
axiomes.

\begin{lemma}
\label{lemma-fpqc-affine-axioms}
Soit $T$ un schéma affine.
\begin{enumerate}
\item Si $T' \to T$ est un isomorphisme, alors $\{T' \to T\}$
est un recouvrement fpqc standard de $T$.
\item Si $\{T_i \to T\}_{i\in I}$ est un recouvrement fpqc standard et si, pour tout
$i$, on a un recouvrement fpqc standard $\{T_{ij} \to T_i\}_{j\in J_i}$, alors
$\{T_{ij} \to T\}_{i \in I, j\in J_i}$ est un recouvrement fpqc standard.
\item Si $\{T_i \to T\}_{i\in I}$ est un recouvrement fpqc standard
et si $T' \to T$ est un morphisme de schémas affines, alors
$\{T' \times_T T_i \to T'\}_{i\in I}$ est un recouvrement fpqc standard.
\end{enumerate}
\end{lemma}

\begin{proof}
Cela résulte formellement du fait que les composés et les changements de base
de morphismes plats sont plats
(Morphismes, lemmes \ref{morphisms-lemma-base-change-flat} et
\ref{morphisms-lemma-composition-flat})
et que les produits fibrés de schémas affines sont affines
(Schémas, lemme \ref{schemes-lemma-fibre-product-affines}).
\end{proof}

\begin{lemma}
\label{lemma-fpqc-covering-affines-mapping-in}
Soit $T$ un schéma. Soit $\{f_i : T_i \to T\}_{i \in I}$ une famille de
morphismes de schémas de but $T$. Supposons que
\begin{enumerate}
\item chaque $f_i$ soit plat, et
\item pour tout schéma affine
$Z$ et tout morphisme $h : Z \to T$, il existe un recouvrement fpqc standard
$\{Z_j \to Z\}_{j = 1, \ldots, n}$ qui raffine la famille
$\{T_i \times_T Z \to Z\}_{i \in I}$.
\end{enumerate}
Alors $\{f_i : T_i \to T\}_{i \in I}$ est un recouvrement fpqc de $T$.
\end{lemma}

\begin{proof}
Soit $T = \bigcup U_\alpha$ un recouvrement ouvert affine.
Pour tout $\alpha$, la famille obtenue par image inverse $\{T_i \times_T U_\alpha \to U_\alpha\}$
peut être raffinée par un recouvrement fpqc standard ; c'est donc un
recouvrement fpqc d'après le lemme
\ref{lemma-family-flat-dominated-covering}.
Comme $\{U_\alpha \to T\}$ est un recouvrement fpqc, on en conclut que
$\{T_i \to T\}$ est un recouvrement fpqc d'après le
lemme \ref{lemma-family-flat-fpqc-local-covering}.
\end{proof}

\begin{definition}
\label{definition-sheaf-property-fpqc}
Soit $F$ un foncteur contravariant de la catégorie
des schémas vers celle des ensembles.
\begin{enumerate}
\item Soit $\{U_i \to T\}_{i \in I}$ une famille de morphismes
de schémas de but fixé.
On dit que $F$ {\it satisfait la propriété de faisceau pour la famille donnée}
si, pour toute famille d'éléments $\xi_i \in F(U_i)$ telle que
$\xi_i|_{U_i \times_T U_j} = \xi_j|_{U_i \times_T U_j}$
il existe un unique élément
$\xi \in F(T)$ tel que $\xi_i = \xi|_{U_i}$ dans $F(U_i)$.
\item On dit que $F$ {\it satisfait la propriété de faisceau pour la
topologie fpqc} s'il satisfait la propriété de faisceau pour tout
recouvrement fpqc.
\end{enumerate}
\end{definition}

\noindent
Nous évitons autant que possible de dire « $F$ est un faisceau » dans cette
situation, puisque nous ne définissons pas de catégorie de faisceaux fpqc,
comme nous l'avons expliqué ci-dessus.

\begin{lemma}
\label{lemma-sheaf-property-fpqc}
Soit $F$ un foncteur contravariant de la catégorie
des schémas vers celle des ensembles. Alors $F$ satisfait
la propriété de faisceau pour la topologie fpqc si et seulement
s'il satisfait
\begin{enumerate}
\item la propriété de faisceau pour tout recouvrement de Zariski, et
\item la propriété de faisceau pour tout recouvrement fpqc standard.
\end{enumerate}
De plus, en présence de (1), la propriété (2) équivaut à la
propriété
\begin{enumerate}
\item[(2')] de faisceau pour $\{V \to U\}$,
où $V$, $U$ sont affines et $V \to U$ est fidèlement plat.
\end{enumerate}
\end{lemma}

\begin{proof}
Supposons (1) et (2) vérifiées.
Soit $\{f_i : T_i \to T\}_{i \in I}$ un recouvrement fpqc. Soit $s_i \in F(T_i)$
une famille d'éléments telle que $s_i$ et $s_j$ aient la même image
dans $F(T_i \times_T T_j)$. Soit $W \subset T$ le plus grand ouvert
tel qu'il existe un unique $s \in F(W)$ satisfaisant
$s|_{f_i^{-1}(W)} = s_i|_{f_i^{-1}(W)}$ pour tout $i$.
Un tel plus grand ouvert existe puisque $F$ satisfait la
propriété de faisceau pour les recouvrements de Zariski ; en fait, $W$ est la
réunion de tous les ouverts qui ont cette propriété. Soit $t \in T$.
Montrons que $t \in W$. Pour cela, choisissons un ouvert affine
$t \in U \subset T$ et montrons qu'il existe un unique
$s \in F(U)$ tel que
$s|_{f_i^{-1}(U)} = s_i|_{f_i^{-1}(U)}$ pour tout $i$.

\medskip\noindent
D'après le lemme \ref{lemma-fpqc-affine}, on peut trouver un recouvrement fpqc standard
$\{U_j \to U\}_{j = 1, \ldots, n}$ qui raffine $\{U \times_T T_i \to U\}$,
disons par des morphismes $h_j : U_j \to T_{i_j}$. D'après (2), on obtient un unique élément
$s \in F(U)$ tel que $s|_{U_j} = F(h_j)(s_{i_j})$. Remarquons que, pour tout
schéma $V \to U$ sur $U$, il existe une unique section $s_V \in F(V)$
dont la restriction vaut $F(h_j \circ \text{pr}_2)(s_{i_j})$ sur
$V \times_U U_j$ pour $j = 1, \ldots, n$. En effet, c'est vrai si $V$
est affine d'après (2), puisque $\{V \times_U U_j \to V\}$ est un recouvrement fpqc standard,
et le cas général résulte de (1) et du cas affine en choisissant un
recouvrement ouvert affine de $V$. En particulier, $s_V = s|_V$.
En prenant maintenant $V = U \times_T T_i$ et en utilisant
$s_{i_j}|_{T_{i_j} \times_T T_i} = s_i|_{T_{i_j} \times_T T_i}$,
on en conclut que $s|_{U \times_T T_i} = s_V = s_i|_{U \times_T T_i}$,
ce qu'il fallait démontrer.

\medskip\noindent
Preuve de l'équivalence de (2) et (2') en présence de (1).
Si $\{T_i \to T\}$ est un recouvrement fpqc standard, alors
$\coprod T_i \to T$ est un morphisme fidèlement plat de schémas affines.
En présence de (1), on a $F(\coprod T_i) = \prod F(T_i)$
et, de même,
$F((\coprod T_i) \times_T (\coprod T_i)) = \prod F(T_i \times_T T_{i'})$.
Ainsi, la condition de faisceau pour $\{T_i \to T\}$ et pour $\{\coprod T_i \to T\}$
est la même.
\end{proof}

\noindent
Le lemme suivant sert seulement à signaler que les difficultés ensemblistes
se présentent bel et bien et que la plupart des lecteurs devraient les ignorer.

\begin{lemma}
\label{lemma-no-set-of-fpqc-covers-is-initial}
Soit $R$ un anneau non nul. Il n'existe aucun ensemble $A$ de
recouvrements fpqc de $\Spec(R)$ tel que tout recouvrement fpqc puisse
être raffiné par un élément de $A$.
\end{lemma}

\begin{proof}
Expliquons d'abord le cas où $R = k$ est un corps. Pour tout ensemble $I$, considérons
l'extension transcendante pure
$k_I = k(\{t_i\}_{i \in I})/k$. Comme $k \to k_I$ est fidèlement plat,
on voit que $\{\Spec(k_I) \to \Spec(k)\}$ est un recouvrement fpqc.
Soit $A$ un ensemble et, pour tout $\alpha \in A$, soit
$\mathcal{U}_\alpha = \{S_{\alpha, j} \to \Spec(k)\}_{j \in J_\alpha}$ un
recouvrement fpqc. Si $\mathcal{U}_\alpha$ raffine $\{\Spec(k_I) \to \Spec(k)\}$,
alors les morphismes $S_{\alpha, j} \to \Spec(k)$ se factorisent par
$\Spec(k_I)$. Comme $\mathcal{U}_\alpha$ est un recouvrement,
au moins un certain $S_{\alpha, j}$ est non vide. Choisissons un
point $s \in S_{\alpha, j}$. Puisque l'on a la factorisation
$S_{\alpha, j} \to \Spec(k_I) \to \Spec(k)$
on obtient un homomorphisme de corps $k_I \to \kappa(s)$.
En particulier, on voit que la cardinalité de $\kappa(s)$
est au moins celle de $I$. Ainsi, si l'on choisit pour $I$ un ensemble
de cardinalité supérieure à celles des corps résiduels
de tous les schémas $S_{\alpha, j}$, une telle factorisation
n'existe pas et le lemme est démontré pour $R = k$.

\medskip\noindent
Cas général. Comme $R$ est non nul, il possède un idéal premier maximal
$\mathfrak m$ de corps résiduel $\kappa$. Soit $I$ un ensemble et
considérons $R_I = S_I^{-1} R[\{t_i\}_{i \in I}]$,
où $S_I \subset R[\{t_i\}_{i \in I}]$ est la partie multiplicative
formée des $f \in R[\{t_i\}_{i \in I}]$ tels que l'image de $f$ soit
un élément non nul de $R/\mathfrak p[\{t_i\}_{i \in I}]$ pour
tout idéal premier $\mathfrak p$ de $R$. Alors $R_I$ est une $R$-algèbre
fidèlement plate et $\{\Spec(R_I) \to \Spec(R)\}$ est un
recouvrement fpqc. Nous laissons au lecteur le soin de montrer que
$R_I \otimes_R \kappa \cong \kappa(\{t_i\}_{i \in I}) = \kappa_I$
avec les notations ci-dessus (indication : utiliser la surjectivité de $R \to \kappa$
et le fait que tout $f \in R[\{t_i\}_{i \in I}]$ dont l'un des monômes apparaît
avec le coefficient $1$ appartient à $S_I$). Soit $A$ un ensemble et,
pour tout $\alpha \in A$, soit
$\mathcal{U}_\alpha = \{S_{\alpha, j} \to \Spec(R)\}_{j \in J_\alpha}$ un
recouvrement fpqc. Si $\mathcal{U}_\alpha$ raffine $\{\Spec(R_I) \to \Spec(R)\}$,
alors, par changement de base, on en conclut que
$\{S_{\alpha, j} \times_{\Spec(R)} \Spec(\kappa) \to \Spec(\kappa)\}$
raffine $\{\Spec(\kappa_I) \to \Spec(\kappa)\}$.
Par le résultat du paragraphe précédent, il existe donc un $I$
pour lequel ce n'est pas le cas, ce qui démontre le lemme.
\end{proof}












\section{La topologie V}
\label{section-V}

\noindent
La topologie V est plus fine que toutes les autres topologies de ce chapitre.
Grosso modo, elle est engendrée par les recouvrements de Zariski et
par les morphismes quasi-compacts satisfaisant une propriété de relèvement
des spécialisations (lemme \ref{lemma-refine-qcqs-V}).
Cependant, le procédé que nous emploierons pour définir les recouvrements V est un peu différent.
Nous définirons d'abord les recouvrements V standards des schémas affines, puis
nous les utiliserons pour définir les recouvrements V en général. Remarque typographique :
dans la littérature, on emploie parfois « $v$-recouvrement » au lieu de « recouvrement V ».

\begin{definition}
\label{definition-standard-V-covering}
Soit $T$ un schéma affine. Un {\it recouvrement V standard} est une famille finie
$\{T_j \to T\}_{j = 1, \ldots, m}$, où les $T_j$ sont affines,
telle que, pour tout morphisme $g : \Spec(V) \to T$, où $V$
est un anneau de valuation, il existe une extension $V \subset W$ d'anneaux de valuation
(Compléments d'algèbre, définition
\ref{more-algebra-definition-extension-valuation-rings}),
un indice $1 \leq j \leq m$ et un diagramme commutatif
$$
\xymatrix{
\Spec(W) \ar[r] \ar[d] & T_j \ar[d] \\
\Spec(V) \ar[r]^g & T
}
$$
\end{definition}

\noindent
Nous démontrons d'abord quelques lemmes élémentaires sur cette notion.

\begin{lemma}
\label{lemma-standard-fpqc-standard-V}
Tout recouvrement fpqc standard est un recouvrement V standard.
\end{lemma}

\begin{proof}
Soit $\{X_i \to X\}_{i = 1, \ldots, n}$ un recouvrement fpqc standard
(définition \ref{definition-standard-fpqc}). Soit $g : \Spec(V) \to X$
un morphisme, où $V$ est un anneau de valuation. Soit $x \in X$
l'image du point fermé de $\Spec(V)$. Choisissons un $i$ et un point
$x_i \in X_i$ d'image $x$. Alors $\Spec(V) \times_X X_i$
possède un point $x'_i$ dont l'image est le point fermé de $\Spec(V)$.
Comme $\Spec(V) \times_X X_i \to \Spec(V)$ est plat,
on peut trouver une spécialisation $x''_i \leadsto x'_i$ de
points de $\Spec(V) \times_X X_i$, où $x''_i$ s'envoie sur
le point générique de $\Spec(V)$ ; voir
Morphismes, lemme \ref{morphisms-lemma-generalizations-lift-flat}. D'après
Schémas, lemme \ref{schemes-lemma-points-specialize},
on peut choisir un anneau de valuation $W$ et un morphisme
$h : \Spec(W) \to \Spec(V) \times_X X_i$ tels que $h$
envoie le point générique de $\Spec(W)$ sur $x''_i$ et le
point fermé de $\Spec(W)$ sur $x'_i$. On obtient un
diagramme commutatif
$$
\xymatrix{
\Spec(W) \ar[r] \ar[d] & X_i \ar[d] \\
\Spec(V) \ar[r] & X
}
$$
où $V \to W$ est une extension d'anneaux de valuation.
Cela démontre le lemme.
\end{proof}

\begin{lemma}
\label{lemma-standard-ph-standard-V}
Tout recouvrement ph standard est un recouvrement V standard.
\end{lemma}

\begin{proof}
Soit $T$ un schéma affine. Soit $f : U \to T$ un morphisme propre
surjectif. Soit $U = \bigcup_{j = 1, \ldots, m} U_j$ un recouvrement ouvert
affine fini. Il faut montrer que $\{U_j \to T\}$ est un recouvrement V
standard ; voir la définition
\ref{definition-standard-ph-covering}.
Soit $g : \Spec(V) \to T$
un morphisme, où $V$ est un anneau de valuation de corps des fractions $K$.
Comme $U \to T$ est surjectif, on peut choisir une extension de corps
$L/K$ et un diagramme commutatif
$$
\xymatrix{
\Spec(L) \ar[rr] \ar[d] & & U \ar[d] \\
\Spec(K) \ar[r] & \Spec(V) \ar[r]^g & T
}
$$
D'après Algèbre, lemme \ref{algebra-lemma-dominate}, on peut choisir un anneau
de valuation $W \subset L$ dominant $V$. Le critère valuatif de
propreté (Morphismes, lemme \ref{morphisms-lemma-characterize-proper})
fournit alors le morphisme $h$ dans le diagramme
commutatif
$$
\xymatrix{
\Spec(L) \ar[r] \ar[d] & \Spec(W) \ar[r]_h \ar[d] & U \ar[d] \\
\Spec(K) \ar[r] & \Spec(V) \ar[r]^g & X
}
$$
Comme $\Spec(W)$ possède un unique point fermé, on voit que $\Im(h)$
est contenue dans $U_j$ pour un certain $j$. Ainsi, $h : \Spec(W) \to U_j$
est le relèvement cherché, et $\{U_j \to T\}$ est donc un recouvrement V standard.
\end{proof}

\begin{lemma}
\label{lemma-base-change-standard-V}
Soit $\{T_j \to T\}_{j = 1, \ldots, m}$ un recouvrement V standard.
Soit $T' \to T$ un morphisme de schémas affines.
Alors $\{T_j \times_T T' \to T'\}_{j = 1, \ldots, m}$ est un
recouvrement V standard.
\end{lemma}

\begin{proof}
Soit $\Spec(V) \to T'$ un morphisme, où $V$ est un anneau de valuation.
Par hypothèse, on peut trouver une extension d'anneaux de valuation $V \subset W$,
un $i$ et un diagramme commutatif
$$
\xymatrix{
\Spec(W) \ar[r] \ar[d] & T_i \ar[d] \\
\Spec(V) \ar[r] & T
}
$$
La propriété universelle des produits fibrés
donne le morphisme $\Spec(W) \to T' \times_T T_i$
voulu.
\end{proof}

\begin{lemma}
\label{lemma-composition-standard-V}
Soit $T$ un schéma affine. Soit $\{T_j \to T\}_{j = 1, \ldots, m}$
un recouvrement V standard. Soit $\{T_{ji} \to T_j\}_{i = 1, \ldots n_j}$
un recouvrement V standard. Alors $\{T_{ji} \to T\}_{i, j}$
est un recouvrement V standard.
\end{lemma}

\begin{proof}
Cela résulte formellement de l'observation suivante : si
$V \subset W$ et $W \subset \Omega$ sont des extensions
d'anneaux de valuation, alors $V \subset \Omega$ est une extension
d'anneaux de valuation.
\end{proof}

\begin{lemma}
\label{lemma-refine-standard-V}
Soit $T$ un schéma affine. Soit $\{T_j \to T\}_{j = 1, \ldots, m}$
une famille de morphismes telle que $T_j$ soit affine pour tout $j$.
Les conditions suivantes sont équivalentes :
\begin{enumerate}
\item $\{T_j \to T\}_{j = 1, \ldots, m}$ est un recouvrement V standard ;
\item il existe un recouvrement V standard qui raffine
$\{T_j \to T\}_{j = 1, \ldots, m}$ ;
\item $\{\coprod_{j = 1, \ldots, m} T_j \to T\}$ est un recouvrement
V standard.
\end{enumerate}
\end{lemma}

\begin{proof}
Omis. Indications : cela résulte presque immédiatement de la définition.
Le seul point qui mérite quelque attention est qu'un
morphisme du spectre d'un anneau local vers
$\coprod_{j = 1, \ldots, m} T_j$ se factorise nécessairement par l'un des $T_j$.
\end{proof}

\begin{definition}
\label{definition-V-covering}
Soit $T$ un schéma. Un {\it recouvrement V de $T$} est une famille
de morphismes de schémas $\{T_i \to T\}_{i \in I}$ telle que, pour tout
ouvert affine $U \subset T$, il existe un recouvrement V standard
$\{U_j \to U\}_{j = 1, \ldots, m}$ qui raffine la famille
$\{T_i \times_T U \to U\}_{i \in I}$.
\end{definition}

\noindent
La topologie V présente les mêmes problèmes ensemblistes que
la topologie fpqc. Nous nous abstiendrons donc de définir des sites V
et nous ne considérerons pas de cohomologie pour la topologie V.
En revanche, étant donné un foncteur
$F : \Sch^{opp} \to \textit{Ensembles}$
il est légitime de demander si $F$ satisfait la propriété de faisceau
pour la topologie V ; voir ci-dessous.
On peut en outre s'interroger sur la descente d'objets
pour la topologie V, etc.

\begin{lemma}
\label{lemma-refine-by-V}
Soit $T$ un schéma. Soit $\{f_i : T_i \to T\}_{i \in I}$ une famille
de morphismes. Les conditions suivantes sont équivalentes :
\begin{enumerate}
\item $\{T_i \to T\}_{i \in I}$ est un recouvrement V ;
\item il existe un recouvrement V qui raffine $\{T_i \to T\}_{i \in I}$ ;
\item $\{\coprod_{i \in I} T_i \to T\}$ est un recouvrement V.
\end{enumerate}
\end{lemma}

\begin{proof}
Omis. Indication : comparer avec la démonstration du
lemme \ref{lemma-refine-by-ph}.
\end{proof}

\begin{lemma}
\label{lemma-V}
Soit $T$ un schéma.
\begin{enumerate}
\item Si $T' \to T$ est un isomorphisme, alors $\{T' \to T\}$
est un recouvrement V de $T$.
\item Si $\{T_i \to T\}_{i\in I}$ est un recouvrement V et si, pour chaque
$i$, on a un recouvrement V $\{T_{ij} \to T_i\}_{j\in J_i}$, alors
$\{T_{ij} \to T\}_{i \in I, j\in J_i}$ est un recouvrement V.
\item Si $\{T_i \to T\}_{i\in I}$ est un recouvrement V
et si $T' \to T$ est un morphisme de schémas, alors
$\{T' \times_T T_i \to T'\}_{i\in I}$ est un recouvrement V.
\end{enumerate}
\end{lemma}

\begin{proof}
L'assertion (1) est claire.

\medskip\noindent
Démonstration de (3).
Soit $U' \subset T'$ un sous-schéma ouvert affine.
Comme $U'$ est quasi-compact, on peut trouver un recouvrement ouvert affine
fini $U' = U'_1 \cup \ldots \cup U'$ tel que
$U'_j \to T$ se factorise par un ouvert affine $U_j \subset T$.
Choisissons un recouvrement V standard $\{U_{jl} \to U_j\}_{l = 1, \ldots, n_j}$
qui raffine $\{T_i \times_T U_j \to U_j\}$.
D'après le lemme \ref{lemma-base-change-standard-V},
le changement de base $\{U_{jl} \times_{U_j} U'_j \to U'_j\}$
est un recouvrement V standard. Notons que $\{U'_j \to U'\}$
est un recouvrement V standard
(par exemple d'après le lemme \ref{lemma-standard-fpqc-standard-V}).
D'après le lemme \ref{lemma-composition-standard-V}, 
la famille $\{U_{jl} \times_{U_j} U'_j \to U'\}$
est un recouvrement V standard. Comme
$\{U_{jl} \times_{U_j} U'_j \to U'\}$ raffine $\{T_i \times_T U' \to U'\}$,
on conclut.

\medskip\noindent
Démonstration de (2).
Soit $U \subset T$ un ouvert affine. Choisissons d'abord un recouvrement V standard
$\{U_k \to U\}_{k = 1, \ldots, m}$ qui raffine
$\{T_i \times_T U \to U\}$.
Supposons que le raffinement soit donné par des morphismes $U_k \to T_{i_k}$ au-dessus de $T$.
Alors
$$
\{T_{i_kj} \times_{T_{i_k}} U_k \to U_k\}_{j \in J_{i_k}}
$$
est un recouvrement V d'après (3). Comme $U_k$ est affine,
on peut trouver un recouvrement V standard
$\{U_{ka} \to U_k\}_{a = 1, \ldots, b_k}$ qui raffine cette famille.
On applique alors le lemme \ref{lemma-composition-standard-V}
pour voir que $\{U_{ka} \to U\}$ est un recouvrement V standard qui
raffine $\{T_{ij} \times_T U \to U\}$. Cela achève la démonstration.
\end{proof}

\begin{lemma}
\label{lemma-zariski-etale-smooth-syntomic-fppf-fpqc-ph-V}
Tout recouvrement fpqc est un recouvrement V. A fortiori,
tout recouvrement fppf, syntomique, lisse, étale ou de Zariski est un recouvrement V.
De même, tout recouvrement ph est un recouvrement V.
\end{lemma}

\begin{proof}
Tout recouvrement fpqc peut, localement sur des ouverts affines, être raffiné par un
recouvrement fpqc standard ; voir le lemme \ref{lemma-fpqc-affine}.
Tout recouvrement fpqc standard est un recouvrement V standard ; voir le
lemme \ref{lemma-standard-fpqc-standard-V}. La première assertion résulte donc
de notre définition des recouvrements V au moyen des recouvrements V standards.
La conclusion pour les recouvrements fppf, syntomiques, lisses, étales ou de Zariski
résulte de ce que ceux-ci sont des recouvrements fpqc ; voir le
lemme \ref{lemma-zariski-etale-smooth-syntomic-fppf-fpqc}.

\medskip\noindent
L'assertion sur les recouvrements ph résulte de la même manière du
lemme \ref{lemma-standard-ph-standard-V}.
\end{proof}

\begin{definition}
\label{definition-sheaf-property-V}
Soit $F$ un foncteur contravariant de la catégorie
des schémas dans celle des ensembles. On dit que
$F$ {\it satisfait la propriété de faisceau pour la topologie V}
s'il satisfait la propriété de faisceau pour tout recouvrement V
(voir la définition \ref{definition-sheaf-property-fpqc}).
\end{definition}

\noindent
Nous évitons d'employer la terminologie « $F$ est un faisceau » dans cette
situation, puisque nous ne définissons pas de catégorie des faisceaux V,
comme nous l'avons expliqué plus haut.

\begin{lemma}
\label{lemma-sheaf-property-V}
Soit $F$ un foncteur contravariant de la catégorie
des schémas dans celle des ensembles. Alors $F$ satisfait
la propriété de faisceau pour la topologie V si et seulement
s'il satisfait aux deux conditions suivantes :
\begin{enumerate}
\item la propriété de faisceau pour tout recouvrement de Zariski ;
\item la propriété de faisceau pour tout recouvrement V standard.
\end{enumerate}
En outre, sous l'hypothèse (1), la propriété (2) équivaut à la
propriété suivante :
\begin{enumerate}
\item[(2')] la propriété de faisceau pour un recouvrement V standard
de la forme $\{V \to U\}$, c'est-à-dire constitué d'une seule flèche.
\end{enumerate}
\end{lemma}

\begin{proof}
Supposons (1) et (2) satisfaites.
Soit $\{f_i : T_i \to T\}_{i \in I}$ un recouvrement V. Soit $s_i \in F(T_i)$
une famille d'éléments telle que $s_i$ et $s_j$ s'envoient sur le même élément
de $F(T_i \times_T T_j)$. Soit $W \subset T$ le plus grand ouvert
tel qu'il existe un unique $s \in F(W)$ vérifiant
$s|_{f_i^{-1}(W)} = s_i|_{f_i^{-1}(W)}$ pour tout $i$.
Un tel plus grand ouvert existe parce que $F$ satisfait la
propriété de faisceau pour les recouvrements de Zariski ; en fait, $W$ est la
réunion de tous les ouverts possédant cette propriété. Soit $t \in T$.
Nous allons montrer que $t \in W$. Pour cela, choisissons un ouvert affine
$t \in U \subset T$ et montrons qu'il existe un unique
$s \in F(U)$ tel que
$s|_{f_i^{-1}(U)} = s_i|_{f_i^{-1}(U)}$ pour tout $i$.

\medskip\noindent
On peut trouver un recouvrement V standard
$\{U_j \to U\}_{j = 1, \ldots, n}$ qui raffine $\{U \times_T T_i \to U\}$,
disons au moyen de morphismes $h_j : U_j \to T_{i_j}$. D'après (2), on obtient un unique élément
$s \in F(U)$ tel que $s|_{U_j} = F(h_j)(s_{i_j})$. Notons que, pour tout
schéma $V \to U$ au-dessus de $U$, il existe une unique section $s_V \in F(V)$
qui se restreint en $F(h_j \circ \text{pr}_2)(s_{i_j})$ sur
$V \times_U U_j$ pour $j = 1, \ldots, n$. En effet, cela est vrai si $V$
est affine d'après (2), puisque $\{V \times_U U_j \to V\}$ est un recouvrement V standard
(lemme \ref{lemma-base-change-standard-V}) ;
dans le cas général, cela résulte de (1) et du cas affine en choisissant un
recouvrement ouvert affine de $V$. En particulier, $s_V = s|_V$.
En prenant maintenant $V = U \times_T T_i$ et en utilisant l'égalité
$s_{i_j}|_{T_{i_j} \times_T T_i} = s_i|_{T_{i_j} \times_T T_i}$
on conclut que $s|_{U \times_T T_i} = s_V = s_i|_{U \times_T T_i}$,
ce qu'il fallait démontrer.

\medskip\noindent
Démonstration de l'équivalence de (2) et (2') sous l'hypothèse (1).
Supposons que $\{T_i \to T\}_{i = 1, \ldots, n}$ soit un recouvrement V standard. Alors
$\coprod_{i = 1, \ldots, n} T_i \to T$ est un morphisme de schémas affines
qui est manifestement aussi un recouvrement V standard.
Sous l'hypothèse (1), on a $F(\coprod T_i) = \prod F(T_i)$
et, de même,
$F((\coprod T_i) \times_T (\coprod T_i)) = \prod F(T_i \times_T T_{i'})$.
Ainsi, la condition de faisceau pour $\{T_i \to T\}$ et pour $\{\coprod T_i \to T\}$
est la même.
\end{proof}

\noindent
Le lemme suivant montre que la propriété d'être un recouvrement V est
liée à la possibilité de relever les spécialisations.

\begin{lemma}
\label{lemma-refine-qcqs-V}
Soit $X \to Y$ un morphisme quasi-compact de schémas.
Les conditions suivantes sont équivalentes :
\begin{enumerate}
\item $\{X \to Y\}$ est un recouvrement V ;
\item pour tout anneau de valuation $V$ et tout morphisme $g : \Spec(V) \to Y$,
il existe une extension d'anneaux de valuation $V \subset W$
et un diagramme commutatif
$$
\xymatrix{
\Spec(W) \ar[r] \ar[d] & X \ar[d] \\
\Spec(V) \ar[r] & Y
}
$$
\item pour tout morphisme $Z \to Y$ et toute spécialisation $z' \leadsto z$
de points de $Z$, il existe une spécialisation $w' \leadsto w$
de points de $Z \times_Y X$ qui s'envoie sur $z' \leadsto z$.
\end{enumerate}
\end{lemma}

\begin{proof}
Supposons (1) et soit $g : \Spec(V) \to Y$ comme dans (2).
Comme $V$ est un anneau local, il existe un ouvert affine $U \subset Y$
tel que $g$ se factorise par $U$. D'après la définition \ref{definition-V-covering},
on peut trouver un recouvrement V standard $\{U_j \to U\}$ qui raffine
$\{X \times_Y U \to U\}$. D'après la définition \ref{definition-standard-V-covering},
on peut trouver un $j$, une extension d'anneaux de valuation $V \subset W$
et un diagramme commutatif
$$
\xymatrix{
\Spec(W) \ar[r] \ar[d] & U_j \ar[d] \ar@{..>}[r] & X \ar[ld] \\
\Spec(V) \ar[r] & Y
}
$$
La propriété de raffinement du recouvrement fournit la flèche pointillée
qui rend le diagramme commutatif, et l'on voit que (2) est satisfaite.

\medskip\noindent
Supposons (2) et soient $Z \to Y$ et $z' \leadsto z$ comme dans (3).
D'après Schémas, lemme \ref{schemes-lemma-points-specialize},
on peut trouver un anneau de valuation $V$ et un morphisme $\Spec(V) \to Z$
tels que le point fermé de $\Spec(V)$ s'envoie sur $z$ et que le
point générique de $\Spec(V)$ s'envoie sur $z'$. D'après (2), on peut trouver une
extension d'anneaux de valuation $V \subset W$ et un
diagramme commutatif
$$
\xymatrix{
\Spec(W) \ar[rr] \ar[d] & & X \ar[d] \\
\Spec(V) \ar[r] & Z \ar[r] & Y
}
$$
Les points générique et fermé de $\Spec(W)$ s'envoient sur des points
$w' \leadsto w$ de $Z \times_Y X$ par le
morphisme induit $\Spec(W) \to Z \times_Y X$.
Cela montre que (3) est satisfaite.

\medskip\noindent
Supposons (3) satisfaite et soit $U \subset Y$ un ouvert affine.
Choisissons un recouvrement ouvert affine fini
$U \times_Y X = \bigcup_{j = 1, \ldots, m} U_j$.
C'est possible puisque $X \to Y$ est quasi-compact.
Nous affirmons que $\{U_j \to U\}$ est un recouvrement V standard.
Cette affirmation entraîne (1) et achève la démonstration du lemme.
Pour la démontrer, soit $V$ un anneau de valuation
et soit $g : \Spec(V) \to U$ un morphisme.
D'après (3), on trouve une spécialisation $w' \leadsto w$ de
points de
$$
T = \Spec(V) \times_X Y = \Spec(V) \times_U (U \times_X Y)
$$
telle que $w'$ s'envoie sur le point générique de $\Spec(V)$
et que $w$ s'envoie sur le point fermé de $\Spec(V)$.
D'après Schémas, lemme \ref{schemes-lemma-points-specialize},
on peut trouver un anneau de valuation $W$ et un morphisme
$\Spec(W) \to T$ tels que le point générique de $\Spec(W)$
s'envoie sur $w'$ et que le point fermé de $\Spec(W)$ s'envoie sur $w$.
La composée $\Spec(W) \to T \to \Spec(V)$ correspond
à une inclusion $V \subset W$ qui fait de $W$ une
extension de l'anneau de valuation $V$.
Comme $T = \bigcup \Spec(V) \times_U U_j$ est un recouvrement
ouvert, on voit que $\Spec(W) \to T$ se factorise par
$\Spec(V) \times_U U_j$ pour un certain $j$. On obtient donc
un diagramme commutatif
$$
\xymatrix{
\Spec(W) \ar[d] \ar[r] & U_j \ar[d] \\
\Spec(V) \ar[r] & U
}
$$
et la démonstration de l'affirmation est achevée.
\end{proof}

\noindent
Tout recouvrement V fournit une famille de morphismes universellement submersive.
La réciproque de ce lemme est fausse ; voir
Exemples, section \ref{examples-section-universally-submersive-not-V}.

\begin{lemma}
\label{lemma-V-covering-universally-submersive}
Soit $\{f_i : X_i \to X\}_{i \in I}$ un recouvrement V.
Alors
$$
\coprod\nolimits_{i \in I} f_i :
\coprod\nolimits_{i \in I} X_i
\longrightarrow
X
$$
est un morphisme de schémas universellement submersif (Morphismes, définition
\ref{morphisms-definition-submersive}).
\end{lemma}

\begin{proof}
Nous utiliserons sans autre mention que le changement de base d'un recouvrement V
est un recouvrement V (lemme \ref{lemma-V}). Il suffit donc
de montrer que le morphisme est submersif.
Le caractère submersif est manifestement local pour la topologie de Zariski sur la base.
On peut donc supposer $X$ affine. Alors $\{X_i \to X\}$
peut être raffiné par un recouvrement V standard $\{Y_j \to X\}$.
S'il est possible de montrer que $\coprod Y_j \to X$ est submersif,
alors, puisqu'il existe une factorisation $\coprod Y_j \to \coprod X_i \to X$,
on conclut que $\coprod X_i \to X$ est submersif.
Posons $Y = \coprod Y_j$ et considérons le morphisme de schémas affines
$f : Y \to X$.
D'après le lemme \ref{lemma-refine-qcqs-V}, on sait que toute
spécialisation $x' \leadsto x$ dans $X$ se relève en une spécialisation
$y' \leadsto y$ dans $Y$.
Ainsi, si $T \subset X$ est une partie telle que $f^{-1}(T)$
soit fermée dans $Y$, alors $T \subset X$ est stable par spécialisation.
Comme $f^{-1}(T) \subset Y$, muni de la structure réduite
de sous-schéma fermé induite, est un schéma affine, on conclut que
$T \subset X$ est fermé d'après Algèbre, lemme
\ref{algebra-lemma-image-stable-specialization-closed}.
Ainsi, $f$ est submersif.
\end{proof}





















\section{Changement de topologies}
\label{section-change-topologies}

\noindent
Soit $f : X \to Y$ un morphisme de schémas au-dessus d'un schéma de base $S$.
On dispose alors des morphismes de sites suivants\footnote{Nous n'avons
pas inclus la comparaison entre la topologie ph et les autres ;
voir à ce sujet
Compléments sur les morphismes, remarque \ref{more-morphisms-remark-change-topologies}.}
(moyennant des choix convenables de sites comme dans la remarque
\ref{remark-choice-sites} ci-dessous) :
\begin{enumerate}
\item $(\Sch/X)_{fppf} \longrightarrow (\Sch/Y)_{fppf}$,
\item $(\Sch/X)_{fppf} \longrightarrow (\Sch/Y)_{syntomic}$,
\item $(\Sch/X)_{fppf} \longrightarrow (\Sch/Y)_{smooth}$,
\item $(\Sch/X)_{fppf} \longrightarrow
(\Sch/Y)_\etale$,
\item $(\Sch/X)_{fppf} \longrightarrow (\Sch/Y)_{Zar}$,
\item $(\Sch/X)_{syntomic} \longrightarrow (\Sch/Y)_{syntomic}$,
\item $(\Sch/X)_{syntomic} \longrightarrow (\Sch/Y)_{smooth}$,
\item $(\Sch/X)_{syntomic} \longrightarrow
(\Sch/Y)_\etale$,
\item $(\Sch/X)_{syntomic} \longrightarrow (\Sch/Y)_{Zar}$,
\item $(\Sch/X)_{smooth} \longrightarrow (\Sch/Y)_{smooth}$,
\item $(\Sch/X)_{smooth} \longrightarrow
(\Sch/Y)_\etale$,
\item $(\Sch/X)_{smooth} \longrightarrow (\Sch/Y)_{Zar}$,
\item $(\Sch/X)_\etale \longrightarrow
(\Sch/Y)_\etale$,
\item $(\Sch/X)_\etale \longrightarrow (\Sch/Y)_{Zar}$,
\item $(\Sch/X)_{Zar} \longrightarrow (\Sch/Y)_{Zar}$,
\item $(\Sch/X)_{fppf} \longrightarrow Y_\etale$,
\item $(\Sch/X)_{syntomic} \longrightarrow Y_\etale$,
\item $(\Sch/X)_{smooth} \longrightarrow Y_\etale$,
\item $(\Sch/X)_\etale \longrightarrow Y_\etale$,
\item $(\Sch/X)_{fppf} \longrightarrow Y_{Zar}$,
\item $(\Sch/X)_{syntomic} \longrightarrow Y_{Zar}$,
\item $(\Sch/X)_{smooth} \longrightarrow Y_{Zar}$,
\item $(\Sch/X)_\etale \longrightarrow Y_{Zar}$,
\item $(\Sch/X)_{Zar} \longrightarrow Y_{Zar}$,
\item $X_\etale \longrightarrow Y_\etale$,
\item $X_\etale \longrightarrow Y_{Zar}$,
\item $X_{Zar} \longrightarrow Y_{Zar}$,
\end{enumerate}
Dans chaque cas, le foncteur continu sous-jacent
$\Sch/Y \to \Sch/X$, ou
$Y_\tau \to \Sch/X$
est le foncteur $Y'/Y \mapsto X \times_Y Y'/X$. En effet, dans les sections
précédentes, nous avons construit les morphismes
$f_{big} : (\Sch/X)_\tau \to (\Sch/Y)_\tau$
et
$f_{small} : X_\tau \to Y_\tau$
pour $\tau$ comme ci-dessus.
Nous avons aussi construit les morphismes de sites
$\pi_Y : (\Sch/Y)_\tau \to Y_\tau$ pour
$\tau \in \{\etale, Zariski\}$.
D'autre part, il est clair que le foncteur identité
$(\Sch/X)_\tau \to (\Sch/X)_{\tau'}$ définit
un morphisme de sites lorsque $\tau$ est une topologie plus fine que
$\tau'$. En composant ces morphismes, on obtient donc la liste des morphismes possibles
ci-dessus.

\medskip\noindent
La description simple du foncteur sous-jacent montre clairement que, pour des
morphismes de schémas $X \to Y \to Z$, la composée de deux des
morphismes de sites ci-dessus, par exemple
$$
(\Sch/X)_{\tau_0} \longrightarrow
(\Sch/Y)_{\tau_1} \longrightarrow
(\Sch/Z)_{\tau_2}
$$
est le morphisme de sites correspondant associé au morphisme de
schémas $X \to Z$.

\begin{remark}
\label{remark-choice-sites}
Prenons une catégorie $\Sch_\alpha$ construite comme dans
Ensembles, lemme \ref{sets-lemma-construct-category},
en partant de l'ensemble de schémas $\{X, Y, S\}$. Choisissons un ensemble de
recouvrements $\text{Cov}_{fppf}$ sur $\Sch_\alpha$ comme dans
Ensembles, lemme \ref{sets-lemma-coverings-site},
en partant de la catégorie $\Sch_\alpha$ et de la classe des recouvrements
fppf. Notons $\Sch_{fppf}$ le gros site fppf ainsi
obtenu. Ensuite, pour $\tau \in \{Zariski, \etale, smooth, syntomic\}$,
munissons $\Sch_\tau$ de la même catégorie sous-jacente que
$\Sch_{fppf}$, avec pour recouvrements
$\text{Cov}_\tau \subset \text{Cov}_{fppf}$, simplement le sous-ensemble des
recouvrements $\tau$. On vérifie immédiatement que cela définit
un gros site $\Sch_\tau$.
\end{remark}









\section{Changement de gros sites}
\label{section-change-alpha}

\noindent
Dans cette section, nous expliquons ce qui se passe lorsqu'on change de gros
site de Zariski, fppf ou étale.

\medskip\noindent
Soient $\tau, \tau' \in \{Zariski, \etale, smooth, syntomic, fppf\}$.
Étant donnés deux gros sites $\Sch_\tau$ et $\Sch'_{\tau'}$,
on dit que
{\it $\Sch_\tau$ est contenu dans $\Sch'_{\tau'}$} si
$\Ob(\Sch_\tau) \subset \Ob(\Sch'_{\tau'})$
et
$\text{Cov}(\Sch_\tau) \subset \text{Cov}(\Sch'_{\tau'})$.
Dans ce cas, $\tau$ est plus fine que $\tau'$ ; par exemple, aucun site
fppf ne peut être contenu dans un site étale.

\begin{lemma}
\label{lemma-contained-in}
Tout ensemble de gros sites de Zariski est contenu dans un même gros site de Zariski.
Il en va de même, mutatis mutandis, pour les gros sites fppf et les gros sites étales.
\end{lemma}

\begin{proof}
Cela résulte de ce que la réunion d'un ensemble d'ensembles est un ensemble et que les
constructions de Ensembles, lemmes \ref{sets-lemma-construct-category} et
\ref{sets-lemma-coverings-site}
permettent de partir de n'importe quel ensemble initialement donné de schémas
et de recouvrements.
\end{proof}

\begin{lemma}
\label{lemma-change-alpha}
Soit $\tau \in \{Zariski, \etale, smooth, syntomic, fppf\}$.
Supposons donnés deux gros sites $\Sch_\tau$ et $\Sch'_\tau$.
Supposons que $\Sch_\tau$ soit contenu dans $\Sch'_\tau$.
Le foncteur d'inclusion $\Sch_\tau \to \Sch'_\tau$ satisfait
aux hypothèses de Sites, lemme \ref{sites-lemma-bigger-site}.
Il existe des morphismes de topos
\begin{eqnarray*}
g : \Sh(\Sch_\tau) &
\longrightarrow &
\Sh(\Sch'_\tau) \\
f : \Sh(\Sch'_\tau) &
\longrightarrow &
\Sh(\Sch_\tau)
\end{eqnarray*}
tels que $f \circ g \cong \text{id}$. Pour tout objet $S$
de $\Sch_\tau$, le foncteur d'inclusion
$(\Sch/S)_\tau \to (\Sch'/S)_\tau$ satisfait lui aussi
aux hypothèses de Sites, lemme \ref{sites-lemma-bigger-site}.
On obtient donc de même des morphismes
\begin{eqnarray*}
g : \Sh((\Sch/S)_\tau) &
\longrightarrow &
\Sh((\Sch'/S)_\tau) \\
f : \Sh((\Sch'/S)_\tau) &
\longrightarrow &
\Sh((\Sch/S)_\tau)
\end{eqnarray*}
avec $f \circ g \cong \text{id}$.
\end{lemma}

\begin{proof}
Les hypothèses (b), (c) et (e) de
Sites, lemme \ref{sites-lemma-bigger-site},
sont immédiates pour les foncteurs
$\Sch_\tau \to \Sch'_\tau$ et
$(\Sch/S)_\tau \to (\Sch'/S)_\tau$. La propriété (a) résulte du
lemme \ref{lemma-zariski-induced},
\ref{lemma-etale-induced},
\ref{lemma-smooth-induced},
\ref{lemma-syntomic-induced}, ou
\ref{lemma-fppf-induced}.
La propriété (d) est satisfaite parce que
les produits fibrés existent dans les catégories $\Sch_\tau$, $\Sch'_\tau$
et sont compatibles avec les produits fibrés dans la catégorie des schémas.
\end{proof}

\noindent
Discussion :
Le foncteur $g^{-1} = f_*$ est simplement le foncteur de restriction qui associe
à un faisceau $\mathcal{G}$ sur $\Sch'_\tau$ sa restriction
$\mathcal{G}|_{\Sch_\tau}$. Ainsi, ce lemme affirme simplement qu'étant donné
un faisceau d'ensembles $\mathcal{F}$ sur $\Sch_\tau$, il existe un
faisceau canonique $\mathcal{F}'$ sur $\Sch'_\tau$ tel que
$\mathcal{F}|_{\Sch'_\tau} = \mathcal{F}'$. En fait, le faisceau
$\mathcal{F}'$ admet la description suivante : c'est le faisceau associé
au préfaisceau
$$
\Sch'_\tau \longrightarrow \textit{Ensembles}, \quad
V \longmapsto \colim_{V \to U} \mathcal{F}(U)
$$
où $U$ est un objet de $\Sch_\tau$. Cela résulte de l'égalité
$\mathcal{F}' = f^{-1}\mathcal{F} = (u_p\mathcal{F})^\#$ d'après
Sites, lemmes \ref{sites-lemma-when-shriek} et \ref{sites-lemma-bigger-site}.






\section{Prolongement de foncteurs}
\label{section-extending-functors}

\noindent
Commençons par un exemple simple qui explique notre démarche.
Soit $R$ un anneau. Supposons que $F$ soit un foncteur défini sur la catégorie
$\mathcal{C}$ des $R$-algèbres de la forme
$$
A = R[x_1, \ldots, x_n]/(f_1, \ldots, f_m)
$$
où $n, m \geq 0$ sont des entiers et $f_1, \ldots, f_m \in R[x_1, \ldots, x_n]$
des éléments. Pour toute $R$-algèbre $B$, on peut alors définir
$$
F'(B) = \colim_{A \to B,\ A \in \mathcal{C}} F(A)
$$
Il se trouve que $F'$ est l'unique foncteur sur la catégorie
de toutes les $R$-algèbres qui prolonge $F$ et commute aux limites
inductives filtrantes. Le même procédé fonctionne dans la catégorie des schémas
si l'on impose que notre foncteur soit un faisceau de Zariski.

\begin{lemma}
\label{lemma-extend}
Soit $S$ un schéma. Soit $\mathcal{C}$ une sous-catégorie pleine
de la catégorie $\Sch/S$ de tous les schémas au-dessus de $S$. Supposons que
\begin{enumerate}
\item si $X \to S$ est un objet de $\mathcal{C}$ et si
$U \subset X$ est un ouvert affine, alors $U \to S$ est isomorphe
à un objet de $\mathcal{C}$ ;
\item si $V$ est un schéma affine au-dessus d'un ouvert affine $U \subset S$
tel que $V \to U$ soit de présentation finie, alors $V \to S$ est isomorphe
à un objet de $\mathcal{C}$.
\end{enumerate}
Soit $F : \mathcal{C}^{opp} \to \textit{Ensembles}$ un foncteur.
Supposons que
\begin{enumerate}
\item[(a)] pour tout recouvrement de Zariski $\{f_i : X_i \to X\}_{i \in I}$
où $X, X_i$ sont des objets de $\mathcal{C}$, la condition de
faisceau soit satisfaite par $F$ et cette famille\footnote{Comme nous
ne savons pas si $X_i \times_X X_j$ appartient à $\mathcal{C}$,
il faut l'interpréter comme suit : d'après la propriété (1),
il existe des recouvrements de Zariski $\{U_{ijk} \to X_i \times_X X_j\}_{k \in K_{ij}}$
où $U_{ijk}$ est un objet de $\mathcal{C}$. La condition de faisceau
signifie alors que $F(X)$ est l'égalisateur des deux applications
de $\prod F(X_i)$ vers $\prod F(U_{ijk})$.} ;
\item[(b)] si $X = \lim X_i$ est une limite projective filtrante de schémas affines
au-dessus de $S$, où $X, X_i$ sont des objets de $\mathcal{C}$, alors
$F(X) = \colim F(X_i)$.
\end{enumerate}
Il existe alors une manière unique de prolonger $F$ en un foncteur
$F' : (\Sch/S)^{opp} \to \textit{Ensembles}$ satisfaisant aux analogues de
(a) et (b), c'est-à-dire que $F'$ satisfait la condition de faisceau pour tout
recouvrement de Zariski et que $F'(X) = \colim F'(X_i)$ dès que $X = \lim X_i$
est une limite projective filtrante de schémas affines au-dessus de $S$.
\end{lemma}

\begin{proof}
L'idée consiste d'abord à prolonger $F$ à une collection suffisamment grande de
schémas affines au-dessus de $S$, puis à utiliser la propriété de faisceau de Zariski pour le prolonger
à tous les schémas.

\medskip\noindent
Supposons que $V$ soit un schéma affine au-dessus de $S$ dont le morphisme structural
$V \to S$ se factorise par un ouvert affine $U \subset S$. Dans ce cas,
on peut écrire
$$
V = \lim V_i
$$
comme une limite projective filtrante, où $V_i \to U$ est de présentation finie
et $V_i$ est affine. Voir Algèbre, lemme \ref{algebra-lemma-ring-colimit-fp}.
D'après les conditions (1) et (2),
on peut remplacer les $V_i$ par des objets de $\mathcal{C}$.
Observons que $V_i \to S$ est localement de présentation finie
(si $S$ est quasi-séparé, ces morphismes sont même
de présentation finie). Posons alors
$$
F'(V) = \colim F(V_i)
$$
On peut en fait donner l'expression plus canonique
$$
F'(V) = \colim_{V \to V'} F(V')
$$
où la limite inductive est prise sur la catégorie des morphismes $V \to V'$ au-dessus de $S$
pour lesquels $V'$ est un objet de $\mathcal{C}$ dont le morphisme
structural $V' \to S$ est localement de présentation finie.
Cette expression coïncide avec la première parce que, d'après
Limits, proposition
\ref{limits-proposition-characterize-locally-finite-presentation}
notre système projectif $V_i$ est cofinal dans cette catégorie !
Enfin, notons que si $V$ était un objet de $\mathcal{C}$,
alors $F'(V) = F(V)$ d'après l'hypothèse (b).

\medskip\noindent
La seconde formule fait de $F'$ un foncteur contravariant
sur la catégorie formée des schémas affines $V$ au-dessus de $S$ dont le
morphisme structural se factorise par un ouvert affine de $S$.
Soit $V$ un tel schéma affine au-dessus de $S$ et
supposons que $V = \bigcup_{k = 1, \ldots, n} V_k$ soit un recouvrement ouvert fini
par des schémas affines. Il est alors légitime de demander si la condition de faisceau
est satisfaite par $F'$ et ce recouvrement ouvert.
C'est vrai et facile à montrer : écrivons $V = \lim V_i$ comme au
paragraphe précédent. D'après Limits, lemme \ref{limits-lemma-descend-opens},
pour tout $i$ suffisamment grand, on peut trouver des ouverts affines
$V_{i, k} \subset V_i$ compatibles aux morphismes de transition
et dont l'image inverse est $V_k$ dans $V$. Ainsi,
$$
F'(V_k) = \colim F(V_{i, k})
\quad\text{et}\quad
F'(V_k \cap V_l) = \colim F(V_{i, k} \cap V_{i, l})
$$
À strictement parler, il faut, dans ces formules, remplacer $V_{i, k}$ et
$V_{i, k} \cap V_{i, l}$ par des objets affines isomorphes de $\mathcal{C}$
avant d'appliquer le foncteur $F$.
Comme $I$ est filtrant, les limites inductives commutent aux égalisateurs.
La condition de faisceau (b) pour $F$ et les recouvrements de Zariski
$\{V_{i, k} \to V_i\}$ entraîne donc la condition de
faisceau pour $F'$ et ce recouvrement.

\medskip\noindent
Soit $X$ un schéma quelconque au-dessus de $S$. Notons $\mathcal{B}_X$
la collection des ouverts affines de $X$ dont le morphisme structural
vers $S$ se factorise par un ouvert affine de $S$. Il est clair que
$\mathcal{B}_X$ est une base de la topologie de $X$.
D'après le résultat du paragraphe précédent et Faisceaux, lemme
\ref{sheaves-lemma-cofinal-systems-coverings-standard-case}
on voit que $F'$ est un faisceau sur $\mathcal{B}_X$.
La restriction de $F'$ à $\mathcal{B}_X$
se prolonge donc uniquement en un faisceau $F'_X$ sur $X$ ; voir
Faisceaux, lemme \ref{sheaves-lemma-extend-off-basis}.
Si $X$ est un objet de $\mathcal{C}$, on dispose d'une
identification canonique $F'_X(X) = F(X)$, car $F'$ et $F$ coïncident
sur les objets où ils sont tous deux définis et
parce que $F$ satisfait la condition de faisceau pour les
recouvrements de Zariski.

\medskip\noindent
Soit $f : X \to Y$ un morphisme de schémas au-dessus de $S$.
On obtient une unique $f$-application de $F'_Y$ vers $F'_X$ compatible
avec les applications $F'(V) \to F'(U)$ pour tous
$U \in \mathcal{B}_X$ et $V \in \mathcal{B}_Y$
tels que $f(U) \subset V$ ; voir
Faisceaux, lemme \ref{sheaves-lemma-f-map-basis-above-and-below-structures}.
Nous omettons la vérification que ces applications se composent
correctement pour des morphismes $X \to Y \to Z$ de schémas au-dessus de $S$.
Nous omettons aussi la vérification que, si $f$ est un morphisme
de $\mathcal{C}$, alors l'application induite $F'_Y(Y) \to F'_X(X)$
coïncide avec l'application $F(Y) \to F(X)$ via les identifications
$F'_X(X) = F(X)$ et $F'_Y(Y) = F(Y)$ ci-dessus.
On voit ainsi que le prolongement cherché de
$F$ est le foncteur qui envoie $X/S$ sur $F'_X(X)$.

\medskip\noindent
La propriété (a) pour le foncteur $X \mapsto F'_X(X)$ est presque immédiate
d'après la construction ; nous omettons les détails.
Supposons que $X = \lim_{i \in I} X_i$
soit une limite projective filtrante de schémas affines au-dessus de $S$. Il faut montrer que
$$
F'_X(X) = \colim_{i \in I} F'_{X_i}(X_i)
$$
Supposons d'abord qu'il existe un $i \in I$ tel que
$X_i \to S$ se factorise par un ouvert affine $U \subset S$.
Alors $F'$ est défini sur $X$ et sur $X_{i'}$ pour $i' \geq i$,
et l'on a $F'_{X_{i'}}(X_{i'}) = F'(X_{i'})$ pour
$i' \geq i$ ainsi que $F'_X(X) = F'(X)$. Dans ce cas, toute flèche
$X \to V$, où $V$ est localement de présentation finie
au-dessus de $S$, se factorise sous la forme $X \to X_{i'} \to V$ pour un certain
$i' \geq i$ ; voir Limits, proposition
\ref{limits-proposition-characterize-locally-finite-presentation}.
On a donc
\begin{align*}
F'_X(X)
& =
F'(X)  \\
& = \colim_{X \to V} F(V) \\
& =
\colim_{i' \geq i} \colim_{X_{i'} \to V} F(V) \\
& =
\colim_{i' \geq i} F'(X_{i'}) \\
& =
\colim_{i' \geq i} F'_{X_{i'}}(X_{i'}) \\
& =
\colim_{i' \in I} F'_{X_{i'}}(X_{i'})
\end{align*}
comme voulu. Enfin, dans le cas général, choisissons un $i \in I$ et un
recouvrement ouvert affine fini $V_i = V_{i, 1} \cup \ldots \cup V_{i, n}$
tel que $V_{i, k} \to S$ se factorise par un ouvert affine de $S$.
Soient $V_k \subset V$ et $V_{i', k}$, pour $i' \geq i$,
les images inverses de $V_{i, k}$.
Le cas précédent donne
$$
F'_{V_k}(V_k) = \colim_{i' \geq i} F'_{V_{i', k}}(V_{i', k})
$$
et
$$
F'_{V_k \cap V_l}(V_k \cap V_l) =
\colim_{i' \geq i}
F'_{V_{i', k} \cap V_{i', l}}(V_{i', k} \cap V_{i', l})
$$
La propriété de faisceau et l'exactitude des limites inductives filtrantes
donnent $F'_X(X) = \colim_{i \in I} F'_{X_i}(X_i)$
également dans ce cas. Cela achève la démonstration de la propriété (b)
et donc celle du lemme.
\end{proof}

\begin{lemma}
\label{lemma-limit-fppf-topology}
Soit $\tau \in \{Zariski, \etale, smooth, syntomic, fppf\}$.
Soit $T$ un schéma affine écrit comme limite
$T = \lim_{i \in I} T_i$ d'un système projectif filtrant de schémas affines.
\begin{enumerate}
\item Soit $\mathcal{V} = \{V_j \to T\}_{j = 1, \ldots, m}$ un
recouvrement $\tau$ standard de $T$; voir les définitions
\ref{definition-standard-Zariski},
\ref{definition-standard-etale},
\ref{definition-standard-smooth},
\ref{definition-standard-syntomic} et
\ref{definition-standard-fppf}.
Il existe alors un indice $i$ et un recouvrement $\tau$ standard
$\mathcal{V}_i = \{V_{i, j} \to T_i\}_{j = 1, \ldots, m}$
dont le changement de base $T \times_{T_i} \mathcal{V}_i$ à $T$
est isomorphe à $\mathcal{V}$.
\item Soient $\mathcal{V}_i$, $\mathcal{V}'_i$ deux recouvrements
$\tau$ standards de $T_i$. Si
$f : T \times_{T_i} \mathcal{V}_i \to T \times_{T_i} \mathcal{V}'_i$ est
un morphisme de recouvrements de $T$, il existe un indice
$i' \geq i$ et un morphisme
$f_{i'} : T_{i'} \times_{T_i} \mathcal{V} \to
T_{i'} \times_{T_i} \mathcal{V}'_i$
dont le changement de base à $T$ est $f$.
\item Si
$f, g : \mathcal{V} \to \mathcal{V}'_i$
sont des morphismes de recouvrements $\tau$ standards de $T_i$ dont
les changements de base $f_T, g_T$ à $T$ sont égaux, alors il existe un
indice $i' \geq i$ tel que $f_{T_{i'}} = g_{T_{i'}}$.
\end{enumerate}
Autrement dit, la catégorie des recouvrements $\tau$ standards de $T$ est
la limite inductive, sur $I$, des catégories des recouvrements $\tau$ standards de $T_i$.
\end{lemma}

\begin{proof}
Démontrons l'énoncé pour $\tau = fppf$.
D'après Limits, lemme \ref{limits-lemma-descend-finite-presentation},
la catégorie des schémas de présentation finie sur $T$ est la
limite inductive, sur $I$, des catégories des schémas de présentation finie sur $T_i$. D'après
Limits, lemmes \ref{limits-lemma-descend-affine-finite-presentation}
et \ref{limits-lemma-descend-flat-finite-presentation},
il en va de même pour la catégorie des schémas affines, plats et
de présentation finie sur $T$.
Pour achever la démonstration du lemme, il suffit de montrer que, si
$\{V_{j, i} \to T_i\}_{j = 1, \ldots, m}$ est une famille finie de
morphismes plats de présentation finie, avec $V_{j, i}$ affine, et si le
changement de base $\coprod_j T \times_{T_i} V_{j, i} \to T$ est surjectif,
alors, pour un certain $i' \geq i$, le morphisme
$\coprod T_{i'} \times_{T_i} V_{j, i} \to T_{i'}$ est surjectif.
Notons $W_{i'} \subset T_{i'}$, respectivement $W \subset T$, l'image.
Par hypothèse, on a bien sûr $W = T$.
Puisque les morphismes sont plats et de présentation finie,
$W_i$ est un ouvert quasi-compact de $T_i$; voir
Morphismes, lemme \ref{morphisms-lemma-fppf-open}.
En outre, $W = T \times_{T_i} W_i$ (la formation de l'image commute
au changement de base). Par conséquent, d'après
Limits, lemme \ref{limits-lemma-descend-opens},
on a $W_{i'} = T_{i'}$ pour un indice $i'$ assez grand,
ce qui conclut.

\medskip\noindent
Pour $\tau \in \{Zariski, \etale, smooth, syntomic\}$, tout recouvrement $\tau$ standard
est un recouvrement fppf standard. La pleine fidélité du foncteur
en résulte. Il reste seulement à montrer que, si un recouvrement fppf standard
$\mathcal{V}_i$ est donné pour un certain $i$ et si $\mathcal{V}_i \times_{T_i} T$
est un recouvrement $\tau$ standard, alors $\mathcal{V}_i \times_{T_i} T_{i'}$
est un recouvrement $\tau$ standard pour tout $i' \gg i$. Cela résulte immédiatement
de Limits, lemmes
\ref{limits-lemma-descend-open-immersion},
\ref{limits-lemma-descend-etale},
\ref{limits-lemma-descend-smooth} et
\ref{limits-lemma-descend-syntomic}.
\end{proof}

\begin{lemma}
\label{lemma-extend-sheaf-general}
Supposons que $S$, $\mathcal{C}$ et $F$ satisfassent aux conditions (1), (2), (a) et (b) du
lemme \ref{lemma-extend}, et notons $F' : (\Sch/S)^{opp} \to \textit{Ensembles}$
le prolongement unique construit dans ce lemme. Soit
$\tau \in \{Zariski, \etale, smooth, syntomic, fppf\}$. Supposons que
\begin{enumerate}
\item[(c)] pour tout recouvrement $\tau$ standard $\{V_i \to V\}_{i = 1, \ldots, n}$
par des schémas affines dans $\Sch/S$, tel que $V \to S$ se factorise par un ouvert affine
$U \subset S$ et que $V \to U$ soit de présentation finie, la condition de faisceau
vaut pour $F$ et $\{V_i \to V\}_{i = 1, \ldots, n}$\footnote{Cela a
un sens, puisque $V$, $V_i$ et $V_i \times_V V_j$ sont isomorphes à des objets
de $\mathcal{C}$ d'après (2).}.
\end{enumerate}
Alors $F'$ satisfait à la condition de faisceau pour tout recouvrement $\tau$.
\end{lemma}

\begin{proof}
Soit $X$ un schéma sur $S$, et soit $\{X_i \to X\}_{i \in I}$
un recouvrement $\tau$. Soient $s_i \in F'(X_i)$ des éléments tels que
$s_i$ et $s_j$ aient même image dans $F'(X_i \times_X X_j)$
pour tous $i, j \in I$. Il faut montrer qu'il existe un unique
élément $s \in F'(X)$ dont la restriction est $s_i \in F'(X_i)$ pour tout $i \in I$.

\medskip\noindent
Cas particulier : $X$ est affine, son morphisme structural
est à valeurs dans un ouvert affine $U$ de $S$, et le recouvrement
$\{X_i \to X\}_{i \in I}$ est un recouvrement $\tau$ standard.
Dans ce cas, on peut écrire
$$
X = \lim V_k
$$
comme limite projective filtrante, où $V_k \to U$ est de présentation finie
et $V_k$ est affine. Voir Algèbre, lemme \ref{algebra-lemma-ring-colimit-fp}.
D'après lemme \ref{lemma-limit-fppf-topology}, il existe un $k$ et un recouvrement
$\tau$ standard $\{V_{k, i} \to V_k\}_{i \in I}$ dont le changement de base
à $X$ est le recouvrement donné. Pour $k' \geq k$, notons
$\{V_{k', i} \to V_{k'}\}_{i \in I}$ le changement de base à $V_{k'}$
de ce recouvrement. On a alors
\begin{align*}
F'(X)
& =
\colim_{k' \geq k} F(V_k) \\
& =
\colim_{k' \geq k}
\text{Equalizer}(
\xymatrix{
\prod F(V_{k', i})
\ar@<1ex>[r] \ar@<-1ex>[r] &
\prod F(V_{k', i} \times_{V_{k'}} V_{k', j})
}
\\
& =
\text{Equalizer}(
\xymatrix{
\colim_{k' \geq k}
\prod F(V_{k', i})
\ar@<1ex>[r] \ar@<-1ex>[r] &
\colim_{k' \geq k}
\prod F(V_{k', i} \times_{V_{k'}} V_{k', j})
}
\\
& =
\text{Equalizer}(
\xymatrix{
\prod F'(X_i)
\ar@<1ex>[r] \ar@<-1ex>[r] &
\prod F'(X_i \times_X X_j)
}
\end{align*}
La première égalité résulte de la construction de $F'$. La deuxième résulte de
l'hypothèse (c). La troisième résulte de l'exactitude des limites inductives filtrantes.
La quatrième résulte encore de la construction de $F'$.
On obtient ainsi la propriété de faisceau pour $F'$
relativement à $\{X_i \to X\}_{i \in I}$.

\medskip\noindent
Cas général. Choisissons un recouvrement ouvert affine $X = \bigcup U_k$ tel que
chaque $U_k$ soit à valeurs dans un ouvert affine de $S$. Pour tout $k$, choisissons
un recouvrement $\tau$ standard $\{V_{k, j} \to U_k\}_{j = 1, \ldots, m_k}$
qui raffine $\{X_i \times_X U_k \to U_k\}_{i \in I}$. Pour chaque
$j \in \{1, \ldots, m_k\}$, choisissons un indice $i_{k, j} \in I$
et un morphisme $g_{k, j} : V_{k, j} \to X_{i_{k, j}}$ au-dessus de $X$.
Soit $s_{k, j}$ l'élément de $F'(V_{k, j})$ obtenu en restreignant
$s_{i_{k, j}}$ par $g_{k, j}$. Observons que $s_{k, j}$ et $s_{k', j'}$
ont même restriction dans $F'(V_{k, j} \times_X V_{k', j'})$
pour tous $k$ et $k'$, tout $j \in \{1, \ldots, m_k\}$
et tout $j' \in \{1, \ldots, m_{k'}\}$; nous omettons la vérification.
En particulier, le paragraphe précédent donne un unique
élément $s_k \in F'(U_k)$ dont la restriction est $s_{k, j}$ pour tout $j$.
Avec ces notations, nous pouvons achever la démonstration.

\medskip\noindent
Unicité de $s$ : elle résulte de ce que $F'$ satisfait à la
propriété de faisceau pour les recouvrements de Zariski et que $s|_{U_k}$ doit être égal
à $s_k$, puisque tous deux ont pour restriction $s_{k, j}$ pour tout $j$.
Cette unicité montre alors que $s_k$ et $s_{k'}$ doivent induire
la même section de $F'$ sur (le schéma non affine) $U_k \cap U_{k'}$,
car ces sections induisent la même section sur le
recouvrement $\tau$ $\{V_{k, j} \times_X V_{k', j'} \to U_k \cap U_{k'}\}$.
La propriété de faisceau pour les recouvrements de Zariski donne donc une unique
section $s$ de $F'$ sur $X$ dont la restriction à $U_k$ est $s_k$.
Nous omettons la vérification (analogue à la précédente) que $s$ se restreint
en $s_i$ sur $X_i$.
\end{proof}

\begin{lemma}
\label{lemma-extend-sheaf}
Soit $\tau \in \{Zariski, \etale, smooth, syntomic, fppf\}$.
Soit $S$ un schéma contenu dans un gros site $\Sch_\tau$.
Soit $F : (\Sch/S)_\tau^{opp} \to \textit{Ensembles}$ un $\tau$-faisceau
satisfaisant à la propriété (b) de lemme \ref{lemma-extend} pour
$\mathcal{C} = (\Sch/S)_\tau$. Alors le prolongement
$F'$ de $F$ à la catégorie de tous les schémas sur $S$
satisfait à la condition de faisceau pour tout recouvrement $\tau$.
\end{lemma}

\begin{proof}
Cela résulte de lemme \ref{lemma-extend-sheaf-general}
appliqué avec $\mathcal{C} = (\Sch/S)_\tau$.
Les conditions (1), (2), (a) et (b) de lemme \ref{lemma-extend}
sont satisfaites; nous omettons les détails. On obtient ainsi l'unique prolongement $F'$
à la catégorie de tous les schémas sur $S$. Observons enfin que
tout recouvrement $\tau$ standard est tautologiquement équivalent à un
recouvrement dans $(\Sch/S)_\tau$; voir Ensembles, lemme \ref{sets-lemma-what-is-in-it}
ainsi que lemmes \ref{lemma-zariski-induced},
\ref{lemma-etale-induced},
\ref{lemma-smooth-induced},
\ref{lemma-syntomic-induced} et
\ref{lemma-fppf-induced}. D'après
Sites, lemme \ref{sites-lemma-tautological-same-sheaf}
la propriété de faisceau se conserve par équivalence tautologique
des recouvrements. Le fait que $F$ soit un $\tau$-faisceau entraîne donc
la propriété (c) de lemme \ref{lemma-extend-sheaf-general},
ce qui conclut.
\end{proof}




\begin{multicols}{2}[\section{Autres chapitres}]
\noindent
Préliminaires
\begin{enumerate}
\item \hyperref[introduction-section-phantom]{Introduction}
\item \hyperref[conventions-section-phantom]{Conventions}
\item \hyperref[sets-section-phantom]{Théorie des ensembles}
\item \hyperref[categories-section-phantom]{Catégories}
\item \hyperref[topology-section-phantom]{Topologie}
\item \hyperref[sheaves-section-phantom]{Faisceaux sur les espaces}
\item \hyperref[sites-section-phantom]{Sites et faisceaux}
\item \hyperref[stacks-section-phantom]{Champs}
\item \hyperref[fields-section-phantom]{Corps}
\item \hyperref[algebra-section-phantom]{Algèbre commutative}
\item \hyperref[brauer-section-phantom]{Groupes de Brauer}
\item \hyperref[homology-section-phantom]{Algèbre homologique}
\item \hyperref[derived-section-phantom]{Catégories dérivées}
\item \hyperref[simplicial-section-phantom]{Méthodes simpliciales}
\item \hyperref[more-algebra-section-phantom]{Compléments d'algèbre}
\item \hyperref[smoothing-section-phantom]{Lissage des morphismes d'anneaux}
\item \hyperref[modules-section-phantom]{Faisceaux de modules}
\item \hyperref[sites-modules-section-phantom]{Modules sur les sites}
\item \hyperref[injectives-section-phantom]{Objets injectifs}
\item \hyperref[cohomology-section-phantom]{Cohomologie des faisceaux}
\item \hyperref[sites-cohomology-section-phantom]{Cohomologie sur les sites}
\item \hyperref[dga-section-phantom]{Algèbre différentielle graduée}
\item \hyperref[dpa-section-phantom]{Algèbre à puissances divisées}
\item \hyperref[sdga-section-phantom]{Faisceaux différentiels gradués}
\item \hyperref[hypercovering-section-phantom]{Hyperrecouvrements}
\end{enumerate}
Schémas
\begin{enumerate}
\setcounter{enumi}{25}
\item \hyperref[schemes-section-phantom]{Schémas}
\item \hyperref[constructions-section-phantom]{Constructions de schémas}
\item \hyperref[properties-section-phantom]{Propriétés des schémas}
\item \hyperref[morphisms-section-phantom]{Morphismes de schémas}
\item \hyperref[coherent-section-phantom]{Cohomologie des schémas}
\item \hyperref[divisors-section-phantom]{Diviseurs}
\item \hyperref[limits-section-phantom]{Limites de schémas}
\item \hyperref[varieties-section-phantom]{Variétés}
\item \hyperref[topologies-section-phantom]{Topologies sur les schémas}
\item \hyperref[descent-section-phantom]{Descente}
\item \hyperref[perfect-section-phantom]{Catégories dérivées de schémas}
\item \hyperref[more-morphisms-section-phantom]{Compléments sur les morphismes}
\item \hyperref[flat-section-phantom]{Compléments sur la platitude}
\item \hyperref[groupoids-section-phantom]{Schémas en groupoïdes}
\item \hyperref[more-groupoids-section-phantom]{Compléments sur les schémas en groupoïdes}
\item \hyperref[etale-section-phantom]{Morphismes étales de schémas}
\end{enumerate}
Thèmes de théorie des schémas
\begin{enumerate}
\setcounter{enumi}{41}
\item \hyperref[chow-section-phantom]{Homologie de Chow}
\item \hyperref[intersection-section-phantom]{Théorie de l'intersection}
\item \hyperref[pic-section-phantom]{Schémas de Picard des courbes}
\item \hyperref[weil-section-phantom]{Théories cohomologiques de Weil}
\item \hyperref[adequate-section-phantom]{Modules adéquats}
\item \hyperref[dualizing-section-phantom]{Complexes dualisants}
\item \hyperref[duality-section-phantom]{Dualité pour les schémas}
\item \hyperref[discriminant-section-phantom]{Discriminants et différentes}
\item \hyperref[derham-section-phantom]{Cohomologie de de Rham}
\item \hyperref[local-cohomology-section-phantom]{Cohomologie locale}
\item \hyperref[algebraization-section-phantom]{Géométrie algébrique et formelle}
\item \hyperref[curves-section-phantom]{Courbes algébriques}
\item \hyperref[resolve-section-phantom]{Résolution des surfaces}
\item \hyperref[models-section-phantom]{Réduction semi-stable}
\item \hyperref[functors-section-phantom]{Foncteurs et morphismes}
\item \hyperref[equiv-section-phantom]{Catégories dérivées de variétés}
\item \hyperref[pione-section-phantom]{Groupes fondamentaux de schémas}
\item \hyperref[etale-cohomology-section-phantom]{Cohomologie étale}
\item \hyperref[crystalline-section-phantom]{Cohomologie cristalline}
\item \hyperref[proetale-section-phantom]{Cohomologie pro-étale}
\item \hyperref[relative-cycles-section-phantom]{Cycles relatifs}
\item \hyperref[more-etale-section-phantom]{Compléments de cohomologie étale}
\item \hyperref[trace-section-phantom]{La formule des traces}
\end{enumerate}
Espaces algébriques
\begin{enumerate}
\setcounter{enumi}{64}
\item \hyperref[spaces-section-phantom]{Espaces algébriques}
\item \hyperref[spaces-properties-section-phantom]{Propriétés des espaces algébriques}
\item \hyperref[spaces-morphisms-section-phantom]{Morphismes d'espaces algébriques}
\item \hyperref[decent-spaces-section-phantom]{Espaces algébriques décents}
\item \hyperref[spaces-cohomology-section-phantom]{Cohomologie des espaces algébriques}
\item \hyperref[spaces-limits-section-phantom]{Limites d'espaces algébriques}
\item \hyperref[spaces-divisors-section-phantom]{Diviseurs sur les espaces algébriques}
\item \hyperref[spaces-over-fields-section-phantom]{Espaces algébriques sur des corps}
\item \hyperref[spaces-topologies-section-phantom]{Topologies sur les espaces algébriques}
\item \hyperref[spaces-descent-section-phantom]{Descente et espaces algébriques}
\item \hyperref[spaces-perfect-section-phantom]{Catégories dérivées d'espaces}
\item \hyperref[spaces-more-morphisms-section-phantom]{Compléments sur les morphismes d'espaces}
\item \hyperref[spaces-flat-section-phantom]{Platitude sur les espaces algébriques}
\item \hyperref[spaces-groupoids-section-phantom]{Groupoïdes dans les espaces algébriques}
\item \hyperref[spaces-more-groupoids-section-phantom]{Compléments sur les groupoïdes dans les espaces}
\item \hyperref[bootstrap-section-phantom]{Amorçage}
\item \hyperref[spaces-pushouts-section-phantom]{Sommes amalgamées d'espaces algébriques}
\end{enumerate}
Thèmes de géométrie
\begin{enumerate}
\setcounter{enumi}{81}
\item \hyperref[spaces-chow-section-phantom]{Groupes de Chow des espaces}
\item \hyperref[groupoids-quotients-section-phantom]{Quotients de groupoïdes}
\item \hyperref[spaces-more-cohomology-section-phantom]{Compléments sur la cohomologie des espaces}
\item \hyperref[spaces-simplicial-section-phantom]{Espaces simpliciaux}
\item \hyperref[spaces-duality-section-phantom]{Dualité pour les espaces}
\item \hyperref[formal-spaces-section-phantom]{Espaces algébriques formels}
\item \hyperref[restricted-section-phantom]{Algébrisation des espaces formels}
\item \hyperref[spaces-resolve-section-phantom]{Résolution des surfaces revisitée}
\end{enumerate}
Théorie des déformations
\begin{enumerate}
\setcounter{enumi}{89}
\item \hyperref[formal-defos-section-phantom]{Théorie formelle des déformations}
\item \hyperref[defos-section-phantom]{Théorie des déformations}
\item \hyperref[cotangent-section-phantom]{Le complexe cotangent}
\item \hyperref[examples-defos-section-phantom]{Problèmes de déformation}
\end{enumerate}
Champs algébriques
\begin{enumerate}
\setcounter{enumi}{93}
\item \hyperref[algebraic-section-phantom]{Champs algébriques}
\item \hyperref[examples-stacks-section-phantom]{Exemples de champs}
\item \hyperref[stacks-sheaves-section-phantom]{Faisceaux sur les champs algébriques}
\item \hyperref[criteria-section-phantom]{Critères de représentabilité}
\item \hyperref[artin-section-phantom]{Axiomes d'Artin}
\item \hyperref[quot-section-phantom]{Espaces de Quot et de Hilbert}
\item \hyperref[stacks-properties-section-phantom]{Propriétés des champs algébriques}
\item \hyperref[stacks-morphisms-section-phantom]{Morphismes de champs algébriques}
\item \hyperref[stacks-limits-section-phantom]{Limites de champs algébriques}
\item \hyperref[stacks-cohomology-section-phantom]{Cohomologie des champs algébriques}
\item \hyperref[stacks-perfect-section-phantom]{Catégories dérivées de champs}
\item \hyperref[stacks-introduction-section-phantom]{Introduction aux champs algébriques}
\item \hyperref[stacks-more-morphisms-section-phantom]{Compléments sur les morphismes de champs}
\item \hyperref[stacks-geometry-section-phantom]{La géométrie des champs}
\end{enumerate}
Thèmes de théorie des espaces de modules
\begin{enumerate}
\setcounter{enumi}{107}
\item \hyperref[moduli-section-phantom]{Champs de modules}
\item \hyperref[moduli-curves-section-phantom]{Espaces de modules de courbes}
\end{enumerate}
Divers
\begin{enumerate}
\setcounter{enumi}{109}
\item \hyperref[examples-section-phantom]{Exemples}
\item \hyperref[exercises-section-phantom]{Exercices}
\item \hyperref[guide-section-phantom]{Guide bibliographique}
\item \hyperref[desirables-section-phantom]{Desiderata}
\item \hyperref[coding-section-phantom]{Style de programmation}
\item \hyperref[obsolete-section-phantom]{Obsolète}
\item \hyperref[fdl-section-phantom]{Licence de documentation libre GNU}
\item \hyperref[index-section-phantom]{Index généré automatiquement}
\end{enumerate}
\end{multicols}



\bibliography{my}
\bibliographystyle{amsalpha}

\end{document}
